% Kompendium TMA4432 Matematikk 4C, høst 2026
% Kompiler med:  latexmk -pdf kompendium.tex
\documentclass[11pt,a4paper]{article}

\usepackage[T1]{fontenc}
\usepackage[utf8]{inputenc}
\usepackage{babel}
\babelprovide[import, main]{norsk}
\usepackage{mathpazo}
\usepackage[scaled=0.92]{helvet}
\usepackage{microtype}
\usepackage[a4paper,margin=2.3cm,headheight=14pt]{geometry}
\usepackage{amsmath,amssymb,mathtools,bm}
\usepackage{xcolor}
\usepackage[most]{tcolorbox}
\usepackage{enumitem}
\usepackage{booktabs,array,tabularx}
\usepackage{etoolbox}
\usepackage{titlesec}
\usepackage{fancyhdr}
\usepackage{ragged2e}
\usepackage{multicol}
\usepackage[colorlinks=true,linkcolor=black,urlcolor=fsfarge]{hyperref}

% ---------- Utseende ----------
\definecolor{blaa}{RGB}{20,60,120}
\definecolor{groenn}{RGB}{25,105,70}
\definecolor{oransje}{RGB}{185,95,15}
\definecolor{roed}{RGB}{165,30,35}
\definecolor{fsfarge}{RGB}{0,110,120}
\definecolor{graa}{RGB}{90,90,90}

\RaggedRight
\setlength{\RaggedRightParindent}{0pt}
\setlength{\parskip}{4pt plus 1pt}
\emergencystretch=2em
\setlist{itemsep=1pt,topsep=3pt,parsep=1pt}
\setlist[itemize]{leftmargin=1.4em}
\setlist[enumerate]{leftmargin=1.6em}

\titleformat{\section}{\Large\sffamily\bfseries\color{blaa}}{\thesection}{0.6em}{}[{\color{blaa}\titlerule[0.8pt]}]
\titleformat{\subsection}{\large\sffamily\bfseries\color{blaa}}{\thesubsection}{0.5em}{}
\titleformat{\subsubsection}{\normalsize\sffamily\bfseries}{}{0em}{}
\titlespacing*{\section}{0pt}{10pt}{8pt}
\titlespacing*{\subsection}{0pt}{9pt}{4pt}
\titlespacing*{\subsubsection}{0pt}{7pt}{2pt}

\pagestyle{fancy}
\fancyhf{}
\fancyhead[L]{\small\sffamily\color{graa}TMA4432 Matematikk 4C -- kompendium}
\fancyhead[R]{\small\sffamily\color{graa}\leftmark}
\fancyfoot[C]{\small\sffamily\thepage}
\renewcommand{\headrulewidth}{0.3pt}
\renewcommand{\sectionmark}[1]{\markboth{#1}{}}

% ---------- Bokser ----------
\tcbset{
  felles/.style={enhanced, breakable, boxrule=0.6pt, arc=2pt, left=6pt, right=6pt, top=4pt, bottom=4pt,
    fonttitle=\sffamily\bfseries\small, before skip=6pt, after skip=6pt, parbox=false, halign=left},
}
\newcommand{\fs}{\hfill{\normalfont\sffamily\scriptsize\colorbox{white}{\color{fsfarge}\textbf{på formelarket}}}}
\newcommand{\boksnavn}[2]{#1\ifstrempty{#2}{}{: #2}}
\newenvironment{defn}[1][]{\begin{tcolorbox}[felles, colback=blaa!4, colframe=blaa!70, title={\boksnavn{Definisjon}{#1}}]}{\end{tcolorbox}}
\newenvironment{teorem}[1][]{\begin{tcolorbox}[felles, colback=groenn!4, colframe=groenn!75, title={\boksnavn{Resultat}{#1}}]}{\end{tcolorbox}}
\newenvironment{formel}[1][]{\begin{tcolorbox}[felles, colback=fsfarge!5, colframe=fsfarge!80, title={#1}]}{\end{tcolorbox}}
\newenvironment{oppskrift}[1][]{\begin{tcolorbox}[felles, colback=oransje!5, colframe=oransje!80, title={\boksnavn{Oppskrift}{#1}}]}{\end{tcolorbox}}
\newenvironment{fallgruve}[1][]{\begin{tcolorbox}[felles, colback=roed!4, colframe=roed!75, title={\boksnavn{Pass på}{#1}}]}{\end{tcolorbox}}
\newenvironment{eksempel}[1][]{\begin{tcolorbox}[felles, colback=graa!4, colframe=graa!70, title={\boksnavn{Eksempel}{#1}}]}{\end{tcolorbox}}
\newenvironment{oving}{\begin{tcolorbox}[felles, colback=white, colframe=graa!50, boxrule=0.4pt, title={Øv på dette (klikk for oppgave og løsningsforslag)}]\small}{\end{tcolorbox}}
\newcommand{\lett}{\textcolor{groenn}{\sffamily\bfseries Lett}}
\newcommand{\middels}{\textcolor{oransje}{\sffamily\bfseries Middels}}
\newcommand{\vanskelig}{\textcolor{roed}{\sffamily\bfseries Vanskelig}}

% ---------- Matematikk ----------
\newcommand{\R}{\mathbb{R}}
\newcommand{\C}{\mathbb{C}}
\newcommand{\N}{\mathbb{N}}
\newcommand{\Z}{\mathbb{Z}}
\newcommand{\T}{\mathsf{T}}
\newcommand{\dd}{\,\mathrm{d}}
\newcommand{\ii}{\mathrm{i}}
\newcommand{\ee}{\mathrm{e}}
\newcommand{\vb}[1]{\mathbf{#1}}
\newcommand{\tn}[1]{\bm{\mathsf{#1}}}
\newcommand{\norm}[1]{\lVert #1\rVert}
\newcommand{\abs}[1]{\lvert #1\rvert}
\newcommand{\ip}[2]{\langle #1,#2\rangle}
\newcommand{\pd}[2]{\frac{\partial #1}{\partial #2}}
\newcommand{\pdd}[2]{\frac{\partial^2 #1}{\partial #2^2}}
\newcommand{\Four}{\mathcal{F}}
\newcommand{\FN}{\mathcal{F}_N}
\newcommand{\kr}[1]{[\![ #1 ]\!]}
\newcommand{\had}{\circledast}
\newcommand{\outp}{\circ}
\DeclareMathOperator{\rank}{rank}
\DeclareMathOperator{\spn}{span}
\DeclareMathOperator{\ran}{range}
\DeclareMathOperator{\kerop}{ker}
\DeclareMathOperator{\supp}{supp}
\DeclareMathOperator{\vect}{vec}
\DeclareMathOperator{\Hess}{Hess}
\DeclareMathOperator{\diag}{diag}
\DeclareMathOperator*{\argmin}{arg\,min}
\newcommand{\sgn}{\operatorname{sign}}

\hypersetup{pdftitle={Kompendium TMA4432 Matematikk 4C},pdfauthor={Studiehjelp}}

\begin{document}

\begin{titlepage}
\centering
\vspace*{3cm}
{\sffamily\Huge\bfseries\color{blaa} Kompendium\par}
\vspace{0.6cm}
{\sffamily\LARGE TMA4432 Matematikk 4C\par}
\vspace{0.3cm}
{\sffamily\large NTNU, høst 2026\par}
\vspace{1.2cm}
{\large Alle temaer, begreper, formler og metoder fra forelesningene,\\ øvingene og prøveeksamenen\par}
\vspace{2cm}
\begin{minipage}{0.82\textwidth}
\small
Bygget på materialet på \texttt{wiki.math.ntnu.no/tma4432/2026h}: forelesningsnotatene 01--13 (Ronny Bergmann), notatene om SVD og tensordekomposisjoner, øvingene 1--12 med løsningsforslag, prøveeksamenen og formelarket. Rekkefølgen følger forelesningene. Formlene er sjekket mot kildene, og regneeksemplene er kontrollregnet med kode. Der lysarkene har feil, står det riktige her, med en merknad.

\medskip
Studiehjelp, ikke offisielt kursmateriell. Ved tvil gjelder forelesningene og formelarket.
\end{minipage}
\vfill
{\small Sist oppdatert 19.\,11.\,2026 \quad·\quad Eksamen 24.\,11.\,2026 kl.\ 15--19\par}
\end{titlepage}

\tableofcontents
\clearpage

\section{Før du begynner: eksamen, formelark og konvensjoner}

\subsection{Eksamen}
\begin{itemize}
  \item \textbf{Tid:} tirsdag 24.\,11.\,2026 kl.\ 15.00--19.00 (4 timer). Oppgavene er på engelsk, men du kan svare på norsk.
  \item \textbf{Hjelpemidler (kode D):} ingen trykte eller håndskrevne hjelpemidler, bare enkel kalkulator. Formelarket ligger ved oppgavesettet.
  \item \textbf{Hva kan komme:} alt fra forelesningene og øvingene. Det som ble gjennomgått mye, er mer sannsynlig. Prøveeksamenen har 10 oppgaver, én per hovedtema, på til sammen rundt 100 poeng.
  \item \textbf{Ingen programmering,} men du kan få spørsmål om algoritmer og om å forstå kode. Om tensorer kommer det ikke oppgaver der du regner ut en konkret dekomposisjon.
  \item \textbf{Alle svar skal begrunnes.} Utregningen teller mer enn svaret. Bruk matematiske uttrykk når du forklarer noe.
\end{itemize}

\subsection{Øvingsoppgaver med lenker}
Etter hvert tema står en boks med oppgaver fra oppgavebanken i repoet. Oppgavenavnet lenker til oppgaven og «LF» til løsningsforslaget, begge som PDF på GitHub. Oppgavene er sortert etter vanskelighetsgrad. Gradene er mine vurderinger ut fra hvor mange steg oppgaven har, hvor mye regning den krever, og om den ber om bevis eller kombinerer flere temaer:
\begin{itemize}[leftmargin=5.6em,labelwidth=5em,labelsep=0.6em,align=left]
  \item[\lett] Én metode brukt direkte, få steg. Bra for å komme i gang.
  \item[\middels] Vanlig eksamensoppgave med flere deloppgaver. Prøveeksamenen ligger stort sett her.
  \item[\vanskelig] Lang regning, et bevis eller flere temaer kombinert, ofte fra mer avanserte emner. Lurt hvis du vil ha god margin.
\end{itemize}
Noen av Bergmanns eksamensoppgaver fra Matte 4N vår 2025 er gjenbrukt nesten ordrett på prøveeksamenen. Da lenker kompendiet bare til prøveeksamen-versjonen.

\subsection{Dette står på formelarket}
Formler merket \colorbox{white}{\sffamily\scriptsize\color{fsfarge}\textbf{på formelarket}} i kompendiet får du utdelt. Alt annet må du kunne selv. Formelarket inneholder:
\begin{enumerate}
  \item Jacobimatrisen $\vb F'(\vb a)=\bigl(\pd{F_i}{x_j}(\vb a)\bigr)_{i,j}$ og lineariseringen $\vb L(\vb x)=\vb F(\vb a)+\vb F'(\vb a)(\vb x-\vb a)$.
  \item Newtons metode for systemer: $\vb F'(\vb x^{(k)})\Delta^{(k)}=-\vb F(\vb x^{(k)})$, $\vb x^{(k+1)}=\vb x^{(k)}+\Delta^{(k)}$.
  \item Fourierrekker for $2\pi$-periodiske funksjoner: formlene for $a_0,a_n,b_n$ og $c_k$.
  \item DFT og IDFT med faktoren $\tfrac1N$ i DFT-en, og Fouriermatrisen $\FN$.
  \item Fouriertransformen og den inverse transformen med faktoren $\tfrac{1}{\sqrt{2\pi}}$, og en tabell over fire transformpar: $\ee^{-ax^2}$, $\ee^{-a\abs x}$, indikatorfunksjonen for $\abs x<a$ og $\sqrt{2/\pi}\,\sin(ax)/x$.
  \item Trigonometriske identiteter: addisjonsformlene, $\cos 2\alpha$, Eulers formel og produkt-til-sum-formlene.
  \item Bølgelikningen på $[0,L]$ med løsningen fra separasjon av variable, inkludert formlene for $A_n$ og $B_n$.
\end{enumerate}
\textbf{Ikke} på formelarket (må kunnes): Fourierrekker med periode $2L$, like og odde utvidelser, Parseval, konvolusjon, varmelikningens løsning, d'Alemberts formel, alle differanseskjemaer, SVD og Moore--Penrose, tensoroperasjoner, kurver og flater, gradient og Lagrange, splines og Bézier, og minste kvadrater.

\subsection{Konvensjoner som brukes i kurset}
\begin{fallgruve}[Normeringer varierer mellom bøker. Kurset bruker disse:]
\begin{itemize}
  \item \textbf{Fourierkoeffisienter:} $c_k=\frac{1}{2\pi}\int_{-\pi}^{\pi}f(x)\ee^{-\ii kx}\dd x$ og $a_0=\frac{1}{2\pi}\int_{-\pi}^{\pi}f$ (altså middelverdien, ikke $\frac{a_0}{2}$).
  \item \textbf{Periodisk konvolusjon} uten faktor: $(f*g)(x)=\int_{-\pi}^{\pi}f(y)g(x-y)\dd y$. Da er $c_k(f*g)=2\pi\,c_k(f)c_k(g)$.
  \item \textbf{DFT med $\frac1N$ foran:} $\hat f_k=\frac1N\sum_j f_j\ee^{-2\pi\ii jk/N}$, så $\hat f_k\approx c_k(f)$. Oppsummeringsforelesningen og prøveeksamenen skriver $\hat{\vb f}=\FN\vb f$ uten $\frac1N$, men regner som om faktoren er med. Bruk formelarket.
  \item \textbf{Fouriertransform:} $\hat f(\omega)=\frac{1}{\sqrt{2\pi}}\int_{-\infty}^\infty f(x)\ee^{-\ii\omega x}\dd x$, med samme faktor i den inverse transformen. Da er $\Four(f*g)=\sqrt{2\pi}\,\Four(f)\Four(g)$.
\end{itemize}
\end{fallgruve}

\subsection{Notasjon}
\begin{itemize}
  \item Tall skrives med små bokstaver ($x$, $\lambda$), vektorer med små fete ($\vb x\in\R^n$, alltid kolonnevektorer), matriser med store ($A\in\R^{m\times n}$) og tensorer med store fete ($\tn X\in\R^{m\times n\times p}$). For hånd markeres fet skrift ofte med strek under: $\underline{x}$.
  \item $A^\T$ er den transponerte, $\bar z$ den komplekskonjugerte, og $I_n$ identitetsmatrisen. $\ip{\vb x}{\vb y}=\vb x^\T\vb y$ er indreproduktet, og $\norm{\vb x}_2=\sqrt{\ip{\vb x}{\vb x}}$.
  \item $\norm{A}_2=\sigma_1$ er spektralnormen (største singulærverdi), og $\norm{A}_F=\bigl(\sum_{i,j}a_{ij}^2\bigr)^{1/2}$ er Frobeniusnormen.
  \item $\vb F'(\vb a)$ er Jacobimatrisen og $\nabla f$ gradienten (en kolonnevektor). $\Hess f$ (eller $\nabla^2 f$) er Hessematrisen.
  \item $U_i^n$ er en numerisk tilnærming til $u(x_i,t_n)$. Her er $n$ en tidsindeks, ikke en potens.
  \item $\mathcal P_k$ er polynomene av grad høyst $k$, $C^k$ er de $k$ ganger kontinuerlig deriverbare funksjonene, og $\mathcal S_{k,\Delta}$ er splinene av grad $k$ med knutevektor $\Delta$.
\end{itemize}

\subsection{Rettelser til lysarkene}\label{sec:rettelser}
Disse feilene i lysarkene er rettet i kompendiet. Alle er kontrollert med regning eller kode.
\begin{center}
\small
\begin{tabularx}{\textwidth}{@{}l X X@{}}
\toprule
Hvor & På lysarket & Riktig \\
\midrule
F1, eksempel 2 & $A=\begin{psmallmatrix}1&3\\2&2\\3&3\end{psmallmatrix}$ & $A=\begin{psmallmatrix}1&3\\2&2\\3&1\end{psmallmatrix}$ (først da stemmer den oppgitte SVD-en) \\
F1, oppsummering & «$m>n$: minst norm blant alle løsninger» & gjelder $m<n$ \\
F6, translasjon & $c_k(f(\cdot-x_0))=\ee^{+\ii kx_0}c_k(f)$ & $c_k(f(\cdot-x_0))=\ee^{-\ii kx_0}c_k(f)$ \\
F7, aliasing & $f(x)=\sum c_k\ee^{2\pi\ii kx}$ & $f(x)=\sum c_k\ee^{\ii kx}$ ($2\pi$-periodisk) \\
F8, gaussisk & $\Four(\ee^{-ax^2})=\frac{1}{\sqrt{2a}}\ee^{-\omega^2/(2a)}$ & $\frac{1}{\sqrt{2a}}\ee^{-\omega^2/(4a)}$ (som på formelarket) \\
F10, Robin-rand, venstre & $U_0^{n+1}=U_0^n+2\alpha\bigl(U_1^n-U_0^n-\tfrac{hb_0}{a_0}U_0^n-\tfrac h{a_0}g_0\bigr)$ & $+\tfrac{hb_0}{a_0}U_0^n$ (resten er likt) \\
F10, implisitt Euler, Neumann & $\vb a^{n+1}=(g_0,0,\dots,0,-g_L)^\T$ & $\vb a^{n+1}=(-g_0,0,\dots,0,g_L)^\T$ \\
F10, implisitt Euler, Dirichlet & $\vb b^{n+1}$ med $g_0(t_{n-1})$, $g_L(t_{n-1})$ & $g_0(t_{n+1})$, $g_L(t_{n+1})$ \\
F11, kvadratisk B-spline & på $[x_{i+1},x_{i+2})$ står bare det første leddet & to ledd, se F11 \\
F11, Bézier & $\vb b'(1)=-n(\vb p_{n-1}-\vb p_n)\vb p_n$, $\vb b''(1)=n(n-1)(\vb p_n-2\vb p_{n-1}+\vb p_{n-1})$ & $\vb b'(1)=n(\vb p_n-\vb p_{n-1})$, $\vb b''(1)=n(n-1)(\vb p_n-2\vb p_{n-1}+\vb p_{n-2})$ \\
\bottomrule
\end{tabularx}
\end{center}
De to F10-rettelsene er testet med en løsning vi kjenner eksakt. Med lysarkenes fortegn konvergerer ikke skjemaet (feilen holder seg på 0,28 og 1,24 når $h$ halveres), med de rettede fortegnene får feilen orden $O(h^2)$.

\clearpage

\section{F1 · Singulærverdidekomposisjon og Moore--Penrose-inversen}

\textbf{Kjernen:} Enhver matrise kan skrives som rotasjon/speiling, så skalering langs aksene, så en ny rotasjon/speiling. Fra dette følger rang, normer, den beste lavrangs-tilnærmingen og en generalisert invers som løser både overbestemte og underbestemte systemer.

\subsection{Motivasjon og forkunnskaper}
\begin{itemize}
  \item \textbf{Lav rang sparer plass og regning:} Hvis $A\approx BC$ med $B\in\R^{n\times k}$, $C\in\R^{k\times n}$ og $k\ll n$, trengs $2nk$ tall i stedet for $n^2$, og $A\vb x=B(C\vb x)$ koster $2nk$ operasjoner i stedet for $n^2$.
  \item \textbf{Egenverdier (repetisjon):} $A\vb v=\lambda\vb v$. Løs $\det(A-\lambda I)=0$ og deretter $(A-\lambda_i I)\vb v_i=\vb 0$. Hvis $A$ er diagonaliserbar, er $A=V\Lambda V^{-1}$. Dette krever en kvadratisk matrise, og noen kvadratiske matriser, for eksempel $\begin{psmallmatrix}1&c\\0&1\end{psmallmatrix}$ med $c\neq 0$, har bare én egenvektor. SVD finnes for alle matriser.
\end{itemize}

\begin{defn}[ortogonal matrise]
$U\in\R^{n\times n}$ er ortogonal hvis $U^\T U=I_n$. Det betyr at kolonnene $\vb u_i$ er ortonormale: $\ip{\vb u_i}{\vb u_j}=1$ hvis $i=j$, ellers $0$. Geometrisk er $U$ en rotasjon og/eller speiling, og den bevarer lengder: $\norm{U\vb x}_2=\norm{\vb x}_2$.
\end{defn}

\subsection{Singulærverdidekomposisjonen}
\begin{teorem}[SVD]
For enhver $A\in\R^{m\times n}$ finnes ortogonale $U\in\R^{m\times m}$ og $V\in\R^{n\times n}$ og tall $\sigma_1\ge\sigma_2\ge\dots\ge\sigma_p\ge 0$, der $p=\min\{m,n\}$, slik at
\[
A=U\Sigma V^\T,\qquad \Sigma=\diag(\sigma_1,\dots,\sigma_p)\in\R^{m\times n}\ \text{(fylt ut med nuller)}.
\]
$\sigma_k$ er \emph{singulærverdiene}, kolonnene $\vb u_k$ i $U$ er de \emph{venstre} og kolonnene $\vb v_k$ i $V$ de \emph{høyre singulærvektorene}.
\end{teorem}

\begin{defn}[redusert SVD]
Kast kolonnene som bare ganges med nuller i $\Sigma$:
\begin{itemize}
  \item $m>n$ (høy matrise): $A=\hat U\hat\Sigma V^\T$ med $\hat U\in\R^{m\times n}$ og $\hat\Sigma\in\R^{n\times n}$.
  \item $m<n$ (bred matrise): $A=U\hat\Sigma\hat V^\T$ med $\hat V\in\R^{n\times m}$ og $\hat\Sigma\in\R^{m\times m}$.
\end{itemize}
Da er $\hat U^\T\hat U=I$ (eller $\hat V^\T\hat V=I$), men $\hat U\hat U^\T$ er bare en projeksjon. Forelesningene skriver ofte $U,\Sigma,V$ uten hatt også i redusert form.
\end{defn}

\begin{teorem}[egenskaper]
\begin{enumerate}
  \item $A^\T=V\Sigma^\T U^\T$, så $A$ og $A^\T$ har de samme singulærverdiene, og $\norm{A}_2=\norm{A^\T}_2=\sigma_1$.
  \item $A^\T A=V\Sigma^\T\Sigma V^\T$ og $AA^\T=U\Sigma\Sigma^\T U^\T$. Begge er positivt semidefinite, $\sigma_k^2$ er egenverdiene deres, $\vb v_k$ er egenvektorer for $A^\T A$ og $\vb u_k$ egenvektorer for $AA^\T$.
  \item Hvis $A=A^\T$, er $\sigma_i=\abs{\lambda_i}$. Hvis $A$ i tillegg er positivt definitt, er $\sigma_i=\lambda_i$.
  \item \textbf{Geometri:} $A$ avbilder enhetssfæren på en ellipse(-oide) med halvakser $\sigma_k\vb u_k$. For $m>n$ ligger den $n$-dimensjonale ellipsoiden inne i $\R^m$, for $m<n$ projiserer $V^\T$ i tillegg.
\end{enumerate}
\end{teorem}

\begin{oppskrift}[regn ut SVD for hånd]
\begin{enumerate}
  \item Velg den minste av $A^\T A$ ($n\times n$) og $AA^\T$ ($m\times m$). Hvis bare singulærverdiene spørres etter, er du ferdig etter steg 2.
  \item Finn egenverdiene $\lambda_1\ge\lambda_2\ge\dots$ og sett $\sigma_k=\sqrt{\lambda_k}$.
  \item Finn normerte egenvektorer. Med $A^\T A$ er de $\vb v_k$, med $AA^\T$ er de $\vb u_k$.
  \item Finn de andre vektorene fra $AV=U\Sigma$ (eller $A^\T U=V\Sigma^\T$):
  \[
  \vb u_k=\frac{1}{\sigma_k}A\vb v_k\qquad\text{eller}\qquad \vb v_k=\frac{1}{\sigma_k}A^\T\vb u_k\qquad(\sigma_k>0).
  \]
  \item Trengs full $U$ (eller $V$)? Utvid med ortonormale vektorer, for eksempel fra $\kerop(A^\T)$ (eller $\kerop(A)$), slik at matrisen blir ortogonal.
  \item Kontroller ved å gange ut $U\Sigma V^\T$ eller i det minste én kolonne.
\end{enumerate}
\end{oppskrift}

\begin{fallgruve}
Ikke regn ut $U$ og $V$ hver for seg som egenvektorer av $AA^\T$ og $A^\T A$. Egenvektorene er bare bestemt opp til fortegn, så $U\Sigma V^\T$ kan bli feil. Regn ut den ene og bruk $\vb u_k=A\vb v_k/\sigma_k$.
\end{fallgruve}

\begin{eksempel}[eksempel 2 fra forelesningen, med rettet matrise]
$A=\begin{psmallmatrix}1&3\\2&2\\3&1\end{psmallmatrix}$. Da er $A^\T A=\begin{psmallmatrix}14&10\\10&14\end{psmallmatrix}$, med egenverdier $24$ og $4$ og egenvektorer $\vb v_1=\tfrac{1}{\sqrt2}(1,1)^\T$ og $\vb v_2=\tfrac1{\sqrt2}(1,-1)^\T$. Altså er $\sigma_1=2\sqrt6$ og $\sigma_2=2$, og
\[
\vb u_1=\tfrac{A\vb v_1}{\sigma_1}=\tfrac{1}{\sqrt3}(1,1,1)^\T,\qquad \vb u_2=\tfrac{A\vb v_2}{\sigma_2}=\tfrac1{\sqrt2}(-1,0,1)^\T,
\]
\[
A=\begin{pmatrix}\frac1{\sqrt3}&-\frac1{\sqrt2}\\[2pt] \frac1{\sqrt3}&0\\[2pt] \frac1{\sqrt3}&\frac1{\sqrt2}\end{pmatrix}
\begin{pmatrix}2\sqrt6&0\\0&2\end{pmatrix}\frac1{\sqrt2}\begin{pmatrix}1&1\\1&-1\end{pmatrix}.
\]
Pseudoinversen blir $A^\dagger=\frac1{\sigma_1}\vb v_1\vb u_1^\T+\frac1{\sigma_2}\vb v_2\vb u_2^\T=\frac1{12}\begin{psmallmatrix}-2&1&4\\4&1&-2\end{psmallmatrix}$, og $A^\dagger A=I_2$ (kontrollregnet).
\end{eksempel}

\subsection{SVD som sum av rang-1-matriser}
\begin{teorem}
Med $r=\max\{k:\sigma_k>0\}$ er
\[
A=\sum_{k=1}^{r}\sigma_k\,\vb u_k\vb v_k^\T,\qquad A\vb x=\sum_{k=1}^{r}\sigma_k\ip{\vb v_k}{\vb x}\,\vb u_k .
\]
Hvert ledd $\vb u_k\vb v_k^\T$ har rang 1. Det følger at
\begin{itemize}
  \item $\rank A=r$ (antall singulærverdier som ikke er null),
  \item $\ran A=\spn\{\vb u_1,\dots,\vb u_r\}$ og $\kerop A=\spn\{\vb v_1,\dots,\vb v_r\}^\perp=\spn\{\vb v_{r+1},\dots,\vb v_n\}$,
  \item $\norm A_2=\sigma_1$ og $\norm A_F=\bigl(\sum_{i,j}a_{ij}^2\bigr)^{1/2}=\bigl(\sigma_1^2+\dots+\sigma_r^2\bigr)^{1/2}$.
\end{itemize}
\end{teorem}

\begin{teorem}[beste tilnærming med lav rang (Eckart--Young)]
Den \emph{avkortede SVD-en} $A_k=\sum_{j=1}^{k}\sigma_j\vb u_j\vb v_j^\T$ ($1\le k\le p$) løser både
$\min_{\rank B\le k}\norm{A-B}_2$ og $\min_{\rank B\le k}\norm{A-B}_F$, og feilene er
\[
\norm{A-A_k}_2=\sigma_{k+1},\qquad \norm{A-A_k}_F=\Bigl(\sum_{j=k+1}^{r}\sigma_j^2\Bigr)^{1/2}.
\]
Å lagre $A_k$ krever $k(m+n+1)$ tall, mot $mn$ for hele $A$.
\end{teorem}

\subsection{Moore--Penrose-inversen (pseudoinversen)}
\begin{defn}[pseudoinvers]
Med SVD-en $A=\sum_{k=1}^r\sigma_k\vb u_k\vb v_k^\T$ er
\[
A^\dagger=\sum_{k=1}^{r}\frac{1}{\sigma_k}\vb v_k\vb u_k^\T=V\Sigma^{+}U^\T\in\R^{n\times m},
\]
der $\Sigma^+$ er $\Sigma^\T$ med hver $\sigma_k>0$ byttet ut med $1/\sigma_k$ (nuller beholdes). Hvis $A$ har full rang, er $A^\dagger=V\Sigma^{-1}U^\T$ i redusert form. Merk at $\vb u$ og $\vb v$ bytter plass.
\end{defn}

\begin{teorem}[de tre tilfellene]
\begin{center}
\small
\begin{tabularx}{\textwidth}{@{}l X X@{}}
\toprule
& Høy, full rang ($m>n=r$) & Bred, full rang ($r=m<n$) \\
\midrule
Identitet & $A^\dagger A=I_n$, \ $AA^\dagger=UU^\T$ (projeksjon på $\ran A$) & $AA^\dagger=I_m$, \ $A^\dagger A=VV^\T$ (projeksjon på $(\kerop A)^\perp$) \\
Formel & $A^\dagger=(A^\T A)^{-1}A^\T$ & $A^\dagger=A^\T(AA^\T)^{-1}$ \\
$\vb x=A^\dagger\vb b$ løser & $\min_{\vb x}\norm{A\vb x-\vb b}_2^2$ (minste kvadrater) & $\min\norm{\vb x}_2^2$ slik at $A\vb x=\vb b$ (minst norm) \\
\bottomrule
\end{tabularx}
\end{center}
\textbf{Generelt} (vilkårlig rang) løser $\vb x^\dagger=A^\dagger\vb b$ et problem i to nivåer: blant alle $\vb x$ som minimerer $\norm{A\vb x-\vb b}_2$, har $\vb x^\dagger$ minst norm.
\end{teorem}

\begin{teorem}[Penrose-likningene]
For alle $A$ gjelder $(A^\dagger)^\dagger=A$, \ $A^\dagger AA^\dagger=A^\dagger$ og \ $AA^\dagger A=A$. I tillegg er $AA^\dagger$ og $A^\dagger A$ symmetriske. Disse fire egenskapene bestemmer $A^\dagger$ entydig. Det brukes for å sjekke en oppgitt kandidat $Q$: vis $AQA=A$, $QAQ=Q$, $(AQ)^\T=AQ$ og $(QA)^\T=QA$. Videre er $(A^\T)^\dagger=(A^\dagger)^\T$.
\end{teorem}

\begin{oppskrift}[minste kvadrater og minst norm med SVD]
\begin{enumerate}
  \item Regn ut den reduserte SVD-en og $A^\dagger=\sum_k\sigma_k^{-1}\vb v_k\vb u_k^\T$.
  \item $\vb x=A^\dagger\vb b$.
  \item Feilen er $\norm{A\vb x-\vb b}_2$. Den er $0$ hvis $\vb b\in\ran A$, ellers er den lengden av komponenten av $\vb b$ som står vinkelrett på $\spn\{\vb u_1,\dots,\vb u_r\}$.
  \item For et underbestemt system har $A^\dagger\vb b$ minst norm blant uendelig mange løsninger. Tolk og sammenlign gjerne med en annen løsning.
\end{enumerate}
\end{oppskrift}

\subsection{Videre: støy og avkortet pseudoinvers}
Hvis $\vb b$ har støy ($\vb b^\delta=\vb b+\vb n^\delta$), blåser små $\sigma_j$ opp støyen gjennom $1/\sigma_j$. Med $\vb x^{(k)}_\delta=A_k^\dagger\vb b^\delta$ blir datatilpasningen $\norm{A\vb x^{(k)}-\vb b^\delta}$ først mye bedre etter hvert som $k$ øker, så flater den ut, og deretter vokser stort sett bare $\norm{\vb x^{(k)}_\delta}$. Å finne balansepunktet kalles L-kurve-metoden.

\begin{oving}
\begin{itemize}[leftmargin=5.6em,labelwidth=5em,labelsep=0.6em,align=left]
\item[\lett] \href{https://github.com/gustavbv/Matte4C---eksamen/blob/main/oppgavesett/F01-svd-og-moore-penrose/1-TMA4432-ovinger-og-proveeksamen/002-TMA4432-oving-1-oppg1-oppgave.pdf}{Øving 1, oppg.~1} {\color{graa}(\href{https://github.com/gustavbv/Matte4C---eksamen/blob/main/oppgavesett/F01-svd-og-moore-penrose/1-TMA4432-ovinger-og-proveeksamen/002-TMA4432-oving-1-oppg1-svar.pdf}{LF})}: SVD av en $2\times2$- og en $3\times2$-matrise
\item[\lett] \href{https://github.com/gustavbv/Matte4C---eksamen/blob/main/oppgavesett/F01-svd-og-moore-penrose/1-TMA4432-ovinger-og-proveeksamen/003-TMA4432-oving-1-oppg2-oppgave.pdf}{Øving 1, oppg.~2} {\color{graa}(\href{https://github.com/gustavbv/Matte4C---eksamen/blob/main/oppgavesett/F01-svd-og-moore-penrose/1-TMA4432-ovinger-og-proveeksamen/003-TMA4432-oving-1-oppg2-svar.pdf}{LF})}: pseudoinvers av en matrise uten full rang
\item[\middels] \href{https://github.com/gustavbv/Matte4C---eksamen/blob/main/oppgavesett/F01-svd-og-moore-penrose/1-TMA4432-ovinger-og-proveeksamen/005-TMA4432-oving-2-oppg1-oppgave.pdf}{Øving 2, oppg.~1} {\color{graa}(\href{https://github.com/gustavbv/Matte4C---eksamen/blob/main/oppgavesett/F01-svd-og-moore-penrose/1-TMA4432-ovinger-og-proveeksamen/005-TMA4432-oving-2-oppg1-svar.pdf}{LF})}: redusert SVD og $A^\dagger$ for en $4\times2$-matrise
\item[\middels] \href{https://github.com/gustavbv/Matte4C---eksamen/blob/main/oppgavesett/F01-svd-og-moore-penrose/3-hele-oppgaven-i-pensum/009-TMA4145-2024-host-oppg3-oppgave.pdf}{TMA4145 høst 2024, oppg.~3} {\color{graa}(\href{https://github.com/gustavbv/Matte4C---eksamen/blob/main/oppgavesett/F01-svd-og-moore-penrose/3-hele-oppgaven-i-pensum/009-TMA4145-2024-host-oppg3-svar.pdf}{LF})}: SVD av en $3\times2$-matrise og minste kvadrater
\item[\middels] \href{https://github.com/gustavbv/Matte4C---eksamen/blob/main/oppgavesett/F01-svd-og-moore-penrose/1-TMA4432-ovinger-og-proveeksamen/001-TMA4432-2026-proveeksamen-oppg1-oppgave.pdf}{Prøveeksamen, oppg.~1} {\color{graa}(\href{https://github.com/gustavbv/Matte4C---eksamen/blob/main/oppgavesett/F01-svd-og-moore-penrose/1-TMA4432-ovinger-og-proveeksamen/001-TMA4432-2026-proveeksamen-oppg1-svar.pdf}{LF})}: SVD, $A^\dagger$, minste feil og løsningen med minst norm
\item[\vanskelig] \href{https://github.com/gustavbv/Matte4C---eksamen/blob/main/oppgavesett/F01-svd-og-moore-penrose/3-hele-oppgaven-i-pensum/019-TMA4205-2019-host-oppg4-oppgave.pdf}{TMA4205 høst 2019, oppg.~4} {\color{graa}(\href{https://github.com/gustavbv/Matte4C---eksamen/blob/main/oppgavesett/F01-svd-og-moore-penrose/3-hele-oppgaven-i-pensum/019-TMA4205-2019-host-oppg4-svar.pdf}{LF})}: beste rang-1-approksimasjon og en Frobenius-skranke
\item[\vanskelig] \href{https://github.com/gustavbv/Matte4C---eksamen/blob/main/oppgavesett/F01-svd-og-moore-penrose/3-hele-oppgaven-i-pensum/032-TMA4205-2009-host-oppg1-oppgave.pdf}{TMA4205 høst 2009, oppg.~1} {\color{graa}(\href{https://github.com/gustavbv/Matte4C---eksamen/blob/main/oppgavesett/F01-svd-og-moore-penrose/3-hele-oppgaven-i-pensum/032-TMA4205-2009-host-oppg1-svar.pdf}{LF})}: fullføre en SVD, beste rang-2-approksimasjon
\item[\vanskelig] \href{https://github.com/gustavbv/Matte4C---eksamen/blob/main/oppgavesett/F01-svd-og-moore-penrose/3-hele-oppgaven-i-pensum/017-TMA4205-2022-host-oppg7-oppgave.pdf}{TMA4205 høst 2022, oppg.~7} {\color{graa}(\href{https://github.com/gustavbv/Matte4C---eksamen/blob/main/oppgavesett/F01-svd-og-moore-penrose/3-hele-oppgaven-i-pensum/017-TMA4205-2022-host-oppg7-svar.pdf}{LF})}: bevis for at avkortet SVD er best i Frobeniusnorm
\end{itemize}
\end{oving}


\section{F2 · Tensorer og tensordekomposisjoner}

\textbf{Kjernen:} Data med flere enn to indekser (fargebilder, video) er tensorer. Vi vil finne en tensorversjon av lavrangs-tilnærmingen fra SVD. Det gir to generaliseringer: CP (sum av rang-1-tensorer) og Tucker (kjernetensor ganget med faktormatriser i hver modus).
\textbf{På eksamen:} forståelse, definisjoner og små regneoperasjoner. Du blir ikke bedt om å regne ut en konkret dekomposisjon.

\subsection{Definisjoner og indeksering}
\begin{defn}[tensor]
For $n_1,\dots,n_d\in\N$ er en tensor $\tn X\in\R^{n_1\times\dots\times n_d}$ en samling tall $x_{i_1,\dots,i_d}$ med $1\le i_k\le n_k$. Tallet $d$ er tensorens \emph{orden}. Orden 0 er et tall, orden 1 en vektor og orden 2 en matrise. Kurset ser mest på $d=3$: $\tn X\in\R^{m\times n\times p}$ med elementer $x_{ijk}$.
\end{defn}

\begin{defn}[fibre og snitt for $d=3$]
\textbf{Fibre} får du ved å holde to indekser fast: kolonnefibre $\vb x_{:jk}\in\R^m$, radfibre $\vb x_{i:k}\in\R^n$ og rørfibre (tube) $\vb x_{ij:}\in\R^p$.\\
\textbf{Snitt} får du ved å holde én indeks fast: frontale snitt $X_{::k}\in\R^{m\times n}$, laterale snitt $X_{:j:}\in\R^{m\times p}$ og horisontale snitt $X_{i::}\in\R^{n\times p}$.
\end{defn}

\begin{defn}[ytre produkt og rang]
Det ytre produktet $\tn X=\vb a^{(1)}\outp\vb a^{(2)}\outp\dots\outp\vb a^{(d)}$ har elementer $x_{i_1\dots i_d}=a^{(1)}_{i_1}a^{(2)}_{i_2}\cdots a^{(d)}_{i_d}$. For $d=2$ er dette $\vb u\vb v^\T$, og for $d=3$ er $(\vb a\outp\vb b\outp\vb c)_{ijk}=a_ib_jc_k$. Slike tensorer har \emph{rang 1}.\\
\textbf{Rangen} til $\tn X$ er det minste $r$ slik at $\tn X=\sum_{s=1}^r\lambda_s\,\vb a^{(1)}_s\outp\dots\outp\vb a^{(d)}_s$.
\end{defn}

\begin{defn}[Kruskal-notasjon]
Samle vektorene for hver modus i matriser $A=(\vb a_1\cdots\vb a_r)\in\R^{m\times r}$, $B\in\R^{n\times r}$ og $C\in\R^{p\times r}$, og skriv
\[
\kr{\boldsymbol\lambda;A,B,C}=\sum_{s=1}^r\lambda_s\,\vb a_s\outp\vb b_s\outp\vb c_s,\qquad x_{ijk}=\sum_{s=1}^r\lambda_s a_{is}b_{js}c_{ks}.
\]
Uten $\boldsymbol\lambda$ er alle $\lambda_s=1$. SVD-en er $A=\kr{\boldsymbol\sigma;U,V}$.
\end{defn}

\subsection{Matriseoperasjoner}
\begin{defn}[vektorisering, Kronecker, Khatri--Rao og Hadamard]
\begin{itemize}
  \item $\vect(A)\in\R^{mn}$: kolonnene i $A$ stablet oppå hverandre.
  \item \textbf{Kronecker-produktet} av $A\in\R^{m\times n}$ og $B\in\R^{p\times q}$ er
  $A\otimes B=\begin{psmallmatrix}a_{11}B&\cdots&a_{1n}B\\ \vdots&&\vdots\\ a_{m1}B&\cdots&a_{mn}B\end{psmallmatrix}\in\R^{mp\times nq}$.
  For vektorer er $\vb a\otimes\vb b=(a_1\vb b;\dots;a_m\vb b)\in\R^{mp}$, og $A\otimes B=(\vb a_1\otimes\vb b_1\ \ \vb a_1\otimes\vb b_2\ \cdots\ \vb a_n\otimes\vb b_q)$. Produktet er alltid definert.
  \item \textbf{Khatri--Rao-produktet} krever like mange kolonner: for $A\in\R^{m\times n}$ og $B\in\R^{p\times n}$ er
  $A\odot B=(\vb a_1\otimes\vb b_1\ \ \vb a_2\otimes\vb b_2\ \cdots\ \vb a_n\otimes\vb b_n)\in\R^{mp\times n}$, altså kolonnevis Kronecker.
  \item \textbf{Hadamard-produktet} krever lik størrelse: $(A\had B)_{ij}=a_{ij}b_{ij}$.
\end{itemize}
\end{defn}
\begin{teorem}[nyttige identiteter (ikke på lysarkene, men brukt i oppgavebanken)]
$(A\otimes B)(C\otimes D)=AC\otimes BD$ (når produktene er definert), $(A\otimes B)^\T=A^\T\otimes B^\T$, $\vect(AXB)=(B^\T\otimes A)\vect(X)$ og $(A\odot B)^\T(A\odot B)=A^\T A\had B^\T B$. Den siste gjør minste kvadrater-steget i ALS billig.
\end{teorem}
\begin{fallgruve}[når er produktene definert? (øving 2)]
For $C\in\R^{m\times n}$: $CC$ krever $m=n$, mens $C^\T C\in\R^{n\times n}$ alltid er definert. $C\had C$ er alltid definert, men $C^\T\had C$ krever $m=n$. $A\odot B$ krever like mange kolonner, og $A\had B$ lik størrelse.
\end{fallgruve}

\subsection{Tensoroperasjoner}
\begin{defn}[matrisering (modus-$k$-utfolding)]
$X_{(k)}\in\R^{n_k\times\prod_{l\neq k}n_l}$: modus-$k$-fibrene blir kolonner. Elementet $x_{i_1\dots i_d}$ havner i rad $i_k$ og kolonne
\[
j=1+\sum_{l\neq k}(i_l-1)J_l,\qquad J_l=\prod_{\substack{q<l\\ q\neq k}}n_q,
\]
altså går de andre indeksene i rekkefølge med den \emph{første indeksen raskest}.
\end{defn}
\begin{eksempel}[$\tn X\in\R^{3\times4\times2}$ med tallene $1,\dots,24$ (fra forelesningen, kontrollregnet)]
$X_{::1}=\begin{psmallmatrix}1&4&7&10\\2&5&8&11\\3&6&9&12\end{psmallmatrix}$ og $X_{::2}=\begin{psmallmatrix}13&16&19&22\\14&17&20&23\\15&18&21&24\end{psmallmatrix}$ gir
\[
X_{(1)}=\begin{psmallmatrix}1&4&7&10&13&16&19&22\\2&5&8&11&14&17&20&23\\3&6&9&12&15&18&21&24\end{psmallmatrix},\quad
X_{(2)}=\begin{psmallmatrix}1&2&3&13&14&15\\4&5&6&16&17&18\\7&8&9&19&20&21\\10&11&12&22&23&24\end{psmallmatrix},\quad
X_{(3)}=\begin{psmallmatrix}1&2&3&\cdots&12\\13&14&15&\cdots&24\end{psmallmatrix}.
\]
\end{eksempel}

\begin{defn}[$k$-modusprodukt]
For $\tn X\in\R^{n_1\times\dots\times n_d}$ og $A\in\R^{m\times n_k}$ er $\tn Y=\tn X\times_k A\in\R^{n_1\times\dots\times n_{k-1}\times m\times n_{k+1}\times\dots\times n_d}$ gitt ved
\[
y_{i_1\dots i_{k-1}\,j\,i_{k+1}\dots i_d}=\sum_{i_k=1}^{n_k}x_{i_1\dots i_d}\,a_{j\,i_k},\qquad\text{kort: }\ Y_{(k)}=A\,X_{(k)}.
\]
Hver modus-$k$-fiber ganges altså med $A$. For en matrise $B$ er $B\times_1A=AB$ og $B\times_2A=BA^\T$. Med en vektor $\vb a\in\R^{n_k}$ forsvinner modus $k$, så ordenen synker med én: $B\times_1\vb v=B^\T\vb v$ og $B\times_2\vb v=B\vb v$.
\end{defn}

\begin{defn}[multilineær transformasjon]
For $\tn X=\kr{A,B,C}$ og $M\in\R^{m\times u}$, $N\in\R^{n\times v}$, $P\in\R^{p\times w}$ er $\tn X(M,N,P)=\sum_s(M^\T\vb a_s)\outp(N^\T\vb b_s)\outp(P^\T\vb c_s)=\kr{M^\T A,N^\T B,P^\T C}$. For matriser er $\kr{A,B}(M,N)=M^\T AB^\T N$. Hvis en ny dimensjon er 1 (for eksempel $w=1$), forsvinner den modusen.
\end{defn}

\begin{teorem}[matrisering av en CP-tensor]
For $\tn X=\kr{A,B,C}$ er
\[
X_{(1)}=A(C\odot B)^\T,\qquad X_{(2)}=B(C\odot A)^\T,\qquad X_{(3)}=C(B\odot A)^\T .
\]
Rekkefølgen i Khatri--Rao-produktet er de to andre modusene \emph{baklengs}.
\end{teorem}

\subsection{CP-dekomposisjon og ALS}
\begin{defn}[CP (CANDECOMP/PARAFAC)]
Med Frobeniusnormen $\norm{\tn X}_F=\bigl(\sum_{i,j,k}x_{ijk}^2\bigr)^{1/2}$ er CP-dekomposisjonen med rang $r$ problemet
\[
\min_{A,B,C}\ \norm{\tn X-\kr{A,B,C}}_F^2,\qquad A\in\R^{m\times r},\ B\in\R^{n\times r},\ C\in\R^{p\times r}.
\]
\end{defn}
\begin{oppskrift}[alternerende minste kvadrater (ALS)]
Hvert delproblem er lineært minste kvadrater i én faktor, så Moore--Penrose (F1) gir løsningen direkte. Start med $A^{(0)},B^{(0)},C^{(0)}$ (for eksempel tilfeldige) og gjenta til et stoppkriterium er oppfylt:
\begin{align*}
A^{(k+1)}&=X_{(1)}\bigl((C^{(k)}\odot B^{(k)})^\T\bigr)^\dagger,\\
B^{(k+1)}&=X_{(2)}\bigl((C^{(k)}\odot A^{(k+1)})^\T\bigr)^\dagger,\\
C^{(k+1)}&=X_{(3)}\bigl((B^{(k+1)}\odot A^{(k+1)})^\T\bigr)^\dagger.
\end{align*}
Hver faktor oppdateres med de nyeste verdiene av de andre. Hvert steg løser $\min_{\tilde A}\norm{X_{(1)}-\tilde A(C\odot B)^\T}_F$.
\end{oppskrift}

\subsection{Tucker-dekomposisjon, HOSVD og HOOI}
\begin{defn}[Tucker]
For gitte ranger $(r_1,r_2,r_3)$:
\[
\min\ \norm{\tn X-\tn Y}_F,\qquad \tn Y=\tn G\times_1A\times_2B\times_3C=\sum_{i,j,k}g_{ijk}\,\vb a_i\outp\vb b_j\outp\vb c_k,
\]
med en \emph{kjernetensor} $\tn G\in\R^{r_1\times r_2\times r_3}$ og faktorer $A\in\R^{m\times r_1}$, $B\in\R^{n\times r_2}$ og $C\in\R^{p\times r_3}$. Matriseringene er $Y_{(1)}=AG_{(1)}(C\otimes B)^\T$, $Y_{(2)}=BG_{(2)}(C\otimes A)^\T$ og $Y_{(3)}=CG_{(3)}(B\otimes A)^\T$ (Kronecker, ikke Khatri--Rao).
\end{defn}
\begin{teorem}[hvordan henger CP, Tucker og SVD sammen?]
\begin{itemize}
  \item CP er et spesialtilfelle av Tucker med $r_1=r_2=r_3=r$ og en superdiagonal kjerne ($g_{ijk}=1$ hvis $i=j=k$, ellers $0$).
  \item SVD-en kan skrives på begge måter: $U\Sigma V^\T=\sum_k\sigma_k\vb u_k\vb v_k^\T$ (CP-form, sum av rang-1) og $U\Sigma V^\T=\Sigma\times_1U\times_2V$ (Tucker-form med kjerne $\Sigma$). For matriser er formene like fordi $\Sigma$ er diagonal. For tensorer er kjernen $\tn G$ som regel ikke superdiagonal, så formene skiller seg.
  \item Lagring: CP trenger $r(m+n+p)$ tall og Tucker $r_1r_2r_3+mr_1+nr_2+pr_3$, mot $mnp$ for hele tensoren.
\end{itemize}
\end{teorem}
\begin{oppskrift}[HOSVD (higher-order SVD)]
For gitte $\tn X$ og $(r_1,r_2,r_3)$:
\begin{enumerate}
  \item $A$ er de $r_1$ ledende venstre singulærvektorene til $X_{(1)}$.
  \item $B$ er de $r_2$ ledende venstre singulærvektorene til $X_{(2)}$.
  \item $C$ er de $r_3$ ledende venstre singulærvektorene til $X_{(3)}$.
  \item $\tn G=\tn X\times_1A^\T\times_2B^\T\times_3C^\T$. Dette tilsvarer $\Sigma=U^\T AV=A\times_1U^\T\times_2V^\T$ for matriser.
\end{enumerate}
\textbf{HOOI} (videre): start fra HOSVD og gjenta. Beregn $\tn Y=\tn X\times_2B^\T\times_3C^\T$ og la $A$ være de $r_1$ ledende venstre singulærvektorene til $Y_{(1)}$. Gjør det samme for $B$ (projiser med $A$ og $C$) og $C$ (projiser med $A$ og $B$). Til slutt er $\tn G=\tn X\times_1A^\T\times_2B^\T\times_3C^\T$.
\end{oppskrift}

\begin{oving}
\begin{itemize}[leftmargin=5.6em,labelwidth=5em,labelsep=0.6em,align=left]
\item[\lett] \href{https://github.com/gustavbv/Matte4C---eksamen/blob/main/oppgavesett/F02-tensorer/3-hele-oppgaven-i-pensum/004-BITS-MATH-F425-2023-var-oppg6-oppgave.pdf}{BITS Pilani 2023, oppg.~6} {\color{graa}(\href{https://github.com/gustavbv/Matte4C---eksamen/blob/main/oppgavesett/F02-tensorer/3-hele-oppgaven-i-pensum/004-BITS-MATH-F425-2023-var-oppg6-svar.pdf}{LF})}: 3-modusprodukt for en liten tensor
\item[\lett] \href{https://github.com/gustavbv/Matte4C---eksamen/blob/main/oppgavesett/F02-tensorer/3b-ovingsark-andre-universiteter/011-MPI-2017-ark2-oppg3-oppgave.pdf}{MPI 2017, ark 2, oppg.~3} {\color{graa}(\href{https://github.com/gustavbv/Matte4C---eksamen/blob/main/oppgavesett/F02-tensorer/3b-ovingsark-andre-universiteter/011-MPI-2017-ark2-oppg3-svar.pdf}{LF})}: rang-1-tensorer og $\kr{A,B,C}$
\item[\middels] \href{https://github.com/gustavbv/Matte4C---eksamen/blob/main/oppgavesett/F02-tensorer/1-TMA4432-ovinger-og-proveeksamen/002-TMA4432-oving-2-oppg2-oppgave.pdf}{Øving 2, oppg.~2} {\color{graa}(\href{https://github.com/gustavbv/Matte4C---eksamen/blob/main/oppgavesett/F02-tensorer/1-TMA4432-ovinger-og-proveeksamen/002-TMA4432-oving-2-oppg2-svar.pdf}{LF})}: Kronecker-, Khatri--Rao- og Hadamard-produkt
\item[\middels] \href{https://github.com/gustavbv/Matte4C---eksamen/blob/main/oppgavesett/F02-tensorer/3b-ovingsark-andre-universiteter/005-MPI-2017-ark1-oppg1-oppgave.pdf}{MPI 2017, ark 1, oppg.~1} {\color{graa}(\href{https://github.com/gustavbv/Matte4C---eksamen/blob/main/oppgavesett/F02-tensorer/3b-ovingsark-andre-universiteter/005-MPI-2017-ark1-oppg1-svar.pdf}{LF})}: matriseringene $X_{(1)},X_{(2)},X_{(3)}$
\item[\middels] \href{https://github.com/gustavbv/Matte4C---eksamen/blob/main/oppgavesett/F02-tensorer/3b-ovingsark-andre-universiteter/014-MPI-2017-ark4-oppg2-oppgave.pdf}{MPI 2017, ark 4, oppg.~2} {\color{graa}(\href{https://github.com/gustavbv/Matte4C---eksamen/blob/main/oppgavesett/F02-tensorer/3b-ovingsark-andre-universiteter/014-MPI-2017-ark4-oppg2-svar.pdf}{LF})}: regne ut $\tn G\times_1A\times_2B\times_3C$
\item[\middels] \href{https://github.com/gustavbv/Matte4C---eksamen/blob/main/oppgavesett/F02-tensorer/1-TMA4432-ovinger-og-proveeksamen/001-TMA4432-2026-proveeksamen-oppg2-oppgave.pdf}{Prøveeksamen, oppg.~2} {\color{graa}(\href{https://github.com/gustavbv/Matte4C---eksamen/blob/main/oppgavesett/F02-tensorer/1-TMA4432-ovinger-og-proveeksamen/001-TMA4432-2026-proveeksamen-oppg2-svar.pdf}{LF})}: Tucker, CP og SVD: hvordan henger de sammen?
\item[\vanskelig] \href{https://github.com/gustavbv/Matte4C---eksamen/blob/main/oppgavesett/F02-tensorer/3b-ovingsark-andre-universiteter/012-MPI-2017-ark3-oppg1-oppgave.pdf}{MPI 2017, ark 3, oppg.~1} {\color{graa}(\href{https://github.com/gustavbv/Matte4C---eksamen/blob/main/oppgavesett/F02-tensorer/3b-ovingsark-andre-universiteter/012-MPI-2017-ark3-oppg1-svar.pdf}{LF})}: maksimal tensorrang via $X_{(1)}=A(C\odot B)^\T$
\item[\vanskelig] \href{https://github.com/gustavbv/Matte4C---eksamen/blob/main/oppgavesett/F02-tensorer/3b-ovingsark-andre-universiteter/022-EPFL-2018-oving7-oppg1-oppgave.pdf}{EPFL 2018, øving 7, oppg.~1} {\color{graa}(\href{https://github.com/gustavbv/Matte4C---eksamen/blob/main/oppgavesett/F02-tensorer/3b-ovingsark-andre-universiteter/022-EPFL-2018-oving7-oppg1-svar.pdf}{LF})}: når er summen av to rang-1-tensorer rang 1?
\end{itemize}
\end{oving}


\section{F3 · Parametriserte kurver og flater}

\textbf{Kjernen:} En kurve er bildet av en funksjon $\vb c:I\to\R^n$ av én variabel, og en flate bildet av $\vb s:D\to\R^n$ av to variable. Derivasjon gir hastighet, tangentlinje og tangentplan. Glatthet betyr at det ikke finnes noe sted der vi «står stille».

\subsection{Kurver}
\begin{defn}[kurve og parametrisering]
En \emph{parametrisering} er en kontinuerlig $\vb c:I\to\R^n$, $I\subseteq\R$ et intervall, med koordinatfunksjoner $c_1(t),\dots,c_n(t)$. \emph{Kurven} er bildemengden (verdimengden) $C=\{\vb c(t):t\in I\}$. En kurve har mange parametriseringer, og et punkt kan besøkes flere ganger. Kontinuitet kan sjekkes komponentvis.
\end{defn}
\begin{itemize}
  \item \textbf{Rett linje:} $\vb c(t)=\vb p_0+t\vb v$ med $\vb v\neq\vb 0$, eller gjennom $\vb a$ og $\vb b$: $\vb p_0=\vb a$ og $\vb v=\vb b-\vb a$.
  \item \textbf{Standardkurver:} sirkel $(r\cos t,r\sin t)$, ellipse $\tfrac{x^2}{a^2}+\tfrac{y^2}{b^2}=1$ gir $(a\cos t,b\sin t)$, $t\in[0,2\pi]$, og spiral (heliks) $(\cos t,\sin t,t)$.
  \item \textbf{Skisse:} eliminer $t$. For eksempel gir $(t^2,4t^4)$ at $y=4x^2$ med $x\ge0$.
\end{itemize}

\begin{defn}[hastighet, fart, akselerasjon og enhetstangent]
\[
\vb v(t)=\vb c'(t)=\bigl(c_1'(t),\dots,c_n'(t)\bigr)^\T,\quad v(t)=\norm{\vb c'(t)}_2,\quad \vb a(t)=\vb v'(t)=\vb c''(t),\quad \vb T(t)=\frac{\vb v(t)}{\norm{\vb v(t)}_2}.
\]
Hastigheten er en vektor, mens farten er et tall (lengden av hastigheten).
\end{defn}

\begin{defn}[glatt kurve]
$\vb c$ er \emph{glatt i $t$} hvis (i) den er $C^1$ nær $t$, (ii) den er injektiv på et intervall rundt $t$ og (iii) $\vb c'(t)\neq\vb 0$.
Kurven $C$ er glatt i et punkt $\vb p$ hvis det finnes en parametrisering som er glatt i en $t$ med $\vb c(t)=\vb p$. $C$ er glatt hvis den er glatt i alle punkter.\\
\textbf{Stykkevis glatt:} glatt og injektiv bortsett fra i endelig mange punkter, og alle $c_i'$ er begrenset (ingen poler).
\end{defn}
\begin{eksempel}[en spiss (forelesningen)]
$\vb c(t)=(2t-t^2,\,3t-t^3)$ gir $\vb c'(t)=(2-2t,\,3-3t^2)$ og $\vb c'(1)=\vb 0$. Vi «kjører», stopper og snur, og kurven får et knekkpunkt der. Parametriseringen er ikke glatt i $t=1$. Om \emph{kurven} er glatt der, må avgjøres separat (her er den ikke det).
\end{eksempel}

\begin{defn}[tangentlinje og buelengde]
\textbf{Tangentlinjen} i $t_0$ er $\vb L(t)=\vb c(t_0)+\vb c'(t_0)(t-t_0)$. En vektor parallell med $\vb c'(t_0)$ er en tangentvektor, og med lengde 1 er den en enhetstangentvektor.\\
\textbf{Buelengden} fra $\vb c(t_0)$ til $\vb c(t)$ er $s(t)=\displaystyle\int_{t_0}^{t}\norm{\vb c'(\tau)}_2\dd\tau$.
En kurve er \emph{parametrisert ved buelengde} hvis $\norm{\vb c'(s)}_2=1$ for alle $s$ (enhetsfart). En glatt kurve kan alltid parametriseres slik.
\end{defn}
\begin{fallgruve}[buelengde i praksis]
Integranden er en kvadratrot. Eksamensoppgaver er nesten alltid laget slik at uttrykket under roten blir et kvadrat, for eksempel $1+4t^2+4t^4=(1+2t^2)^2$. Se etter et slikt kvadrat.
\end{fallgruve}

\begin{oving}
\begin{itemize}[leftmargin=5.6em,labelwidth=5em,labelsep=0.6em,align=left]
\item[\lett] \href{https://github.com/gustavbv/Matte4C---eksamen/blob/main/oppgavesett/F03-kurver-og-flater/3-hele-oppgaven-i-pensum/033-TMA4105-2017-var-oppg1-oppgave.pdf}{TMA4105 vår 2017, oppg.~1} {\color{graa}(\href{https://github.com/gustavbv/Matte4C---eksamen/blob/main/oppgavesett/F03-kurver-og-flater/3-hele-oppgaven-i-pensum/033-TMA4105-2017-var-oppg1-svar.pdf}{LF})}: buelengde
\item[\middels] \href{https://github.com/gustavbv/Matte4C---eksamen/blob/main/oppgavesett/F03-kurver-og-flater/3-hele-oppgaven-i-pensum/046-TMA4105-2015-var-oppg1-oppgave.pdf}{TMA4105 vår 2015, oppg.~1} {\color{graa}(\href{https://github.com/gustavbv/Matte4C---eksamen/blob/main/oppgavesett/F03-kurver-og-flater/3-hele-oppgaven-i-pensum/046-TMA4105-2015-var-oppg1-svar.pdf}{LF})}: enhetstangent og buelengde
\item[\middels] \href{https://github.com/gustavbv/Matte4C---eksamen/blob/main/oppgavesett/F03-kurver-og-flater/1-TMA4432-ovinger-og-proveeksamen/002-TMA4432-oving-3-oppg1-oppgave.pdf}{Øving 3, oppg.~1} {\color{graa}(\href{https://github.com/gustavbv/Matte4C---eksamen/blob/main/oppgavesett/F03-kurver-og-flater/1-TMA4432-ovinger-og-proveeksamen/002-TMA4432-oving-3-oppg1-svar.pdf}{LF})}: parametrisering, fart, akselerasjon og glatthet
\item[\vanskelig] \href{https://github.com/gustavbv/Matte4C---eksamen/blob/main/oppgavesett/F03-kurver-og-flater/3-hele-oppgaven-i-pensum/057-MA1103-2009-kont-oppg3-oppgave.pdf}{MA1103 kont.\ 2009, oppg.~3} {\color{graa}(\href{https://github.com/gustavbv/Matte4C---eksamen/blob/main/oppgavesett/F03-kurver-og-flater/3-hele-oppgaven-i-pensum/057-MA1103-2009-kont-oppg3-svar.pdf}{LF})}: skjæringskurve og bevegelse med konstant fart
\end{itemize}
\end{oving}


\subsection{Partiellderiverte}
\begin{defn}
$\pd F{x_i}(\vb y)=\displaystyle\lim_{h\to0}\frac{F(y_1,\dots,y_i+h,\dots,y_m)-F(\vb y)}{h}$: deriver med hensyn på $x_i$ og hold de andre variablene konstante.
\end{defn}

\subsection{Flater}
\begin{defn}[flate og parametrisering]
En parametrisering er en kontinuerlig $\vb s:D\to\R^n$, $D\subseteq\R^2$, og flaten er $S=\{\vb s(u,v):(u,v)\in D\}$.
\end{defn}
\begin{itemize}
  \item \textbf{Plan:} $\vb r(u,v)=\vb p_0+u\vb a+v\vb b$ med $\vb a,\vb b$ lineært uavhengige. Dette kan skrives om til $Ax_1+Bx_2+Cx_3+D=0$.
  \item \textbf{Sylinder:} $(\cos u,\sin u,v)$. \textbf{Kule} med radius $r$: $(r\cos u\sin v,\ r\sin u\sin v,\ r\cos v)$, $u\in[0,2\pi]$, $v\in[0,\pi]$. \textbf{Ellipsoide} $\tfrac{x^2}{a^2}+\tfrac{y^2}{b^2}+\tfrac{z^2}{c^2}=1$: $(a\cos u\sin v,\ b\sin u\sin v,\ c\cos v)$.
  \item \textbf{Graf} av $F:\R^2\to\R$: $\vb h(u,v)=(u,v,F(u,v))$.
\end{itemize}

\begin{defn}[hastighetsvektorer, totalderivert og glatthet]
\[
\pd{\vb s}u=\Bigl(\pd{s_1}u,\dots,\pd{s_n}u\Bigr)^\T,\qquad \pd{\vb s}v=\Bigl(\pd{s_1}v,\dots,\pd{s_n}v\Bigr)^\T,\qquad
\vb s'(u,v)=\Bigl(\pd{\vb s}u\ \ \pd{\vb s}v\Bigr)\in\R^{n\times2}.
\]
$\vb s'$ er \emph{totalderivert} (et spesialtilfelle av Jacobimatrisen i F4). $\vb s$ er \emph{glatt i $(u,v)$} hvis $\pd{\vb s}u$ og $\pd{\vb s}v$ er lineært uavhengige, altså $\rank\vb s'(u,v)=2$. I $\R^3$ er det det samme som at $\pd{\vb s}u\times\pd{\vb s}v\neq\vb 0$.\\
\textbf{Stykkevis glatt flate:} glatt med begrensede partiellderiverte utenfor endelig mange stykkevis glatte kurver i $D$, og injektiv bortsett fra i endelig mange punkter.
\end{defn}

\begin{teorem}[tangentplan og normalvektor]
I et glatt punkt $(u_0,v_0)$:
\[
\vb t(u,v)=\vb s(u_0,v_0)+\pd{\vb s}u(u_0,v_0)(u-u_0)+\pd{\vb s}v(u_0,v_0)(v-v_0)=\vb s(u_0,v_0)+\vb s'(u_0,v_0)\begin{pmatrix}u-u_0\\v-v_0\end{pmatrix}.
\]
I $\R^3$ står $\vb n=\pd{\vb s}u\times\pd{\vb s}v$ vinkelrett på tangentplanet, og $\vb n/\norm{\vb n}_2$ er enhetsnormalen. Kryssproduktet er
$\vb a\times\vb b=(a_2b_3-a_3b_2,\ a_3b_1-a_1b_3,\ a_1b_2-a_2b_1)^\T$.
\end{teorem}

\begin{teorem}[flate gitt ved en likning $g(x,y,z)=c$ (prøveeksamen, oppg.~3)]
Gradienten $\nabla g(\vb p)$ står vinkelrett på nivåflaten (se F5). Tangentplanet i $\vb p$ er derfor
\[
\ip{\nabla g(\vb p)}{\vb x-\vb p}=0.
\]
\textbf{Graf} $z=F(x,y)$: $\vb h_u=(1,0,F_u)$, $\vb h_v=(0,1,F_v)$ og $\vb n=(-F_u,-F_v,1)$. Tangentplanet er $z=F(u_0,v_0)+F_u(u_0,v_0)(x-u_0)+F_v(u_0,v_0)(y-v_0)$. En graf av en deriverbar $F$ er alltid glatt.
\end{teorem}

\begin{fallgruve}[polene på en kule]
Med kuleparametriseringen er $\pd{\vb s}u=\vb 0$ når $v=0$ eller $v=\pi$, så parametriseringen er ikke glatt der. \emph{Flaten} er likevel glatt: bruk en annen parametrisering eller gradienten til likningen. Skill mellom at parametriseringen ikke er glatt og at flaten ikke er glatt.
\end{fallgruve}

\begin{oppskrift}[typiske oppgaver]
\begin{enumerate}
  \item \emph{Parametriser:} ellipse, kule eller ellipsoide som over. For skjæringskurver, løs den ene likningen for én variabel og sett inn.
  \item \emph{Hastighet, fart og akselerasjon:} deriver komponentvis.
  \item \emph{Glatthet:} finn $t$ der $\vb c'(t)=\vb 0$, eller $(u,v)$ der kolonnene i $\vb s'$ er avhengige.
  \item \emph{Tangentlinje eller -plan:} bruk formlene over. For en flate gitt implisitt, bruk gradienten.
\end{enumerate}
\end{oppskrift}

\begin{oving}
\begin{itemize}[leftmargin=5.6em,labelwidth=5em,labelsep=0.6em,align=left]
\item[\lett] \href{https://github.com/gustavbv/Matte4C---eksamen/blob/main/oppgavesett/F03-kurver-og-flater/1-TMA4432-ovinger-og-proveeksamen/001-TMA4432-2026-proveeksamen-oppg3-oppgave.pdf}{Prøveeksamen, oppg.~3} {\color{graa}(\href{https://github.com/gustavbv/Matte4C---eksamen/blob/main/oppgavesett/F03-kurver-og-flater/1-TMA4432-ovinger-og-proveeksamen/001-TMA4432-2026-proveeksamen-oppg3-svar.pdf}{LF})}: tangentplan til en flate gitt ved en likning
\item[\middels] \href{https://github.com/gustavbv/Matte4C---eksamen/blob/main/oppgavesett/F03-kurver-og-flater/3-hele-oppgaven-i-pensum/019-TMA4105-2023-var-oppg2-oppgave.pdf}{TMA4105 vår 2023, oppg.~2} {\color{graa}(\href{https://github.com/gustavbv/Matte4C---eksamen/blob/main/oppgavesett/F03-kurver-og-flater/3-hele-oppgaven-i-pensum/019-TMA4105-2023-var-oppg2-svar.pdf}{LF})}: tangentplan til en ellipsoide
\item[\middels] \href{https://github.com/gustavbv/Matte4C---eksamen/blob/main/oppgavesett/F03-kurver-og-flater/1-TMA4432-ovinger-og-proveeksamen/004-TMA4432-oving-3-oppg3-oppgave.pdf}{Øving 3, oppg.~3} {\color{graa}(\href{https://github.com/gustavbv/Matte4C---eksamen/blob/main/oppgavesett/F03-kurver-og-flater/1-TMA4432-ovinger-og-proveeksamen/004-TMA4432-oving-3-oppg3-svar.pdf}{LF})}: parametrisering, tangentplan og normalvektor
\item[\vanskelig] \href{https://github.com/gustavbv/Matte4C---eksamen/blob/main/oppgavesett/F03-kurver-og-flater/3-hele-oppgaven-i-pensum/066-TMA4105-2005-var-oppg5-oppgave.pdf}{TMA4105 vår 2005, oppg.~5} {\color{graa}(\href{https://github.com/gustavbv/Matte4C---eksamen/blob/main/oppgavesett/F03-kurver-og-flater/3-hele-oppgaven-i-pensum/066-TMA4105-2005-var-oppg5-svar.pdf}{LF})}: tangentplan og høyeste punkt på en skjæringskurve
\end{itemize}
\end{oving}


\section{F4 · Jacobimatrisen og Newtons metode for systemer}

\textbf{Kjernen:} Den deriverte av $\vb F:\R^n\to\R^m$ er en matrise, Jacobimatrisen, og den gir den beste lineære tilnærmingen. For å løse $\vb F(\vb x)=\vb 0$ setter vi lineariseringen lik null og løser et lineært system, og så gjentar vi. Det er Newtons metode.

\subsection{Vektorfunksjoner og Jacobimatrisen}
\begin{defn}[vektorfunksjon og kontinuitet]
$\vb F:D\to\R^m$, $D\subseteq\R^n$, med komponenter $F_1,\dots,F_m:D\to\R$. $\vb F$ er kontinuerlig i $\vb y$ hvis $\lim_{\vb x\to\vb y}\vb F(\vb x)=\vb F(\vb y)$ (med $\norm\cdot_2$). Dette er det samme som at hver komponent er kontinuerlig.
\end{defn}

\begin{formel}[Jacobimatrise og linearisering\fs]
\[
\vb F'(\vb a)=\begin{pmatrix}\pd{F_1}{x_1}(\vb a)&\cdots&\pd{F_1}{x_n}(\vb a)\\ \vdots&&\vdots\\ \pd{F_m}{x_1}(\vb a)&\cdots&\pd{F_m}{x_n}(\vb a)\end{pmatrix}\in\R^{m\times n},\qquad
\vb L(\vb x)=\vb F(\vb a)+\vb F'(\vb a)(\vb x-\vb a).
\]
Rad $i$ er $\nabla F_i(\vb a)^\T$ (gradienten til komponent $i$), og kolonne $j$ er $\pd{\vb F}{x_j}$. Størrelsen er alltid (antall komponenter) $\times$ (antall variable).
\end{formel}

\begin{defn}[deriverbarhet]
$\vb F$ er deriverbar i $\vb a$ hvis $\vb F(\vb x)=\vb F(\vb a)+\vb F'(\vb a)(\vb x-\vb a)+\vb E(\vb x)$ med $\displaystyle\lim_{\vb x\to\vb a}\frac{\vb E(\vb x)}{\norm{\vb x-\vb a}}=\vb 0$. Lineariseringsfeilen går altså fortere mot null enn $\norm{\vb x-\vb a}$.
\end{defn}
\begin{teorem}
(1) Er $\vb F$ deriverbar i $\vb a$, så er den kontinuerlig i $\vb a$. (2) Finnes alle $\pd{F_i}{x_j}$ i en omegn av $\vb a$, og er de kontinuerlige i $\vb a$, så er $\vb F$ deriverbar i $\vb a$. At de partiellderiverte bare finnes, er ikke nok.
\end{teorem}
Lineariseringen generaliserer $f(a)+f'(a)(x-a)$ ($m=n=1$), tangentlinjen ($n=1$) og tangentplanet ($n=2$).

\subsection{Newtons metode}
\textbf{I én variabel (repetisjon):} $x^{(k+1)}=x^{(k)}-\dfrac{f(x^{(k)})}{f'(x^{(k)})}$. Hvis $f\in C^2$, $f'(x^*)\neq0$ og startpunktet er nær nok, er konvergensen kvadratisk: $\abs{e^{(k+1)}}\le M\abs{e^{(k)}}^2$.

\begin{formel}[Newtons metode for systemer\fs]
For $\vb F:\R^n\to\R^n$ og startpunkt $\vb x^{(0)}$, gjenta for $k=0,1,2,\dots$:
\[
\vb F'(\vb x^{(k)})\,\Delta^{(k)}=-\vb F(\vb x^{(k)}),\qquad \vb x^{(k+1)}=\vb x^{(k)}+\Delta^{(k)}.
\]
\textbf{Utledning:} sett lineariseringen rundt $\hat{\vb x}$ lik null, $\vb F(\hat{\vb x})+\vb F'(\hat{\vb x})(\vb x-\hat{\vb x})=\vb 0$, og løs for oppdateringen $\Delta=\vb x-\hat{\vb x}$.
\end{formel}

\begin{oppskrift}[ett Newton-steg for hånd]
\begin{enumerate}
  \item Skriv systemet på formen $\vb F(\vb x)=\vb 0$ ved å flytte alt over på én side. For eksempel blir $\begin{psmallmatrix}x^2+\cos y\\ \sin x+y^3\end{psmallmatrix}=\begin{psmallmatrix}x\\y\end{psmallmatrix}$ til $\vb F=\begin{psmallmatrix}x^2-x+\cos y\\ \sin x+y^3-y\end{psmallmatrix}$.
  \item Regn ut Jacobimatrisen $\vb F'(\vb x)$ generelt, og sett inn $\vb x^{(k)}$.
  \item Løs $\vb F'\Delta=-\vb F$ med Gauss-eliminasjon. For $2\times2$ kan du bruke $\begin{psmallmatrix}a&b\\c&d\end{psmallmatrix}^{-1}=\frac{1}{ad-bc}\begin{psmallmatrix}d&-b\\-c&a\end{psmallmatrix}$.
  \item $\vb x^{(k+1)}=\vb x^{(k)}+\Delta^{(k)}$. Sett gjerne inn i $\vb F$ for å se at verdien har blitt mindre.
\end{enumerate}
\end{oppskrift}

\begin{eksempel}[forelesningen: $x^3-y+\tfrac14=0$, $x^2+y^2-1=0$, start i $(1,1)$]
$\vb F(1,1)=(\tfrac14,\,1)^\T$ og $\vb F'(x,y)=\begin{psmallmatrix}3x^2&-1\\2x&2y\end{psmallmatrix}$, så $\vb F'(1,1)=\begin{psmallmatrix}3&-1\\2&2\end{psmallmatrix}$ med determinant $8$. Da er
\[
\Delta^{(0)}=\tfrac18\begin{pmatrix}2&1\\-2&3\end{pmatrix}\begin{pmatrix}-\tfrac14\\-1\end{pmatrix}=\begin{pmatrix}-0{,}1875\\-0{,}3125\end{pmatrix},\qquad \vb x^{(1)}=(0{,}8125,\ 0{,}6875)^\T.
\]
Iterasjonene konvergerer mot $(0{,}7463,\ 0{,}6656)$ (kontrollregnet).
\end{eksempel}

\begin{teorem}[konvergens (QSS teorem 7.1)]
La $\vb F\in C^1$ ha et nullpunkt $\vb x^*$, og anta at det finnes $C,R,L>0$ slik at
(1) $\vb F'(\vb x^*)^{-1}$ finnes og $\norm{\vb F'(\vb x^*)^{-1}}_2\le C$, og
(2) $\norm{\vb F'(\vb x)-\vb F'(\vb y)}_2\le L\norm{\vb x-\vb y}_2$ for alle $\vb x,\vb y$ i kulen med radius $R$ om $\vb x^*$ (Lipschitz).
Da er Newton entydig definert for startpunkter i kulen og konvergerer med
\[
\norm{\vb x^{(k+1)}-\vb x^*}_2\le CL\,\norm{\vb x^{(k)}-\vb x^*}_2^2\qquad\text{(kvadratisk konvergens).}
\]
Kort sagt: er $\vb F$ glatt nok og $\vb F'(\vb x^*)$ ikke-singulær, konvergerer metoden kvadratisk for startpunkter nær nok. Konvergensen er bare lokal.
\end{teorem}

\begin{fallgruve}
\begin{itemize}
  \item Newton krever at $\vb F'(\vb x^{(k)})$ er inverterbar (full rang). Kan du starte i et punkt, sjekker du med $\det\vb F'(\vb x^{(0)})\neq0$. I øving 4 er $\vb F=(x_1^2+x_2,\ x_1+x_2^2)$ med $\det\vb F'=4x_1x_2-1$, som er null i $(\tfrac12,\tfrac12)$.
  \item Ikke inverter Jacobimatrisen numerisk. Løs det lineære systemet.
  \item Stoppkriterier: $\norm{\Delta^{(k)}}$ liten, $\norm{\vb F(\vb x^{(k)})}$ liten eller et maksimalt antall iterasjoner.
  \item Hvilket nullpunkt metoden finner, avhenger av startpunktet.
\end{itemize}
\end{fallgruve}

\begin{oving}
\begin{itemize}[leftmargin=5.6em,labelwidth=5em,labelsep=0.6em,align=left]
\item[\lett] \href{https://github.com/gustavbv/Matte4C---eksamen/blob/main/oppgavesett/F04-jacobi-og-newton/3-hele-oppgaven-i-pensum/013-TMA4130-2010-host-oppg8-oppgave.pdf}{TMA4130 høst 2010, oppg.~8} {\color{graa}(\href{https://github.com/gustavbv/Matte4C---eksamen/blob/main/oppgavesett/F04-jacobi-og-newton/3-hele-oppgaven-i-pensum/013-TMA4130-2010-host-oppg8-svar.pdf}{LF})}: ett Newton-steg
\item[\lett] \href{https://github.com/gustavbv/Matte4C---eksamen/blob/main/oppgavesett/F04-jacobi-og-newton/1-TMA4432-ovinger-og-proveeksamen/001-TMA4432-2026-proveeksamen-oppg4-oppgave.pdf}{Prøveeksamen, oppg.~4} {\color{graa}(\href{https://github.com/gustavbv/Matte4C---eksamen/blob/main/oppgavesett/F04-jacobi-og-newton/1-TMA4432-ovinger-og-proveeksamen/001-TMA4432-2026-proveeksamen-oppg4-svar.pdf}{LF})}: ett Newton-steg for et $2\times2$-system
\item[\middels] \href{https://github.com/gustavbv/Matte4C---eksamen/blob/main/oppgavesett/F04-jacobi-og-newton/1-TMA4432-ovinger-og-proveeksamen/002-TMA4432-oving-4-oppg1-oppgave.pdf}{Øving 4, oppg.~1} {\color{graa}(\href{https://github.com/gustavbv/Matte4C---eksamen/blob/main/oppgavesett/F04-jacobi-og-newton/1-TMA4432-ovinger-og-proveeksamen/002-TMA4432-oving-4-oppg1-svar.pdf}{LF})}: Jacobimatrise, linearisering og singulær Jacobimatrise
\item[\middels] \href{https://github.com/gustavbv/Matte4C---eksamen/blob/main/oppgavesett/F04-jacobi-og-newton/3-hele-oppgaven-i-pensum/053-MA2501-2017-var-oppg4-oppgave.pdf}{MA2501 vår 2017, oppg.~4} {\color{graa}(\href{https://github.com/gustavbv/Matte4C---eksamen/blob/main/oppgavesett/F04-jacobi-og-newton/3-hele-oppgaven-i-pensum/053-MA2501-2017-var-oppg4-svar.pdf}{LF})}: sette opp Newton og ta to steg
\item[\vanskelig] \href{https://github.com/gustavbv/Matte4C---eksamen/blob/main/oppgavesett/F04-jacobi-og-newton/3-hele-oppgaven-i-pensum/062-TMA4215-2013-kont-oppg2-oppgave.pdf}{TMA4215 kont.\ 2013, oppg.~2} {\color{graa}(\href{https://github.com/gustavbv/Matte4C---eksamen/blob/main/oppgavesett/F04-jacobi-og-newton/3-hele-oppgaven-i-pensum/062-TMA4215-2013-kont-oppg2-svar.pdf}{LF})}: Newton-steg og bevis for konvergens
\item[\vanskelig] \href{https://github.com/gustavbv/Matte4C---eksamen/blob/main/oppgavesett/F04-jacobi-og-newton/3-hele-oppgaven-i-pensum/050-TMA4215-2018-host-oppg3-oppgave.pdf}{TMA4215 høst 2018, oppg.~3} {\color{graa}(\href{https://github.com/gustavbv/Matte4C---eksamen/blob/main/oppgavesett/F04-jacobi-og-newton/3-hele-oppgaven-i-pensum/050-TMA4215-2018-host-oppg3-svar.pdf}{LF})}: feilestimat for Newton for systemer
\end{itemize}
\end{oving}


\section{F5 · Kjerneregel, gradient og optimering}

\textbf{Kjernen:} Gradienten peker dit funksjonen stiger raskest og står vinkelrett på nivåkurvene. I et minimum er gradienten null. Gradient descent går i retning $-\nabla f$, og Newton løser $\nabla f=\vb 0$ med Hessematrisen. Med bibetingelse må $\nabla f$ være parallell med $\nabla g$ (Lagrange).

\subsection{Kjerneregelen}
\begin{teorem}[kjerneregelen]
$\vb F(\vb x)=\vb G(\vb U(\vb x))$ med $\vb U:\R^m\to\R^n$ deriverbar i $\vb a$ og $\vb G:\R^n\to\R^p$ deriverbar i $\vb U(\vb a)$. Da er
\[
\underbrace{\vb F'(\vb a)}_{p\times m}=\underbrace{\vb G'(\vb U(\vb a))}_{p\times n}\ \underbrace{\vb U'(\vb a)}_{n\times m}.
\]
\textbf{Med skalarer:} for $f(x,y,z)=g(u(x,y,z),v(x,y,z))$ er $\pd fx=\pd gu\pd ux+\pd gv\pd vx$.
\textbf{Langs en kurve:} $h(t)=f(\vb c(t))$ gir $h'(t)=f'(\vb c(t))\,\vb c'(t)=\ip{\nabla f(\vb c(t))}{\vb c'(t)}$.
\end{teorem}
\begin{fallgruve}[øving 5]
Regn ut $\vb U'$ og $\vb G'$ hver for seg, sett inn $\vb U(\vb a)$ i $\vb G'$ (ikke $\vb a$), og gang så sammen. I et bestemt punkt er dette ofte enklere enn å skrive ut $\vb F$ og derivere direkte.
\end{fallgruve}

\subsection{Optimering: begreper}
\begin{defn}[minimerer]
Problemet er $\vb x^*\in\argmin_{\vb x\in\R^n}f(\vb x)$ for en kostnadsfunksjon (skalarfelt) $f:\R^n\to\R$. For å maksimere $h$ minimerer du $f=-h$. $\vb x^*$ er en \emph{global} minimerer hvis $f(\vb x^*)\le f(\vb x)$ for alle $\vb x$, og en \emph{lokal} minimerer hvis det gjelder i en omegn, for eksempel $\{\vb x:\norm{\vb x-\vb x^*}_2\le R\}$.
\end{defn}

\begin{defn}[gradient, retningsderivert og Hessematrise]
\[
\nabla f(\vb y)=\Bigl(\pd f{x_1}(\vb y),\dots,\pd f{x_n}(\vb y)\Bigr)^\T\in\R^n,\qquad
D_{\vb r}f(\vb y)=\lim_{h\to0}\frac{f(\vb y+h\vb r)-f(\vb y)}{h}\quad(\norm{\vb r}_2=1),
\]
\[
\Hess f(\vb x)=\bigl(\nabla f\bigr)'(\vb x)=\Bigl(\frac{\partial^2f}{\partial x_i\partial x_j}(\vb x)\Bigr)_{i,j=1}^n.
\]
Gradienten er en kolonnevektor, mens Jacobimatrisen $f'(\vb y)=\nabla f(\vb y)^\T$ er en radvektor. For $f\in C^2$ er Hessematrisen symmetrisk.
\end{defn}

\begin{teorem}[gradientens egenskaper]
\begin{enumerate}
  \item Er $f$ deriverbar, så er $D_{\vb r}f(\vb y)=\ip{\nabla f(\vb y)}{\vb r}$.
  \item Med Cauchy--Schwarz: $\nabla f$ peker i retningen med raskest stigning, og $-\nabla f$ i retningen med raskest nedgang. Største stigning er $\norm{\nabla f}_2$.
  \item \textbf{Nivåmengder:} $\{\vb x:f(\vb x)=c\}$ (nivåkurve for $n=2$, nivåflate for $n=3$). Ligger kurven $\vb c(t)$ i en nivåmengde, er $\ip{\nabla f(\vb c(t))}{\vb c'(t)}=0$, så $\nabla f$ står vinkelrett på nivåmengden. Bevis: deriver $f(\vb c(t))=\alpha$ med kjerneregelen.
  \item \textbf{Nødvendig betingelse:} er $\vb y$ et lokalt minimum (eller maksimum) i det indre av definisjonsmengden, så er $\nabla f(\vb y)=\vb 0$ (kritisk eller stasjonært punkt). Det omvendte gjelder ikke: et kritisk punkt kan også være et maksimum eller et sadelpunkt.
\end{enumerate}
\end{teorem}

\begin{formel}[Nyttige gradienter (øving 5)]
\begin{center}
\begin{tabular}{@{}lll@{}}
\toprule
$f(\vb x)$ & $\nabla f(\vb x)$ & $\Hess f(\vb x)$\\
\midrule
$\ip{\vb a}{\vb x}$ & $\vb a$ & $0$\\
$\norm{\vb x}_2$ & $\vb x/\norm{\vb x}_2$ \ ($\vb x\neq\vb 0$) & \\
$\norm{A\vb x}_2^2$ & $2A^\T A\vb x$ & $2A^\T A$\\
$\vb x^\T A\vb x$ & $(A+A^\T)\vb x$ \ ($=2A\vb x$ hvis $A=A^\T$) & $A+A^\T$\\
$\tfrac12\norm{A\vb x-\vb b}_2^2$ & $A^\T(A\vb x-\vb b)$ & $A^\T A$\\
$\tfrac12\norm{A\vb x-\vb y}_2^2+\mu\norm{\vb x}_2^2$ & $A^\T(A\vb x-\vb y)+2\mu\vb x$ & $A^\T A+2\mu I$\\
\bottomrule
\end{tabular}
\end{center}
Den siste er \textbf{Tikhonov-regularisering} (ridge-regresjon). Minimumet løser $(A^\T A+2\mu I)\vb x=A^\T\vb y$, som er entydig løsbart for alle $\mu>0$, også når $A$ ikke har full rang. Når $\mu\to0^+$, går løsningen mot $A^\dagger\vb y$.
\end{formel}

\begin{teorem}[klassifisering (nyttig, men ikke gjennomgått i F5)]
Er $\nabla f(\vb x^*)=\vb 0$ og $\Hess f(\vb x^*)$ positivt definitt (alle egenverdier $>0$), er $\vb x^*$ et lokalt minimum. Negativt definitt gir lokalt maksimum, og egenverdier med begge fortegn gir et sadelpunkt. For $2\times2$: $\det>0$ og $f_{xx}>0$ gir minimum.
\end{teorem}

\subsection{Gradient descent}
\begin{oppskrift}[gradient descent (bratteste nedstigning)]
Gitt $f$, $\nabla f$, startpunkt $\vb x^{(0)}$ og steglengder $s_k$:
\[
\vb x^{(k+1)}=\vb x^{(k)}-s_k\nabla f(\vb x^{(k)}),\qquad k=0,1,2,\dots
\]
\begin{itemize}
  \item \textbf{Steglengde:} konstant $s$ (enkelt). Eksakt linjesøk $\min_t f(\vb x-t\nabla f(\vb x))$ er et nytt optimeringsproblem og som regel for dyrt. I praksis brukes \textbf{backtracking}: prøv $\beta^\ell u$, $\ell=0,1,2,\dots$ (med $0<\beta<1$), og ta det første steget som gir nok nedgang (Armijo-betingelsen, ofte skrevet $f(\vb x-s\nabla f)\le f(\vb x)-c\,s\norm{\nabla f}_2^2$ med for eksempel $c=10^{-4}$).
  \item \textbf{Stopp} når $\norm{\nabla f(\vb x^{(k)})}_2<\varepsilon$ eller etter et maksimalt antall steg.
  \item For stor $s$: iteratene svinger eller divergerer, og $f$ kan øke. For liten $s$: konvergensen blir treg.
  \item Starter du i et kritisk punkt, er $\nabla f=\vb 0$ og metoden står stille (prøveeksamen, oppg.~5c).
\end{itemize}
\end{oppskrift}

\begin{oppskrift}[Newton for optimering]
Bruk Newton (F4) på $\vb F=\nabla f$, slik at Jacobimatrisen blir Hessematrisen:
\[
\Hess f(\vb x^{(k)})\,\Delta^{(k)}=-\nabla f(\vb x^{(k)}),\qquad \vb x^{(k+1)}=\vb x^{(k)}+\Delta^{(k)}.
\]
Metoden trenger andrederiverte og at $\Hess f(\vb x^{(k)})$ er inverterbar. Den konvergerer kvadratisk mot et kritisk punkt, som også kan være et maksimum eller et sadelpunkt. For en kvadratisk $f$ er ett steg nok.
\end{oppskrift}

\begin{oving}
\begin{itemize}[leftmargin=5.6em,labelwidth=5em,labelsep=0.6em,align=left]
\item[\lett] \href{https://github.com/gustavbv/Matte4C---eksamen/blob/main/oppgavesett/F05-gradient-og-optimering/3-hele-oppgaven-i-pensum/007-TMA4135-2016-host-oppg5-oppgave.pdf}{TMA4135 høst 2016, oppg.~5} {\color{graa}(\href{https://github.com/gustavbv/Matte4C---eksamen/blob/main/oppgavesett/F05-gradient-og-optimering/3-hele-oppgaven-i-pensum/007-TMA4135-2016-host-oppg5-svar.pdf}{LF})}: gradient, retningsderivert og nivåkurve
\item[\middels] \href{https://github.com/gustavbv/Matte4C---eksamen/blob/main/oppgavesett/F05-gradient-og-optimering/1-TMA4432-ovinger-og-proveeksamen/002-TMA4432-oving-5-oppg1-oppgave.pdf}{Øving 5, oppg.~1} {\color{graa}(\href{https://github.com/gustavbv/Matte4C---eksamen/blob/main/oppgavesett/F05-gradient-og-optimering/1-TMA4432-ovinger-og-proveeksamen/002-TMA4432-oving-5-oppg1-svar.pdf}{LF})}: kjerneregelen med Jacobimatriser
\item[\middels] \href{https://github.com/gustavbv/Matte4C---eksamen/blob/main/oppgavesett/F05-gradient-og-optimering/1-TMA4432-ovinger-og-proveeksamen/003-TMA4432-oving-5-oppg2-oppgave.pdf}{Øving 5, oppg.~2} {\color{graa}(\href{https://github.com/gustavbv/Matte4C---eksamen/blob/main/oppgavesett/F05-gradient-og-optimering/1-TMA4432-ovinger-og-proveeksamen/003-TMA4432-oving-5-oppg2-svar.pdf}{LF})}: gradienter av normer og kvadratiske former, Tikhonov
\end{itemize}
\end{oving}

\begin{oving}
\begin{itemize}[leftmargin=5.6em,labelwidth=5em,labelsep=0.6em,align=left]
\item[\middels] \href{https://github.com/gustavbv/Matte4C---eksamen/blob/main/oppgavesett/F05-gradient-og-optimering/1-TMA4432-ovinger-og-proveeksamen/001-TMA4432-2026-proveeksamen-oppg5-oppgave.pdf}{Prøveeksamen, oppg.~5} {\color{graa}(\href{https://github.com/gustavbv/Matte4C---eksamen/blob/main/oppgavesett/F05-gradient-og-optimering/1-TMA4432-ovinger-og-proveeksamen/001-TMA4432-2026-proveeksamen-oppg5-svar.pdf}{LF})}: gradient descent for hånd, steglengde og Newton
\item[\middels] \href{https://github.com/gustavbv/Matte4C---eksamen/blob/main/oppgavesett/F05-gradient-og-optimering/1-TMA4432-ovinger-og-proveeksamen/004-TMA4432-oving-5-oppg3-oppgave.pdf}{Øving 5, oppg.~3} {\color{graa}(\href{https://github.com/gustavbv/Matte4C---eksamen/blob/main/oppgavesett/F05-gradient-og-optimering/1-TMA4432-ovinger-og-proveeksamen/004-TMA4432-oving-5-oppg3-svar.pdf}{LF})}: gradient descent og Lagrange
\item[\vanskelig] \href{https://github.com/gustavbv/Matte4C---eksamen/blob/main/oppgavesett/F05-gradient-og-optimering/2-Bergmann-eksamener/005-TMA4215-2022-host-oppg2-oppgave.pdf}{TMA4215 høst 2022, oppg.~2} {\color{graa}(\href{https://github.com/gustavbv/Matte4C---eksamen/blob/main/oppgavesett/F05-gradient-og-optimering/2-Bergmann-eksamener/005-TMA4215-2022-host-oppg2-svar.pdf}{LF})}: gradient descent sammenlignet med normallikningen
\item[\vanskelig] \href{https://github.com/gustavbv/Matte4C---eksamen/blob/main/oppgavesett/F05-gradient-og-optimering/4-delvis-i-pensum/154-TMA4215-2019-kont-oppg3-oppgave.pdf}{TMA4215 kont.\ 2019, oppg.~3} {\color{graa}(\href{https://github.com/gustavbv/Matte4C---eksamen/blob/main/oppgavesett/F05-gradient-og-optimering/4-delvis-i-pensum/154-TMA4215-2019-kont-oppg3-svar.pdf}{LF})}: kritiske punkter og Newton for optimering (delvis pensum)
\end{itemize}
\end{oving}


\subsection{Optimering med bibetingelse: Lagrange-multiplikatorer}
\begin{defn}
Problemet er $\min f(\vb x)$ slik at $g(\vb x)=0$. Punkter med $g(\vb x)=0$ er \emph{tillatte} (feasible). Ved et optimum kan vi ikke kreve $\nabla f=\vb 0$, men ingen tillatt retning (tangent til $\{g=0\}$) kan gi forbedring. Derfor må $\nabla f$ stå vinkelrett på nivåmengden, altså være parallell med $\nabla g$.
\end{defn}
\begin{teorem}[Lagrange]
La $f,g:\R^n\to\R$ ha kontinuerlige gradienter. Hvis $\vb x^*$ er et ekstremalpunkt for $f$ under bibetingelsen $g=0$ og $\nabla g(\vb x^*)\neq\vb 0$, finnes en $\lambda\in\R$ (Lagrange-multiplikator) slik at
\[
\nabla f(\vb x^*)=\lambda\nabla g(\vb x^*).
\]
Med \textbf{Lagrangefunksjonen} $\mathcal L(\vb x,\lambda)=f(\vb x)+\lambda g(\vb x)$ betyr $\nabla\mathcal L=\vb 0$ at $\nabla f+\lambda\nabla g=\vb 0$ og $g(\vb x)=0$. Fortegnet på $\lambda$ er bare en konvensjon. Er $\vb x^*$ også et ubetinget minimum, er $\lambda=0$. Kritiske punkter for $\mathcal L$ er som regel sadelpunkter for $\mathcal L$.
\end{teorem}
\begin{oppskrift}[Lagrange for hånd]
\begin{enumerate}
  \item Sett opp $\nabla f=\lambda\nabla g$ og $g=0$: $n+1$ likninger med $n+1$ ukjente.
  \item Løs, ofte ved å eliminere $\lambda$, for eksempel ved å dele to likninger på hverandre eller dele opp i tilfeller.
  \item Ta også med tillatte punkter der $\nabla g=\vb 0$. Der gjelder ikke teoremet.
  \item Regn ut $f$ i alle kandidatene. Er mengden $\{g=0\}$ lukket og begrenset (for eksempel en sirkel), er den største verdien maksimum og den minste minimum.
\end{enumerate}
\end{oppskrift}
\begin{eksempel}[øving 5, oppg.~3: $f=x^2+4y^2-4x$ på $x^2+y^2=1$]
$\nabla f=(2x-4,\,8y)$ og $\nabla g=(2x,\,2y)$. Likningene $2x-4=2\lambda x$, $8y=2\lambda y$ og $x^2+y^2=1$ gir enten $y=0$, $x=\pm1$ eller $\lambda=4$, $x=-\tfrac23$, $y=\pm\tfrac{\sqrt5}3$. Verdiene er $f(1,0)=-3$ (minimum), $f(-1,0)=5$ og $f(-\tfrac23,\pm\tfrac{\sqrt5}{3})=\tfrac{16}3$ (maksimum). Kontrollregnet.
\end{eksempel}
\textbf{Videre (ikke eksamen):} flere bibetingelser $\vb H(\vb x)=\vb 0$, ulikheter $\vb G(\vb x)\le\vb 0$ (KKT), straffemetoder, utvidet Lagrange og indrepunktsmetoder.

\begin{oving}
\begin{itemize}[leftmargin=5.6em,labelwidth=5em,labelsep=0.6em,align=left]
\item[\lett] \href{https://github.com/gustavbv/Matte4C---eksamen/blob/main/oppgavesett/F04-jacobi-og-newton/3-hele-oppgaven-i-pensum/026-TMA4105-2024-var-oppg1-oppgave.pdf}{TMA4105 vår 2024, oppg.~1} {\color{graa}(\href{https://github.com/gustavbv/Matte4C---eksamen/blob/main/oppgavesett/F04-jacobi-og-newton/3-hele-oppgaven-i-pensum/026-TMA4105-2024-var-oppg1-svar.pdf}{LF})}: minimum med én lineær bibetingelse
\item[\middels] \href{https://github.com/gustavbv/Matte4C---eksamen/blob/main/oppgavesett/F05-gradient-og-optimering/3-hele-oppgaven-i-pensum/056-TMA4105-2017-var-oppg3-oppgave.pdf}{TMA4105 vår 2017, oppg.~3} {\color{graa}(\href{https://github.com/gustavbv/Matte4C---eksamen/blob/main/oppgavesett/F05-gradient-og-optimering/3-hele-oppgaven-i-pensum/056-TMA4105-2017-var-oppg3-svar.pdf}{LF})}: største og minste verdi på en sirkel
\item[\middels] \href{https://github.com/gustavbv/Matte4C---eksamen/blob/main/oppgavesett/F05-gradient-og-optimering/3-hele-oppgaven-i-pensum/028-TMA4105-2025-var-oppg7-oppgave.pdf}{TMA4105 vår 2025, oppg.~7} {\color{graa}(\href{https://github.com/gustavbv/Matte4C---eksamen/blob/main/oppgavesett/F05-gradient-og-optimering/3-hele-oppgaven-i-pensum/028-TMA4105-2025-var-oppg7-svar.pdf}{LF})}: minste avstand fra origo til en kurve
\item[\vanskelig] \href{https://github.com/gustavbv/Matte4C---eksamen/blob/main/oppgavesett/F03-kurver-og-flater/3-hele-oppgaven-i-pensum/066-TMA4105-2005-var-oppg5-oppgave.pdf}{TMA4105 vår 2005, oppg.~5} {\color{graa}(\href{https://github.com/gustavbv/Matte4C---eksamen/blob/main/oppgavesett/F03-kurver-og-flater/3-hele-oppgaven-i-pensum/066-TMA4105-2005-var-oppg5-svar.pdf}{LF})}: høyeste punkt på en skjæringskurve
\end{itemize}
\end{oving}


\section{F6 · Fourierrekker}

\textbf{Kjernen:} En periodisk funksjon er en sum av sinus- og cosinusbølger (eller $\ee^{\ii kx}$). Koeffisientene er projeksjoner på et ortogonalt system, som koordinater i en ortogonal basis. Delsummene er den beste tilnærmingen i $L^2$-normen.

\subsection{Periodiske funksjoner og trigonometriske polynomer}
\begin{defn}[periode]
$f:\R\to\C$ er $T$-periodisk hvis $f(x+T)=f(x)$ for alle $x$. Den minste slike $T>0$ er \emph{fundamentalperioden}. Er $g$ $\tfrac Tn$-periodisk, kalles $n$ frekvensen (på $[0,T]$). For eksempel har $\sin(nx)$ fundamentalperiode $2\pi/n$, og en lineærkombinasjon av $T$-periodiske funksjoner er $T$-periodisk. Kurset bruker stort sett $T=2\pi$.
\end{defn}
\begin{teorem}[Euler]
$\ee^{\ii x}=\cos x+\ii\sin x$, \quad $\cos x=\dfrac{\ee^{\ii x}+\ee^{-\ii x}}{2}$, \quad $\sin x=\dfrac{\ee^{\ii x}-\ee^{-\ii x}}{2\ii}$, \quad $\ee^{\pm\ii n\pi}=\cos(n\pi)=(-1)^n$, \quad $\sin(n\pi)=0$.
\end{teorem}
\begin{defn}[trigonometriske polynomer av grad $\le N$]
$V_N=\bigl\{a_0+\sum_{n=1}^N a_n\cos(nx)+b_n\sin(nx)\bigr\}=\bigl\{\sum_{k=-N}^{N}c_k\ee^{\ii kx}\bigr\}$. Den reelle og den komplekse formen beskriver samme rom.
\end{defn}

\subsection{Indreprodukt, ortogonale systemer og beste tilnærming}
\begin{defn}[indreprodukt og norm for $2\pi$-periodiske funksjoner]
\[
\ip fg=\int_{-\pi}^{\pi}f(x)\overline{g(x)}\dd x,\qquad \norm f=\sqrt{\ip ff},\qquad \abs{\ip fg}\le\norm f\norm g\ \text{(Cauchy--Schwarz)}.
\]
Rommet er $L^2$: funksjoner med $\int_{-\pi}^\pi\abs f^2<\infty$. Indreproduktet er lineært i første argument og konjugert symmetrisk.
\end{defn}
\begin{defn}[ortogonalt og ortonormalt system]
$\{\varphi_k\}$ er ortogonalt hvis $\ip{\varphi_k}{\varphi_n}=0$ for $k\neq n$ og $=d_n\neq0$ for $k=n$, og ortonormalt hvis i tillegg $d_n=1$. Del på $\sqrt{d_k}$ for å normere.
\end{defn}
\begin{teorem}[det trigonometriske systemet er ortogonalt]
\begin{gather*}
\int_{-\pi}^{\pi}\ee^{\ii kx}\overline{\ee^{\ii nx}}\dd x=2\pi\,\delta_{kn},\\
\int_{-\pi}^{\pi}\cos(kx)\cos(nx)\dd x=\int_{-\pi}^{\pi}\sin(kx)\sin(nx)\dd x=\pi\,\delta_{kn}\quad(k,n\ge1),
\end{gather*}
\[
\int_{-\pi}^{\pi}\cos(kx)\sin(nx)\dd x=0,\qquad \int_{-\pi}^{\pi}1\dd x=2\pi.
\]
\end{teorem}
\begin{teorem}[generelle Fourierkoeffisienter og projeksjon]
For et ortogonalt system er $\hat f(k)=c_k=\dfrac{\ip f{\varphi_k}}{\norm{\varphi_k}^2}$, og $\Pi_Nf=\sum_{k}c_k\varphi_k$ (summert over de første $N$) er den ortogonale projeksjonen på $V_N$:
\begin{itemize}
  \item $\ip{f-\Pi_Nf}g=0$ for alle $g\in V_N$, og $\norm{f-\Pi_Nf}=\min_{g\in V_N}\norm{f-g}$ (beste tilnærming).
  \item $\norm{\Pi_Nf}\le\norm f$. For et ortonormalt system gjelder \textbf{Bessels ulikhet} $\sum_k\abs{\hat f(k)}^2\le\norm f^2$ og \textbf{Riemann--Lebesgue}: $\hat f(k)\to0$.
  \item Delsummen $S_Nf=\Pi_Nf$. Fourierrekken er den formelle rekken $\sum_k\hat f(k)\varphi_k$.
\end{itemize}
\end{teorem}

\subsection{Reelle og komplekse Fourierrekker}
\begin{formel}[Fourierrekke for $2\pi$-periodisk $f$\fs]
\[
f\sim\sum_{k=-\infty}^{\infty}c_k\ee^{\ii kx}=a_0+\sum_{n=1}^\infty a_n\cos(nx)+b_n\sin(nx),
\]
\begin{gather*}
c_k=\frac1{2\pi}\int_{-\pi}^{\pi}f(x)\ee^{-\ii kx}\dd x,\qquad a_0=\frac1{2\pi}\int_{-\pi}^{\pi}f(x)\dd x,\\
a_n=\frac1\pi\int_{-\pi}^{\pi}f(x)\cos(nx)\dd x,\qquad b_n=\frac1\pi\int_{-\pi}^{\pi}f(x)\sin(nx)\dd x.
\end{gather*}
$a_0$ er middelverdien av $f$. Hvis $f$ er reell, er alle $a_n$ og $b_n$ reelle. Integrasjonsintervallet kan være en hvilken som helst periode, for eksempel $[0,2\pi]$.
\end{formel}

\begin{teorem}[omregning mellom reell og kompleks form]
\begin{gather*}
a_0=c_0,\qquad a_n=c_n+c_{-n},\qquad b_n=\ii(c_n-c_{-n}),\\ c_n=\frac{a_n-\ii b_n}{2},\qquad c_{-n}=\frac{a_n+\ii b_n}{2}\qquad(n\ge1).
\end{gather*}
Delsummene er like: $S_Nf=\sum_{k=-N}^N c_k\ee^{\ii kx}=a_0+\sum_{n=1}^N(a_n\cos nx+b_n\sin nx)$. For reell $f$ er $c_{-k}=\overline{c_k}$.
\end{teorem}

\begin{oppskrift}[finn koeffisientene]
\begin{enumerate}
  \item \textbf{Er $f$ allerede et trigonometrisk polynom?} Les av koeffisientene uten å integrere. Bruk Euler og identitetene for å skrive om, for eksempel $\sin^2x=\tfrac{1-\cos2x}2$, $\cos^2x=\tfrac{1+\cos2x}2$ og $2\cos(4x)=\ee^{4\ii x}+\ee^{-4\ii x}$. For eksempel gir $f=2\cos4x+3\ee^{6\ii x}$ at $c_{\pm4}=1$, $c_6=3$ og ellers $0$.
  \item \textbf{Bruk symmetri.} Er $f$ like, er $b_n=0$ og $a_n=\frac2\pi\int_0^\pi f\cos(nx)\dd x$ (cosinusrekke). Er $f$ odde, er $a_0=a_n=0$ og $b_n=\frac2\pi\int_0^\pi f\sin(nx)\dd x$ (sinusrekke).
  \item \textbf{Integrer}, ofte med delvis integrasjon: $\int x\cos(nx)\dd x=\frac{x\sin(nx)}n+\frac{\cos(nx)}{n^2}$ og $\int x\sin(nx)\dd x=-\frac{x\cos(nx)}n+\frac{\sin(nx)}{n^2}$. Del opp stykkevise funksjoner.
  \item \textbf{Forenkle} med $\cos(n\pi)=(-1)^n$, og skill gjerne mellom like og odde $n$.
\end{enumerate}
\end{oppskrift}

\begin{defn}[like og odde funksjoner]
$f$ er like hvis $f(-x)=f(x)$ og odde hvis $f(-x)=-f(x)$. Like$\cdot$like og odde$\cdot$odde er like, og like$\cdot$odde er odde. Er $f$ odde, er $\int_{-L}^Lf=0$. Er $f$ like, er $\int_{-L}^Lf=2\int_0^Lf$.
\end{defn}

\begin{eksempel}[{standardrekker på $[-\pi,\pi)$, kontrollregnet}]
\begin{itemize}
  \item $f(x)=x$: $a_n=0$, $b_n=\dfrac{2(-1)^{n+1}}{n}$, $c_0=0$, $c_k=\dfrac{\ii(-1)^k}{k}$.
  \item $f(x)=\abs x$: $a_0=\dfrac\pi2$, $a_n=\dfrac{2((-1)^n-1)}{\pi n^2}$ ($=-\frac{4}{\pi n^2}$ for odde $n$, $0$ for like $n$), $b_n=0$.
  \item $f(x)=x^2$: $a_0=\dfrac{\pi^2}{3}$, $a_n=\dfrac{4(-1)^n}{n^2}$, $b_n=0$.
  \item Heaviside $u(x)=1$ for $x\ge0$ og $0$ for $x<0$: $u\sim\frac12+\frac2\pi\sum_{n\ge1}\frac{\sin((2n-1)x)}{2n-1}$.
\end{itemize}
\end{eksempel}

\subsection{Konvergens}
\begin{teorem}[punktvis konvergens]
La $f$ være $2\pi$-periodisk og stykkevis $C^1$, med ensidige grenser $f(x^\pm)$ og ensidige deriverte i hoppene. Da er
\[
\lim_{N\to\infty}S_Nf(x)=\tfrac12\bigl(f(x^-)+f(x^+)\bigr).
\]
Der $f$ er kontinuerlig, konvergerer rekken altså mot $f(x)$, og i et hopp mot midtpunktet. Husk å sjekke endepunktene $\pm\pi$ til den \emph{periodiske} utvidelsen.
\end{teorem}
\begin{itemize}
  \item \textbf{Gibbs-fenomenet:} nær et hopp skyter $S_Nf$ over med omtrent 9\,\% av hoppet, og oversvingen forsvinner ikke når $N\to\infty$. Den blir bare smalere.
  \item \textbf{Hvor fort koeffisientene avtar:} har $f$ et hopp, avtar de som $1/k$. Er $f$ kontinuerlig med knekk, avtar de som $1/k^2$. Glattere $f$ gir raskere avtakende koeffisienter.
\end{itemize}

\begin{defn}[Dirichlet-kjernen]
\begin{gather*}
D_N(x)=\sum_{k=-N}^{N}\ee^{\ii kx}=1+2\sum_{n=1}^N\cos(nx)=\frac{\sin\bigl((N+\frac12)x\bigr)}{\sin(x/2)},\\
\int_{-\pi}^\pi D_N(x)\dd x=2\pi,\qquad D_N(0)=2N+1.
\end{gather*}
Den normerte kjernen $\tilde D_N=D_N/(2\pi)$ har integral 1, og $S_Nf=\tilde D_N*f$. Den lukkede formen kommer fra den geometriske rekken.
\end{defn}

\subsection{Konvolusjon}
\begin{defn}[periodisk konvolusjon]
$(f*g)(x)=\displaystyle\int_{-\pi}^{\pi}f(y)g(x-y)\dd y$, og $f*g=g*f$.
\textbf{Tolkning:} med $h_\varepsilon=\frac1{2\varepsilon}$ på $[-\varepsilon,\varepsilon]$ (og $0$ ellers) er $(f*h_\varepsilon)(0)=\frac1{2\varepsilon}\int_{-\varepsilon}^{\varepsilon}f$, en middelverdi som går mot $f(0)$ når $\varepsilon\to0$ hvis $f$ er kontinuerlig.
\end{defn}
\begin{teorem}[konvolusjon gir produkt av koeffisientene]
$c_k(f*g)=2\pi\,c_k(f)\,c_k(g)$. (Noen bøker har en faktor $\frac1{2\pi}$ i definisjonen av $*$, og da forsvinner $2\pi$.)
\end{teorem}
\begin{oppskrift}[regn ut $h=f*g$ (prøveeksamen 6c, øving 7)]
Finn $c_k(f)$ og $c_k(g)$, gang dem sammen med faktoren $2\pi$, og sett sammen $h=\sum_k c_k(h)\ee^{\ii kx}$. Bare frekvenser som finnes i \emph{både} $f$ og $g$, overlever. Gjør om til reell form til slutt.
For eksempel gir $f=\sin3x+5\cos2x$ og $g=x^2$ på $[-\pi,\pi]$ (med $c_k(g)=\frac{2(-1)^k}{k^2}$ for $k\neq0$) at $h=-\frac{4\pi}9\sin3x+5\pi\cos2x$.
\end{oppskrift}

\subsection{Periode $2L$, utvidelser og Parseval}
\begin{formel}[Fourierrekke med periode $2L$ (ikke på formelarket)]
Bytt $x$ med $\frac{\pi x}{L}$:
\[
f\sim\sum_{k\in\Z}c_k\ee^{\ii k\pi x/L}=a_0+\sum_{n=1}^\infty a_n\cos\frac{n\pi x}L+b_n\sin\frac{n\pi x}L,
\]
\begin{gather*}
c_k=\frac1{2L}\int_{-L}^{L}f(x)\ee^{-\ii k\pi x/L}\dd x,\qquad a_0=\frac1{2L}\int_{-L}^Lf(x)\dd x,\\
a_n=\frac1L\int_{-L}^Lf(x)\cos\frac{n\pi x}L\dd x,\qquad b_n=\frac1L\int_{-L}^Lf(x)\sin\frac{n\pi x}L\dd x.
\end{gather*}
\end{formel}
\begin{defn}[like og odde utvidelse (halvperiodeutvidelse)]
For $f$ på $[0,L]$: den \emph{odde} utvidelsen er $f_o(x)=-f(-x)$ for $x<0$, og den \emph{like} er $f_e(x)=f(-x)$ for $x<0$. Utvid deretter $2L$-periodisk.
\begin{align*}
&\text{odde (sinusrekke):} && b_n=\frac2L\int_0^Lf(x)\sin\frac{n\pi x}L\dd x,\\
&\text{like (cosinusrekke):} && a_0=\frac1L\int_0^Lf(x)\dd x,\qquad a_n=\frac2L\int_0^Lf(x)\cos\frac{n\pi x}L\dd x.
\end{align*}
Dette brukes direkte når initialbetingelser i PDE-er skal utvikles (F9 og F10).
\end{defn}
\begin{teorem}[Parsevals identitet]
For $f\in L^2(-L,L)$:
\[
\frac1{2L}\int_{-L}^{L}\abs{f(x)}^2\dd x=\sum_{k=-\infty}^{\infty}\abs{c_k}^2=a_0^2+\frac12\sum_{n=1}^{\infty}\bigl(a_n^2+b_n^2\bigr).
\]
\textbf{Rekkesummer:} sett inn kjente koeffisienter. For $f(x)=x$ på $[-\pi,\pi)$ gir Parseval $\frac{\pi^2}3=\frac12\sum\frac{4}{n^2}$, altså $\sum\frac1{n^2}=\frac{\pi^2}6$. Punktvis konvergens for Heaviside i $x=\frac\pi2$ gir $1-\frac13+\frac15-\dots=\frac\pi4$.
\end{teorem}

\subsection{Spekter og regneregler}
\begin{defn}[spekter, amplitude og fase]
For en $2L$-periodisk $f$ er spekteret paret $\bigl(\hat f(k),\frac k{2L}\bigr)_{k\in\Z}$. Skriv $c_k=\abs{c_k}\ee^{\ii\theta_k}$: $\abs{c_k}$ er amplituden, $\theta_k$ fasen og $\frac k{2L}$ frekvensen.
\end{defn}
\begin{teorem}[regneregler ($2\pi$-periodisk)]
\begin{itemize}
  \item \textbf{Reell $f$:} $c_{-k}(f)=\overline{c_k(f)}$.
  \item \textbf{Linearitet:} $c_k(\alpha f+\beta g)=\alpha c_k(f)+\beta c_k(g)$, og tilsvarende for $a_n$ og $b_n$.
  \item \textbf{Forskyvning:} $h(x)=f(x-x_0)$ gir $c_k(h)=\ee^{-\ii kx_0}c_k(f)$. Lysarket har feil fortegn (se rettelsene i avsnitt~\ref{sec:rettelser}).
  \item \textbf{Modulasjon:} $h(x)=\ee^{-\ii k_0x}f(x)$ gir $c_k(h)=c_{k+k_0}(f)$.
  \item \textbf{Derivasjon:} $c_k(f')=\ii k\,c_k(f)$. Reelt er $a_0(f')=0$, $a_n(f')=n\,b_n(f)$ og $b_n(f')=-n\,a_n(f)$, forutsatt at $f$ er periodisk og kontinuerlig. Har den periodiske utvidelsen et hopp, gjelder ikke formelen.
\end{itemize}
\end{teorem}

\begin{oving}
\begin{itemize}[leftmargin=5.6em,labelwidth=5em,labelsep=0.6em,align=left]
\item[\lett] \href{https://github.com/gustavbv/Matte4C---eksamen/blob/main/oppgavesett/F06-fourierrekker/1-TMA4432-ovinger-og-proveeksamen/003-TMA4432-oving-6-oppg2-oppgave.pdf}{Øving 6, oppg.~2} {\color{graa}(\href{https://github.com/gustavbv/Matte4C---eksamen/blob/main/oppgavesett/F06-fourierrekker/1-TMA4432-ovinger-og-proveeksamen/003-TMA4432-oving-6-oppg2-svar.pdf}{LF})}: koeffisienter ved avlesning og omregning
\item[\lett] \href{https://github.com/gustavbv/Matte4C---eksamen/blob/main/oppgavesett/F06-fourierrekker/2-Bergmann-eksamener/013-TMA4125-2024-var-oppg3-oppgave.pdf}{Matte 4N vår 2024, oppg.~3} {\color{graa}(\href{https://github.com/gustavbv/Matte4C---eksamen/blob/main/oppgavesett/F06-fourierrekker/2-Bergmann-eksamener/013-TMA4125-2024-var-oppg3-svar.pdf}{LF})}: avlesning av komplekse koeffisienter
\item[\middels] \href{https://github.com/gustavbv/Matte4C---eksamen/blob/main/oppgavesett/F06-fourierrekker/1-TMA4432-ovinger-og-proveeksamen/007-TMA4432-oving-7-oppg1-oppgave.pdf}{Øving 7, oppg.~1} {\color{graa}(\href{https://github.com/gustavbv/Matte4C---eksamen/blob/main/oppgavesett/F06-fourierrekker/1-TMA4432-ovinger-og-proveeksamen/007-TMA4432-oving-7-oppg1-svar.pdf}{LF})}: like og odde utvidelse
\item[\middels] \href{https://github.com/gustavbv/Matte4C---eksamen/blob/main/oppgavesett/F06-fourierrekker/2-Bergmann-eksamener/014-TMA4125-2023-var-oppg3-oppgave.pdf}{Matte 4N vår 2023, oppg.~3} {\color{graa}(\href{https://github.com/gustavbv/Matte4C---eksamen/blob/main/oppgavesett/F06-fourierrekker/2-Bergmann-eksamener/014-TMA4125-2023-var-oppg3-svar.pdf}{LF})}: like utvidelse, reell og kompleks rekke
\item[\middels] \href{https://github.com/gustavbv/Matte4C---eksamen/blob/main/oppgavesett/F06-fourierrekker/3-hele-oppgaven-i-pensum/043-TMA4125-2018-var-oppg1-oppgave.pdf}{TMA4125 vår 2018, oppg.~1} {\color{graa}(\href{https://github.com/gustavbv/Matte4C---eksamen/blob/main/oppgavesett/F06-fourierrekker/3-hele-oppgaven-i-pensum/043-TMA4125-2018-var-oppg1-svar.pdf}{LF})}: sinusrekke og Parseval gir $\sum 1/n^2$
\item[\middels] \href{https://github.com/gustavbv/Matte4C---eksamen/blob/main/oppgavesett/F06-fourierrekker/1-TMA4432-ovinger-og-proveeksamen/008-TMA4432-oving-7-oppg2-oppgave.pdf}{Øving 7, oppg.~2} {\color{graa}(\href{https://github.com/gustavbv/Matte4C---eksamen/blob/main/oppgavesett/F06-fourierrekker/1-TMA4432-ovinger-og-proveeksamen/008-TMA4432-oving-7-oppg2-svar.pdf}{LF})}: konvolusjon via koeffisientene
\item[\vanskelig] \href{https://github.com/gustavbv/Matte4C---eksamen/blob/main/oppgavesett/F06-fourierrekker/2-Bergmann-eksamener/018-TMA4125-30-35-2022-kont-oppg6-oppgave.pdf}{Matte 4N kont.\ 2022, oppg.~6} {\color{graa}(\href{https://github.com/gustavbv/Matte4C---eksamen/blob/main/oppgavesett/F06-fourierrekker/2-Bergmann-eksamener/018-TMA4125-30-35-2022-kont-oppg6-svar.pdf}{LF})}: odde og like utvidelse, konvergens i hopp
\item[\vanskelig] \href{https://github.com/gustavbv/Matte4C---eksamen/blob/main/oppgavesett/F06-fourierrekker/2-Bergmann-eksamener/012-TMA4125-30-2025-kont-oppg4-oppgave.pdf}{Matte 4N kont.\ 2025, oppg.~4} {\color{graa}(\href{https://github.com/gustavbv/Matte4C---eksamen/blob/main/oppgavesett/F06-fourierrekker/2-Bergmann-eksamener/012-TMA4125-30-2025-kont-oppg4-svar.pdf}{LF})}: Dirichlet-kjernen og konvolusjon
\item[\vanskelig] \href{https://github.com/gustavbv/Matte4C---eksamen/blob/main/oppgavesett/F06-fourierrekker/3-hele-oppgaven-i-pensum/019-TMA4125-2026-var-oppg1-oppgave.pdf}{TMA4125 vår 2026, oppg.~1} {\color{graa}(\href{https://github.com/gustavbv/Matte4C---eksamen/blob/main/oppgavesett/F06-fourierrekker/3-hele-oppgaven-i-pensum/019-TMA4125-2026-var-oppg1-svar.pdf}{LF})}: utlede Dirichlet-kjernen
\end{itemize}
\end{oving}


\section{F7 · Den diskrete Fouriertransformen (DFT) og FFT}

\textbf{Kjernen:} DFT-en er Fourierkoeffisientene regnet ut fra $N$ punktverdier med trapesregelen, altså en matrise-vektor-multiplikasjon. Frekvenser som skiller seg med $N$, kan ikke skilles på gitteret (aliasing). FFT regner ut det samme på $O(N\log N)$ operasjoner.

\subsection{Fra Fourierkoeffisienter til DFT}
Trapesregelen på $[0,2\pi)$ med $N$ punkter $x_j=\frac{2\pi j}N$ gir
\[
c_k(f)=\frac1{2\pi}\int_0^{2\pi}f(x)\ee^{-\ii kx}\dd x\approx\frac1N\sum_{j=0}^{N-1}f\Bigl(\frac{2\pi j}{N}\Bigr)\ee^{-2\pi\ii jk/N}.
\]
Vi ser derfor på et signal $\vb f=(f_0,\dots,f_{N-1})^\T\in\C^N$ med $f_j=f(2\pi j/N)$.

\begin{formel}[DFT og IDFT\fs]
\[
\hat f_k=\frac1N\sum_{j=0}^{N-1}f_j\,\ee^{-2\pi\ii jk/N},\quad k=0,\dots,N-1,\qquad
f_j=\sum_{k=0}^{N-1}\hat f_k\,\ee^{2\pi\ii jk/N},\quad j=0,\dots,N-1.
\]
Med $w_N=\ee^{-2\pi\ii/N}$ og Fouriermatrisen $\FN=\bigl(w_N^{jk}\bigr)_{j,k=0}^{N-1}$:
\[
\hat{\vb f}=\tfrac1N\FN\vb f,\qquad \vb f=\overline{\FN}\,\hat{\vb f},\qquad \FN^{-1}=\tfrac1N\overline{\FN},\qquad \FN=\FN^\T.
\]
\end{formel}
\begin{teorem}[den sentrale summen (øving 8, oppg.~6)]
$\displaystyle\sum_{k=0}^{N-1}w_N^{jk}=\begin{cases}N,& j\equiv0\pmod N,\\0,&\text{ellers}\end{cases}$ \quad (geometrisk rekke med $q=w_N^j$, der $q^N=1$).
Det gir $\frac1N\overline{\FN}\FN=I$, altså at IDFT inverterer DFT.
\end{teorem}
\begin{eksempel}[små Fouriermatriser (kontrollregnet)]
\[
\mathcal F_2=\begin{pmatrix}1&1\\1&-1\end{pmatrix},\qquad
\mathcal F_4=\begin{pmatrix}1&1&1&1\\1&-\ii&-1&\ii\\1&-1&1&-1\\1&\ii&-1&-\ii\end{pmatrix}\quad(w_4=-\ii).
\]
For $N=5$ (Matte 4N vår 2023) skriver du $w_5^{jk}$ og reduserer eksponenten modulo 5: $(\mathcal F_5)_{jk}=w_5^{(jk\bmod 5)}$.
\end{eksempel}

\subsection{Egenskaper}
\begin{teorem}
\begin{itemize}
  \item \textbf{$N$-periodisk:} $\hat f_{k+N}=\hat f_k$ for alle $k\in\Z$, siden $w_N^N=1$. Indeksene $k=0,\dots,N-1$ (eller $-\frac N2,\dots,\frac N2-1$) er derfor nok.
  \item \textbf{Reelt signal:} $\hat f_{-k}=\overline{\hat f_k}$, altså $\hat f_{N-k}=\overline{\hat f_k}$. Spekteret er speilet, så bare halvparten trengs.
  \item \textbf{Aliasing-formelen:} er $f$ kontinuerlig med $\sum_k\abs{c_k(f)}<\infty$, så er
  \[
  \hat f_k=\sum_{\ell\in\Z}c_{k+\ell N}(f).
  \]
  På gitteret kan ikke frekvensene $k$ og $k+N$ skilles fra hverandre, fordi $\ee^{\ii(k+N)x_j}=\ee^{\ii kx_j}$.
\end{itemize}
\end{teorem}

\begin{teorem}[båndbegrensning, rekonstruksjon og Nyquist]
$f$ er \emph{båndbegrenset} med $B$ hvis $c_k(f)=0$ for $\abs k>B$. Hvis $N\ge2B+1$, er det høyst ett ledd i aliasing-summen for hver $k$, og
\[
c_k(f)=\hat f_k\ \ (0\le k\le B),\qquad c_{-k}(f)=\hat f_{N-k}\ \ (1\le k\le B),\qquad c_k(f)=0\ \text{ellers}.
\]
Da kan $f$ rekonstrueres eksakt fra prøvene (trigonometrisk interpolasjon). $2m+1$ prøver for høyeste frekvens $m$ kalles Nyquist-raten. Er $N\le 2B$, finnes det flere båndbegrensede funksjoner med de samme prøvene, og rekonstruksjonen er ikke entydig.
\end{teorem}
\begin{eksempel}[prøveeksamen, oppg.~7, kort]
$N=6$ og $\hat{\vb f}=(2,1,0,0,0,1)$. Med $B=2\le\frac{N-1}2$ blir $c_0=2$, $c_1=\hat f_1=1$ og $c_{-1}=\hat f_5=1$, så $f(x)=2+2\cos x$. Med $B=4$ er $N=6<2\cdot4+1$, og $f$ er ikke entydig.
\end{eksempel}

\begin{oppskrift}[tolke og filtrere et spekter (øving 8, oppg.~3)]
\begin{itemize}
  \item Indeks $k>\frac N2$ svarer til negativ frekvens $k-N$. \emph{fftshift} flytter dem slik at rekkefølgen blir $-\frac N2,\dots,\frac N2-1$.
  \item Med prøveavstand $\Delta t$ svarer indeks $k$ til frekvensen $\frac{k}{N\Delta t}$ Hz. Plott $\abs{\hat f_k}$ for $k=0,\dots,\lfloor N/2\rfloor$.
  \item \textbf{Filtrering eller støyfjerning:} sett $\hat f_k=0$ for uønskede frekvenser (eller der $\abs{\hat f_k}$ er liten) og ta IDFT. For et reelt signal må både $k$ og $N-k$ fjernes, ellers blir resultatet komplekst.
\end{itemize}
\end{oppskrift}

\subsection{Fast Fourier Transform (FFT)}
\begin{itemize}
  \item DFT som matrise-vektor-produkt koster $O(N^2)$. Med lyd i 44,1\,kHz er $N=44\,100$ per sekund, og det blir for dyrt.
  \item \textbf{Radix-2 (Cooley--Tukey, $N=2^n$, $M=N/2$):} del summen i første og andre halvdel, og se på like og odde $k$ hver for seg (uten faktoren $\frac1N$):
  \[
  \hat f_{2\kappa}=\sum_{j=0}^{M-1}\bigl(f_j+f_{j+M}\bigr)w_M^{j\kappa},\qquad
  \hat f_{2\kappa+1}=\sum_{j=0}^{M-1}\bigl(f_j-f_{j+M}\bigr)w_N^{j}\,w_M^{j\kappa}.
  \]
  Det gir to DFT-er av lengde $N/2$, $N$ addisjoner og $N/2$ multiplikasjoner. Gjentatt $n=\log_2N$ ganger gir det totalt $O(N\log N)$ operasjoner. Formlene er kontrollregnet.
  \item \textbf{Historie:} Gauss (ca.\ 1801, banen til Ceres, $N=12=4\cdot3$), Runge (1905) og Cooley--Tukey (1965). Andre faktoriseringer $N=N_1N_2$ fungerer også.
\end{itemize}

\subsection{Diskret sinus- og cosinustransform (DST og DCT)}
Samme idé som like og odde utvidelser. Et reelt signal gir reelle koeffisienter. Det finnes fire varianter av hver (I--IV), avhengig av hvordan signalet utvides. DCT-II er mest brukt, blant annet i bildekomprimering (JPEG). Å kaste små koeffisienter gir sterk komprimering, som i Barbara-eksempelet fra forelesningen (172\,kB ned til 3\,kB).

\begin{oving}
\begin{itemize}[leftmargin=5.6em,labelwidth=5em,labelsep=0.6em,align=left]
\item[\lett] \href{https://github.com/gustavbv/Matte4C---eksamen/blob/main/oppgavesett/F07-dft-og-fft/2-Bergmann-eksamener/010-TMA4125-2022-var-oppg7-oppgave.pdf}{Matte 4N vår 2022, oppg.~7} {\color{graa}(\href{https://github.com/gustavbv/Matte4C---eksamen/blob/main/oppgavesett/F07-dft-og-fft/2-Bergmann-eksamener/010-TMA4125-2022-var-oppg7-svar.pdf}{LF})}: $\mathcal F_4$, DFT og invers DFT
\item[\lett] \href{https://github.com/gustavbv/Matte4C---eksamen/blob/main/oppgavesett/F07-dft-og-fft/2-Bergmann-eksamener/008-TMA4125-30-2023-kont-oppg3-oppgave.pdf}{Matte 4N kont.\ 2023, oppg.~3} {\color{graa}(\href{https://github.com/gustavbv/Matte4C---eksamen/blob/main/oppgavesett/F07-dft-og-fft/2-Bergmann-eksamener/008-TMA4125-30-2023-kont-oppg3-svar.pdf}{LF})}: DFT av samplet $\ee^{3\ii x}$
\item[\middels] \href{https://github.com/gustavbv/Matte4C---eksamen/blob/main/oppgavesett/F07-dft-og-fft/1-TMA4432-ovinger-og-proveeksamen/001-TMA4432-2026-proveeksamen-oppg7-oppgave.pdf}{Prøveeksamen, oppg.~7} {\color{graa}(\href{https://github.com/gustavbv/Matte4C---eksamen/blob/main/oppgavesett/F07-dft-og-fft/1-TMA4432-ovinger-og-proveeksamen/001-TMA4432-2026-proveeksamen-oppg7-svar.pdf}{LF})}: båndbegrensning og rekonstruksjon
\item[\middels] \href{https://github.com/gustavbv/Matte4C---eksamen/blob/main/oppgavesett/F07-dft-og-fft/2-Bergmann-eksamener/007-TMA4125-2023-var-oppg4-oppgave.pdf}{Matte 4N vår 2023, oppg.~4} {\color{graa}(\href{https://github.com/gustavbv/Matte4C---eksamen/blob/main/oppgavesett/F07-dft-og-fft/2-Bergmann-eksamener/007-TMA4125-2023-var-oppg4-svar.pdf}{LF})}: $\mathcal F_5$, båndbegrensning og aliasing
\item[\middels] \href{https://github.com/gustavbv/Matte4C---eksamen/blob/main/oppgavesett/F07-dft-og-fft/2-Bergmann-eksamener/005-TMA4125-30-2025-kont-oppg5-oppgave.pdf}{Matte 4N kont.\ 2025, oppg.~5} {\color{graa}(\href{https://github.com/gustavbv/Matte4C---eksamen/blob/main/oppgavesett/F07-dft-og-fft/2-Bergmann-eksamener/005-TMA4125-30-2025-kont-oppg5-svar.pdf}{LF})}: aliasing og Nyquist for $3+2\cos 2x$
\item[\vanskelig] \href{https://github.com/gustavbv/Matte4C---eksamen/blob/main/oppgavesett/F07-dft-og-fft/1b-TMA4432-ekstraoppgaver/003-TMA4432-oving-8-oppg6-oppgave.pdf}{Øving 8, oppg.~6 (ekstra)} {\color{graa}(\href{https://github.com/gustavbv/Matte4C---eksamen/blob/main/oppgavesett/F07-dft-og-fft/1b-TMA4432-ekstraoppgaver/003-TMA4432-oving-8-oppg6-svar.pdf}{LF})}: bevis: periodisitet og summen av $w_N^{jk}$
\item[\vanskelig] \href{https://github.com/gustavbv/Matte4C---eksamen/blob/main/oppgavesett/F07-dft-og-fft/3-hele-oppgaven-i-pensum/012-TMA4130-35-2023-host-oppg7-oppgave.pdf}{Matte 4N/4D høst 2023, oppg.~7} {\color{graa}(\href{https://github.com/gustavbv/Matte4C---eksamen/blob/main/oppgavesett/F07-dft-og-fft/3-hele-oppgaven-i-pensum/012-TMA4130-35-2023-host-oppg7-svar.pdf}{LF})}: tolke et spekter og kode for et filter
\end{itemize}
\end{oving}


\section{F8 · Fouriertransformen}

\textbf{Kjernen:} Fourierrekker for funksjoner som ikke er periodiske, ved å la $L\to\infty$. Summen over $k$ blir et integral over frekvensen $\omega$. Derivasjon blir multiplikasjon med $\ii\omega$, og konvolusjon blir produkt. Det gjør PDE-er på hele $\R$ om til ODE-er (F10).

\begin{formel}[Fouriertransform og invers\fs]
For $f\in L^1(\R)$ (det vil si $\int_{-\infty}^\infty\abs f<\infty$):
\begin{gather*}
\hat f(\omega)=\Four(f)(\omega)=\frac1{\sqrt{2\pi}}\int_{-\infty}^{\infty}f(x)\ee^{-\ii\omega x}\dd x,\\
f(x)=\Four^{-1}(\hat f)(x)=\frac1{\sqrt{2\pi}}\int_{-\infty}^{\infty}\hat f(\omega)\ee^{\ii\omega x}\dd\omega .
\end{gather*}
Inversjonen gjelder når både $f$ og $\hat f$ er absolutt integrerbare. $f$ kalles tidsdomenet og $\hat f$ frekvensdomenet.
\end{formel}
\begin{fallgruve}[konvensjoner]
Andre bøker plasserer $2\pi$ annerledes: $\int f\ee^{-\ii\omega x}$ (med $\frac1{2\pi}$ i inversen), $\frac1{2\pi}\int\dots$ eller $\int f\ee^{-2\pi\ii\omega x}$ (frekvens i Hz), og noen har minuset i inversen. Sjekk alltid hvilken definisjon en tabell eller en gammel eksamen bruker. Kurset bruker $\frac1{\sqrt{2\pi}}$ begge veier.
\end{fallgruve}

\begin{defn}[cosinus- og sinustransform]
For like $f$: $\hat f_c(\omega)=\sqrt{\frac2\pi}\int_0^\infty f(x)\cos(\omega x)\dd x$, med invers $\check g_c(x)=\sqrt{\frac2\pi}\int_0^\infty g(\omega)\cos(\omega x)\dd\omega$.\\
For odde $f$: $\hat f_s(\omega)=\sqrt{\frac2\pi}\int_0^\infty f(x)\sin(\omega x)\dd x$, med invers på samme form.\\
En reell funksjon gir en reell transform. Sammenhengen er $\hat f=\hat f_c$ for like $f$ og $\hat f=-\ii\hat f_s$ for odde $f$.
\end{defn}

\begin{formel}[Transformtabellen\fs]
\begin{center}
\begin{tabular}{@{}ll@{}}
\toprule
$f(x)$ & $\hat f(\omega)$\\
\midrule
$\ee^{-ax^2}$ \ ($a>0$) & $\dfrac1{\sqrt{2a}}\ee^{-\omega^2/(4a)}$ \quad (gaussisk gir gaussisk)\\[6pt]
$\ee^{-a\abs x}$ \ ($a>0$) & $\sqrt{\dfrac2\pi}\,\dfrac{a}{\omega^2+a^2}$\\[6pt]
$1$ for $\abs x<a$, ellers $0$ & $\sqrt{\dfrac2\pi}\,\dfrac{\sin(a\omega)}{\omega}$\\[6pt]
$\sqrt{\dfrac2\pi}\,\dfrac{\sin(ax)}x$ & $1$ for $\abs\omega<a$, ellers $0$\\
\bottomrule
\end{tabular}
\end{center}
Ikke på arket, men nyttig: $\chi_{[a,b]}$ (1 på $[a,b]$, ellers 0) har $\hat f(\omega)=\dfrac{\ee^{-\ii\omega a}-\ee^{-\ii\omega b}}{\ii\omega\sqrt{2\pi}}$ for $\omega\neq0$ og $\hat f(0)=\dfrac{b-a}{\sqrt{2\pi}}$.
\end{formel}

\begin{teorem}[regneregler (de tre første er fra forelesningen)]
\begin{itemize}
  \item \textbf{Linearitet:} $\Four(af+bg)=a\Four(f)+b\Four(g)$.
  \item \textbf{Derivasjon:} hvis $f,f'\in L^1$ og $f\to0$ for $x\to\pm\infty$, er $\Four(f')=\ii\omega\,\Four(f)$, og dermed $\Four(f'')=-\omega^2\Four(f)$.
  \item \textbf{Konvolusjon:} $(f*g)(x)=\displaystyle\int_{-\infty}^{\infty}f(y)g(x-y)\dd y$, med $f*g=g*f$ og
  \[
  \Four(f*g)=\sqrt{2\pi}\;\Four(f)\,\Four(g).
  \]
  Konvolusjon i tid blir multiplikasjon i frekvens.
\end{itemize}
\textbf{Ekstra (ikke på lysarkene, kontrollregnet):} forskyvning $\Four(f(x-a))=\ee^{-\ii\omega a}\hat f(\omega)$, modulasjon $\Four(\ee^{\ii ax}f)=\hat f(\omega-a)$, skalering $\Four(f(bx))=\frac1{\abs b}\hat f(\omega/b)$ og $\Four(xf(x))=\ii\,\frac{\dd}{\dd\omega}\hat f(\omega)$.
\end{teorem}

\begin{teorem}[gaussiske integraler (øving 8, oppg.~5)]
$\displaystyle\int_{-\infty}^\infty\ee^{-x^2}\dd x=\sqrt\pi$ (kvadrer og bytt til polarkoordinater), og
$\displaystyle\int_{-\infty}^\infty\ee^{-ax^2+bx}\dd x=\sqrt{\frac\pi a}\,\ee^{b^2/(4a)}$ for $a>0$. $\ee^{-x^2/2}$ er sin egen transform.
\end{teorem}

\begin{oppskrift}[typiske oppgaver]
\begin{enumerate}
  \item \textbf{FT fra definisjonen:} for stykkevise funksjoner, integrer over intervallene der $f\neq0$ og bruk delvis integrasjon for $x^n$-ledd. Behandle $\omega=0$ for seg. Sjekk med symmetri: like $f$ gir reell og like $\hat f$, og odde $f$ gir rent imaginær og odde $\hat f$.
  \item \textbf{Med tabell og regler:} skriv $f$ som en kombinasjon av tabellfunksjoner. For eksempel gir $t^2\ee^{-t^2}$ at du deriverer $\ee^{-t^2}$ to ganger, eller bruker $\Four(x^2f)=-\hat f''$.
  \item \textbf{Konvolusjon:} $f*g=\Four^{-1}\bigl(\sqrt{2\pi}\hat f\hat g\bigr)$. For eksempel:
  \begin{itemize}
    \item $\ee^{-x^2}*\ee^{-2x^2}=\sqrt{\pi/3}\,\ee^{-2x^2/3}$ (øving 7, oppg.~5, kontrollregnet).
    \item $\operatorname{sinc}*\operatorname{sinc}=\pi\operatorname{sinc}$ med $\operatorname{sinc}x=\frac{\sin x}x$ (øving 8, oppg.~4).
    \item $\chi_{(0,2)}*\chi_{(0,2)}$ er en trekantfunksjon (øving 8, oppg.~2).
  \end{itemize}
\end{enumerate}
\end{oppskrift}

\begin{oving}
\begin{itemize}[leftmargin=5.6em,labelwidth=5em,labelsep=0.6em,align=left]
\item[\lett] \href{https://github.com/gustavbv/Matte4C---eksamen/blob/main/oppgavesett/F08-fouriertransform/2-Bergmann-eksamener/010-TMA4125-30-35-2022-kont-oppg7-oppgave.pdf}{Matte 4N kont.\ 2022, oppg.~7} {\color{graa}(\href{https://github.com/gustavbv/Matte4C---eksamen/blob/main/oppgavesett/F08-fouriertransform/2-Bergmann-eksamener/010-TMA4125-30-35-2022-kont-oppg7-svar.pdf}{LF})}: konvolusjon av gaussfunksjoner med tabellen
\item[\lett] \href{https://github.com/gustavbv/Matte4C---eksamen/blob/main/oppgavesett/F08-fouriertransform/3-hele-oppgaven-i-pensum/046-TMA4135-2015-host-oppg2-oppgave.pdf}{TMA4135 høst 2015, oppg.~2} {\color{graa}(\href{https://github.com/gustavbv/Matte4C---eksamen/blob/main/oppgavesett/F08-fouriertransform/3-hele-oppgaven-i-pensum/046-TMA4135-2015-host-oppg2-svar.pdf}{LF})}: $f_a*f_b$ for gaussfunksjoner
\item[\middels] \href{https://github.com/gustavbv/Matte4C---eksamen/blob/main/oppgavesett/F08-fouriertransform/2-Bergmann-eksamener/009-TMA4125-2022-var-oppg6-oppgave.pdf}{Matte 4N vår 2022, oppg.~6} {\color{graa}(\href{https://github.com/gustavbv/Matte4C---eksamen/blob/main/oppgavesett/F08-fouriertransform/2-Bergmann-eksamener/009-TMA4125-2022-var-oppg6-svar.pdf}{LF})}: FT av en avkortet cosinus
\item[\middels] \href{https://github.com/gustavbv/Matte4C---eksamen/blob/main/oppgavesett/F08-fouriertransform/1-TMA4432-ovinger-og-proveeksamen/006-TMA4432-oving-8-oppg2-oppgave.pdf}{Øving 8, oppg.~2} {\color{graa}(\href{https://github.com/gustavbv/Matte4C---eksamen/blob/main/oppgavesett/F08-fouriertransform/1-TMA4432-ovinger-og-proveeksamen/006-TMA4432-oving-8-oppg2-svar.pdf}{LF})}: $f*f$ fra definisjonen, og FT av $f$ og $g$
\item[\middels] \href{https://github.com/gustavbv/Matte4C---eksamen/blob/main/oppgavesett/F08-fouriertransform/3-hele-oppgaven-i-pensum/013-TMA4130-35-2025-kont-2-sett-oppg4-oppgave.pdf}{TMA4130 kont.\ 2025 (2.\ sett), oppg.~4} {\color{graa}(\href{https://github.com/gustavbv/Matte4C---eksamen/blob/main/oppgavesett/F08-fouriertransform/3-hele-oppgaven-i-pensum/013-TMA4130-35-2025-kont-2-sett-oppg4-svar.pdf}{LF})}: FT av stykkevis funksjon og derivasjonsregelen
\item[\vanskelig] \href{https://github.com/gustavbv/Matte4C---eksamen/blob/main/oppgavesett/F08-fouriertransform/1-TMA4432-ovinger-og-proveeksamen/005-TMA4432-oving-8-oppg1-oppgave.pdf}{Øving 8, oppg.~1} {\color{graa}(\href{https://github.com/gustavbv/Matte4C---eksamen/blob/main/oppgavesett/F08-fouriertransform/1-TMA4432-ovinger-og-proveeksamen/005-TMA4432-oving-8-oppg1-svar.pdf}{LF})}: FT av $x^2$ på $\abs x<1$ (mye delvis integrasjon)
\item[\vanskelig] \href{https://github.com/gustavbv/Matte4C---eksamen/blob/main/oppgavesett/F08-fouriertransform/3-hele-oppgaven-i-pensum/012-TMA4130-35-2025-kont-oppg7-oppgave.pdf}{TMA4130 kont.\ 2025, oppg.~7} {\color{graa}(\href{https://github.com/gustavbv/Matte4C---eksamen/blob/main/oppgavesett/F08-fouriertransform/3-hele-oppgaven-i-pensum/012-TMA4130-35-2025-kont-oppg7-svar.pdf}{LF})}: skaleringsregel og produkt som konvolusjon
\item[\vanskelig] \href{https://github.com/gustavbv/Matte4C---eksamen/blob/main/oppgavesett/F08-fouriertransform/3-hele-oppgaven-i-pensum/059-TMA4135-2011-kont-oppg2-oppgave.pdf}{TMA4135 kont.\ 2011, oppg.~2} {\color{graa}(\href{https://github.com/gustavbv/Matte4C---eksamen/blob/main/oppgavesett/F08-fouriertransform/3-hele-oppgaven-i-pensum/059-TMA4135-2011-kont-oppg2-svar.pdf}{LF})}: delbrøk, FT og et vanskelig integral
\end{itemize}
\end{oving}


\section{F9 · Partielle differensiallikninger og bølgelikningen}

\textbf{Kjernen:} Klassifiser PDE-en, still opp rand- og initialbetingelser, prøv $u=F(x)G(t)$ (separasjon av variable), finn egenfunksjonene fra randbetingelsene og bruk superposisjon og Fourierrekker for å tilpasse initialbetingelsene.

\subsection{Begreper og klassifisering}
\begin{defn}
En PDE er en likning med partiellderiverte av en ukjent funksjon $u$ av minst to variable, for eksempel $u(x,t)$ eller $u(x,y,t)$.
\begin{itemize}
  \item \textbf{Orden:} høyeste deriverte som forekommer.
  \item \textbf{Lineær:} $u$ og de deriverte forekommer bare i første potens og ikke multiplisert med hverandre. Koeffisientene kan avhenge av $x$ og $t$. Ellers er den \textbf{ikke-lineær}, for eksempel med $u^2$, $uu_x$ eller $\sin u$.
  \item \textbf{Homogen:} hvert ledd inneholder $u$ eller en derivert av $u$. Da er $u\equiv0$ en løsning. Et ledd uten $u$, som $f(x,t)$, en konstant $C\neq0$ eller $x$, gjør likningen \textbf{inhomogen}.
\end{itemize}
\end{defn}
\begin{eksempel}[øving 9, oppg.~1]
\begin{tabular}{@{}ll@{}}
$u_x+u^2=x$ & ikke-lineær ($u^2$), inhomogen ($x$), 1.\ orden\\
$u_t=u_{xx}$ & lineær, homogen, 2.\ orden (varmelikningen)\\
$\frac1r\partial_r(r\partial_ru)+\partial_\theta^2u=0$ & lineær, homogen (Laplace i polarkoordinater)\\
$u_{xx}+u_{yy}+C=0$ & lineær, inhomogen for $C\neq0$ (Poisson)\\
$u_t+u_x+xu^n=0$ & lineær hvis $n=1$, ellers ikke-lineær, homogen\\
$u_t+uu_x=u_{xx}$ & ikke-lineær ($uu_x$), homogen (Burgers)
\end{tabular}
\end{eksempel}

\begin{defn}[viktige likninger]
\begin{tabular}{@{}ll@{}}
Varmelikningen & $u_t-\alpha u_{xx}=q(x,t)$\\
Bølgelikningen & $u_{tt}=c^2u_{xx}$ ($c$ er bølgefarten), i 2D $u_{tt}=c^2(u_{xx}+u_{yy})$\\
Bjelke eller stav & $u_{tt}+k\,u_{xxxx}=q(x,t)$ (4.\ orden)\\
Laplace og Poisson & $\Delta u=u_{xx}+u_{yy}=0$ \ og \ $\Delta u=f$
\end{tabular}
\end{defn}

\begin{defn}[løsning, randbetingelser og initialbetingelser]
En løsning er en funksjon med de nødvendige deriverte som oppfyller likningen i området. For å få én bestemt løsning trengs \textbf{randbetingelser} (RB, verdier på randen av området) og \textbf{initialbetingelser} (IB, ved $t=0$). Med $\partial_t^k$ som høyeste tidsderiverte trengs $k$ initialbetingelser: én for varmelikningen og to (posisjon og hastighet) for bølgelikningen.
\end{defn}
\begin{teorem}[superposisjonsprinsippet]
Er $u_1$ og $u_2$ løsninger av en \emph{lineær homogen} PDE, så er $au_1+bu_2$ også en løsning. For en inhomogen likning gjelder ikke dette, men løsningen er en partikulær løsning pluss en homogen løsning (øving 9, oppg.~2).
\end{teorem}

\subsection{Bølgelikningen}
Modell: en streng på $[0,L]$ som er festet i endene. Strengen er homogen, tyngdekraften neglisjeres, og hvert punkt beveger seg bare loddrett. Bølgefarten er $c^2=\frac{\text{strekkraft}}{\text{masse per lengde}}$, så $c^2=\frac{FL}m$ i øving 9.

\begin{formel}[{Bølgelikningen på $[0,L]$ og løsningen\fs}]
\[
\begin{cases}
u_{tt}=c^2u_{xx}, & 0<x<L,\ t>0,\\
u(0,t)=u(L,t)=0 & \text{(RB)},\\
u(x,0)=f(x),\quad u_t(x,0)=g(x) & \text{(IB: posisjon og hastighet)},
\end{cases}
\]
\begin{gather*}
u(x,t)=\sum_{n=1}^\infty\Bigl(A_n\cos\frac{cn\pi t}{L}+B_n\sin\frac{cn\pi t}L\Bigr)\sin\frac{n\pi x}L,\\
A_n=\frac2L\int_0^Lf(x)\sin\frac{n\pi x}L\dd x,\qquad B_n=\frac2{cn\pi}\int_0^Lg(x)\sin\frac{n\pi x}L\dd x.
\end{gather*}
\end{formel}

\begin{oppskrift}[separasjon av variable (utledningen må kunnes)]
\begin{enumerate}
  \item \textbf{Ansats:} $u(x,t)=F(x)G(t)$. Innsatt gir det $FG''=c^2F''G$, altså
  \[
  \frac{F''(x)}{F(x)}=\frac{G''(t)}{c^2G(t)}=k .
  \]
  $k$ er en konstant, fordi venstre side bare avhenger av $x$ og høyre bare av $t$.
  Det gir to ODE-er: $F''-kF=0$ og $G''-c^2kG=0$. Randbetingelsene gir $F(0)=F(L)=0$ (ellers er $G\equiv0$).
  \item \textbf{Tilfellet $k=0$:} $F=Ax+B$, og RB gir $F\equiv0$.
  \item \textbf{Tilfellet $k=p^2>0$:} $F=A\ee^{px}+B\ee^{-px}$, og RB gir $A=B=0$.
  \item \textbf{Tilfellet $k=-p^2<0$:} $F=A\cos px+B\sin px$. $F(0)=0$ gir $A=0$, og $F(L)=0$ gir $\sin pL=0$, så $p=\frac{n\pi}L$. Altså er $F_n=\sin\frac{n\pi x}L$.
  \item \textbf{Tidsdelen:} $G_n''+\lambda_n^2G_n=0$ med $\lambda_n=\frac{cn\pi}L$, så $G_n=A_n\cos\lambda_nt+B_n\sin\lambda_nt$.
  \item \textbf{Superposisjon:} $u=\sum_nG_n(t)F_n(x)$. Sett $t=0$: $f(x)=\sum A_n\sin\frac{n\pi x}L$ er sinusrekken til $f$ (odde utvidelse). Deriver i $t$: $g(x)=\sum B_n\lambda_n\sin\frac{n\pi x}L$.
\end{enumerate}
\textbf{Begreper:} $u_n$ er egenfunksjonene, $\lambda_n=\frac{cn\pi}L$ egenverdiene, og $\{\lambda_n\}$ spekteret. Grunntonen har periode $\frac{2\pi}{\lambda_1}=\frac{2L}c$.
\end{oppskrift}

\begin{fallgruve}[effektive triks]
\begin{itemize}
  \item \textbf{Les av:} er $f$ en endelig sum av $\sin\frac{n\pi x}L$, er $A_n$ bare koeffisientene. For eksempel har $f=\sin(4\pi x)+3\sin(7\pi x)$ på $[0,1]$ $A_4=1$, $A_7=3$ og ellers $0$, så du slipper å integrere.
  \item \textbf{Er $g\equiv0$,} er alle $B_n=0$.
  \item \textbf{Et annet intervall,} for eksempel $[-L,L]$ med $u(\pm L,t)=0$ (prøveeksamen, oppg.~8): bruk $F=A\cos px+B\sin px$. Da gir $A\cos pL=0$ og $B\sin pL=0$ to familier: $\sin\frac{n\pi x}L$ og $\cos\frac{(2m-1)\pi x}{2L}$. Alternativt kan du forskyve til $[0,2L]$.
  \item \textbf{Neumann-rand} ($u_x=0$ i endene) gir $\cos\frac{n\pi x}L$, inkludert en konstant mode for $n=0$.
\end{itemize}
\end{fallgruve}

\begin{eksempel}[strengen løftes på midten (forelesningen)]
$f(x)=H\bigl(1-\abs{\frac{2x}L-1}\bigr)$ (trekant med topp $H$ i $x=\frac L2$) og $g=0$. Da er $B_n=0$ og $A_n=\frac{8H}{n^2\pi^2}\sin\frac{n\pi}{2}$, så bare odde $n$ bidrar, med vekslende fortegn.
\end{eksempel}

\begin{oving}
\begin{itemize}[leftmargin=5.6em,labelwidth=5em,labelsep=0.6em,align=left]
\item[\lett] \href{https://github.com/gustavbv/Matte4C---eksamen/blob/main/oppgavesett/F09-pde-og-bolgelikning/1-TMA4432-ovinger-og-proveeksamen/003-TMA4432-oving-9-oppg1-oppgave.pdf}{Øving 9, oppg.~1} {\color{graa}(\href{https://github.com/gustavbv/Matte4C---eksamen/blob/main/oppgavesett/F09-pde-og-bolgelikning/1-TMA4432-ovinger-og-proveeksamen/003-TMA4432-oving-9-oppg1-svar.pdf}{LF})}: klassifisering av PDE-er
\item[\middels] \href{https://github.com/gustavbv/Matte4C---eksamen/blob/main/oppgavesett/F09-pde-og-bolgelikning/2-Bergmann-eksamener/011-TMA4125-2024-var-oppg4-oppgave.pdf}{Matte 4N vår 2024, oppg.~4} {\color{graa}(\href{https://github.com/gustavbv/Matte4C---eksamen/blob/main/oppgavesett/F09-pde-og-bolgelikning/2-Bergmann-eksamener/011-TMA4125-2024-var-oppg4-svar.pdf}{LF})}: klassifisering og verifisering
\item[\middels] \href{https://github.com/gustavbv/Matte4C---eksamen/blob/main/oppgavesett/F09-pde-og-bolgelikning/1-TMA4432-ovinger-og-proveeksamen/001-TMA4432-2026-proveeksamen-oppg8-oppgave.pdf}{Prøveeksamen, oppg.~8} {\color{graa}(\href{https://github.com/gustavbv/Matte4C---eksamen/blob/main/oppgavesett/F09-pde-og-bolgelikning/1-TMA4432-ovinger-og-proveeksamen/001-TMA4432-2026-proveeksamen-oppg8-svar.pdf}{LF})}: separasjon på $[-1,1]$: sinus- og cosinusmoder
\item[\middels] \href{https://github.com/gustavbv/Matte4C---eksamen/blob/main/oppgavesett/F09-pde-og-bolgelikning/3-hele-oppgaven-i-pensum/023-TMA4130-35-2025-kont-oppg6-oppgave.pdf}{TMA4130 kont.\ 2025, oppg.~6} {\color{graa}(\href{https://github.com/gustavbv/Matte4C---eksamen/blob/main/oppgavesett/F09-pde-og-bolgelikning/3-hele-oppgaven-i-pensum/023-TMA4130-35-2025-kont-oppg6-svar.pdf}{LF})}: streng sluppet fra trekantform
\item[\vanskelig] \href{https://github.com/gustavbv/Matte4C---eksamen/blob/main/oppgavesett/F09-pde-og-bolgelikning/1-TMA4432-ovinger-og-proveeksamen/005-TMA4432-oving-9-oppg3-oppgave.pdf}{Øving 9, oppg.~3} {\color{graa}(\href{https://github.com/gustavbv/Matte4C---eksamen/blob/main/oppgavesett/F09-pde-og-bolgelikning/1-TMA4432-ovinger-og-proveeksamen/005-TMA4432-oving-9-oppg3-svar.pdf}{LF})}: trapesformet startprofil og svingeperiode
\item[\vanskelig] \href{https://github.com/gustavbv/Matte4C---eksamen/blob/main/oppgavesett/F09-pde-og-bolgelikning/2-Bergmann-eksamener/017-TMA4125-2022-var-oppg9-oppgave.pdf}{Matte 4N vår 2022, oppg.~9} {\color{graa}(\href{https://github.com/gustavbv/Matte4C---eksamen/blob/main/oppgavesett/F09-pde-og-bolgelikning/2-Bergmann-eksamener/017-TMA4125-2022-var-oppg9-svar.pdf}{LF})}: separasjon for en fjerdeordens PDE
\end{itemize}
\end{oving}


\section{F10 · d'Alembert, varmelikningen og differansemetoder}

\textbf{Kjernen:} To analytiske metoder til (d'Alembert og Fouriertransform på hele $\R$), varmelikningen løst med separasjon, og numerikk. Vi bytter deriverte med differanser på et gitter og får et skjema som går fremover i tid. Stabilitet krever små nok tidssteg (CFL, $\alpha\le\frac12$), med mindre skjemaet er implisitt.

\subsection{d'Alemberts løsning av bølgelikningen}
\begin{teorem}[d'Alembert]
Alle $u(x,t)=\varphi(x+ct)+\psi(x-ct)$ med $\varphi,\psi\in C^2$ løser $u_{tt}=c^2u_{xx}$: to bølger som går mot venstre og mot høyre med fart $c$. Initialbetingelsene $u(x,0)=f$ og $u_t(x,0)=g$ gir $\varphi(x)=\frac12\bigl(f(x)+\frac1c\int_{x_0}^xg\bigr)$ og $\psi(x)=\frac12\bigl(f(x)-\frac1c\int_{x_0}^xg\bigr)$, og dermed
\[
u(x,t)=\frac12\bigl(f(x+ct)+f(x-ct)\bigr)+\frac1{2c}\int_{x-ct}^{x+ct}g(s)\dd s .
\]
Integralleddet forsvinner hvis $g=0$. På $[0,L]$ med faste ender bruker du den odde $2L$-periodiske utvidelsen av $f$ og $g$. Da stemmer svaret med sinusrekkeløsningen.
\end{teorem}
\begin{eksempel}[øving 10, oppg.~1: $c=2$, $f=\sin x$, $g=\ee^{-x/4}$ (kontrollregnet)]
$u=\frac12\bigl(\sin(x+2t)+\sin(x-2t)\bigr)+\frac14\int_{x-2t}^{x+2t}\ee^{-s/4}\dd s=\sin x\cos2t+\ee^{-(x-2t)/4}-\ee^{-(x+2t)/4}$.
\end{eksempel}

\begin{oving}
\begin{itemize}[leftmargin=5.6em,labelwidth=5em,labelsep=0.6em,align=left]
\item[\lett] \href{https://github.com/gustavbv/Matte4C---eksamen/blob/main/oppgavesett/F10-varmelikning-og-differanser/2-Bergmann-eksamener/016-TMA4125-2024-var-oppg5-oppgave.pdf}{Matte 4N vår 2024, oppg.~5} {\color{graa}(\href{https://github.com/gustavbv/Matte4C---eksamen/blob/main/oppgavesett/F10-varmelikning-og-differanser/2-Bergmann-eksamener/016-TMA4125-2024-var-oppg5-svar.pdf}{LF})}: d'Alemberts formel på $\R$
\item[\lett] \href{https://github.com/gustavbv/Matte4C---eksamen/blob/main/oppgavesett/F10-varmelikning-og-differanser/1-TMA4432-ovinger-og-proveeksamen/003-TMA4432-oving-10-oppg1-oppgave.pdf}{Øving 10, oppg.~1} {\color{graa}(\href{https://github.com/gustavbv/Matte4C---eksamen/blob/main/oppgavesett/F10-varmelikning-og-differanser/1-TMA4432-ovinger-og-proveeksamen/003-TMA4432-oving-10-oppg1-svar.pdf}{LF})}: d'Alembert med $f=\sin x$, $g=\ee^{-x/4}$
\item[\middels] \href{https://github.com/gustavbv/Matte4C---eksamen/blob/main/oppgavesett/F10-varmelikning-og-differanser/3-hele-oppgaven-i-pensum/033-TMA4130-35-2024-host-oppg6-oppgave.pdf}{TMA4130 høst 2024, oppg.~6} {\color{graa}(\href{https://github.com/gustavbv/Matte4C---eksamen/blob/main/oppgavesett/F10-varmelikning-og-differanser/3-hele-oppgaven-i-pensum/033-TMA4130-35-2024-host-oppg6-svar.pdf}{LF})}: separert løsning sammenlignet med d'Alembert
\item[\vanskelig] \href{https://github.com/gustavbv/Matte4C---eksamen/blob/main/oppgavesett/F10-varmelikning-og-differanser/3-hele-oppgaven-i-pensum/031-TMA4125-30-35-2024-kont-oppg7-oppgave.pdf}{Matte 4N kont.\ 2024, oppg.~7} {\color{graa}(\href{https://github.com/gustavbv/Matte4C---eksamen/blob/main/oppgavesett/F10-varmelikning-og-differanser/3-hele-oppgaven-i-pensum/031-TMA4125-30-35-2024-kont-oppg7-svar.pdf}{LF})}: d'Alembert via Fouriertransformen
\end{itemize}
\end{oving}


\subsection{Varmelikningen}
\textbf{Utledning (i korte trekk):} energibevaring for et område $\Omega$ pluss divergensteoremet gir $\partial_te+\nabla\cdot\vb F=p$. Bruk to konstitutive lover: $e=e_0+\sigma u$ (varmekapasitet $\sigma$, relativ temperatur $u$) og Fouriers lov $\vb F=-\lambda\nabla u$ (varmeledningsevne $\lambda$). Da blir $\sigma u_t-\nabla\cdot(\lambda\nabla u)=p$. Med konstant $\lambda$ og $c^2=\lambda/\sigma$ får vi $u_t=c^2\Delta u+p$, som regel med $p=0$.

\begin{defn}[randbetingelser]
\begin{tabular}{@{}lll@{}}
Robin & $\lambda\partial_nu+\kappa(u-u_a)=0$ & varmeovergang til omgivelsestemperatur $u_a$\\
Neumann & $\partial_nu=0$ (homogen) eller $u_x=g$ & $\kappa=0$: perfekt isolert\\
Dirichlet & $u=u_a$ (eller $u=g(t)$) & $\kappa\to\infty$: temperaturen er gitt
\end{tabular}
\end{defn}

\begin{formel}[{Varmelikningen på $[0,L]$ (ikke på formelarket)}]
\textbf{Dirichlet} $u(0,t)=u(L,t)=0$ og $u(x,0)=f(x)$:
\[
u(x,t)=\sum_{n=1}^\infty B_n\,\ee^{-(cn\pi/L)^2t}\sin\frac{n\pi x}L,\qquad B_n=\frac2L\int_0^Lf(x)\sin\frac{n\pi x}L\dd x .
\]
\textbf{Neumann} $u_x(0,t)=u_x(L,t)=0$:
\begin{gather*}
u(x,t)=A_0+\sum_{n=1}^\infty A_n\,\ee^{-(cn\pi/L)^2t}\cos\frac{n\pi x}L,\\
A_0=\frac1L\int_0^Lf(x)\dd x,\qquad A_n=\frac2L\int_0^Lf(x)\cos\frac{n\pi x}L\dd x .
\end{gather*}
Utledningen er som for bølgelikningen, men med $G'=-c^2p^2G$, som gir $G=\ee^{-c^2p^2t}$. Når $t\to\infty$, går Dirichlet-løsningen mot $0$ og Neumann-løsningen mot middelverdien $A_0$. Høye frekvenser dempes raskest.
\end{formel}
\begin{oppskrift}[konstante temperaturer $u(0,t)=T_1$, $u(L,t)=T_2$ (Kreyszig, ekstra)]
Skriv $u=u_s+w$ med den stasjonære løsningen $u_s(x)=T_1+(T_2-T_1)\frac xL$. Da løser $w$ varmelikningen med null-RB og startverdi $f-u_s$. Tilsvarende kan du løse inhomogene likninger ved å finne en partikulær løsning og legge til en homogen.
\end{oppskrift}

\begin{formel}[Stasjonær varme: Laplace-likningen på et rektangel (Kreyszig 12.6)]
$u_{xx}+u_{yy}=0$ på $0<x<a$, $0<y<b$, med $u=0$ på tre sider og $u(x,b)=f(x)$:
\[
u(x,y)=\sum_{n=1}^\infty A_n\sin\frac{n\pi x}a\,\sinh\frac{n\pi y}a,\qquad A_n=\frac{2}{a\sinh(n\pi b/a)}\int_0^af(x)\sin\frac{n\pi x}a\dd x .
\]
Separasjonen $u=F(x)G(y)$ gir $F''=-p^2F$ (sinus) og $G''=p^2G$ (sinh, fordi $G(0)=0$). Er flere sider ulik null, løs ett problem per side og legg sammen (superposisjon). Kontrollregnet.
\end{formel}

\begin{oving}
\begin{itemize}[leftmargin=5.6em,labelwidth=5em,labelsep=0.6em,align=left]
\item[\middels] \href{https://github.com/gustavbv/Matte4C---eksamen/blob/main/oppgavesett/F09-pde-og-bolgelikning/2-Bergmann-eksamener/010-TMA4125-30-2025-kont-oppg6-oppgave.pdf}{Matte 4N kont.\ 2025, oppg.~6} {\color{graa}(\href{https://github.com/gustavbv/Matte4C---eksamen/blob/main/oppgavesett/F09-pde-og-bolgelikning/2-Bergmann-eksamener/010-TMA4125-30-2025-kont-oppg6-svar.pdf}{LF})}: varmelikningen med separasjon
\item[\middels] \href{https://github.com/gustavbv/Matte4C---eksamen/blob/main/oppgavesett/F09-pde-og-bolgelikning/3-hele-oppgaven-i-pensum/022-TMA4130-35-2025-kont-oppg5-oppgave.pdf}{TMA4130 kont.\ 2025, oppg.~5} {\color{graa}(\href{https://github.com/gustavbv/Matte4C---eksamen/blob/main/oppgavesett/F09-pde-og-bolgelikning/3-hele-oppgaven-i-pensum/022-TMA4130-35-2025-kont-oppg5-svar.pdf}{LF})}: isolerte ender (Neumann)
\item[\vanskelig] \href{https://github.com/gustavbv/Matte4C---eksamen/blob/main/oppgavesett/F09-pde-og-bolgelikning/2-Bergmann-eksamener/016-TMA4125-2022-var-oppg8-oppgave.pdf}{Matte 4N vår 2022, oppg.~8} {\color{graa}(\href{https://github.com/gustavbv/Matte4C---eksamen/blob/main/oppgavesett/F09-pde-og-bolgelikning/2-Bergmann-eksamener/016-TMA4125-2022-var-oppg8-svar.pdf}{LF})}: inhomogen varmelikning
\item[\vanskelig] \href{https://github.com/gustavbv/Matte4C---eksamen/blob/main/oppgavesett/F09-pde-og-bolgelikning/3-hele-oppgaven-i-pensum/027-TMA4125-30-35-2024-kont-oppg6-oppgave.pdf}{Matte 4N kont.\ 2024, oppg.~6} {\color{graa}(\href{https://github.com/gustavbv/Matte4C---eksamen/blob/main/oppgavesett/F09-pde-og-bolgelikning/3-hele-oppgaven-i-pensum/027-TMA4125-30-35-2024-kont-oppg6-svar.pdf}{LF})}: Laplace-likningen på et kvadrat med separasjon
\end{itemize}
\end{oving}


\subsection{Differanser}
\begin{formel}[Differanseformler (fra Taylorrekker)]
\begin{align*}
f'(x)&=\frac{f(x+h)-f(x)}h+O(h) && \text{(forover)},\\
f'(x)&=\frac{f(x)-f(x-h)}h+O(h) && \text{(bakover)},\\
f'(x)&=\frac{f(x+h)-f(x-h)}{2h}+O(h^2) && \text{(sentral)},\\
f''(x)&=\frac{f(x-h)-2f(x)+f(x+h)}{h^2}+O(h^2) && \text{(andre ordens sentral)}.
\end{align*}
\textbf{Feilorden} betyr hvor fort feilen går mot null: med $O(h^2)$ gir halv $h$ en fjerdedel av feilen. Begrunnelse: Taylor gir $f(x\pm h)=f\pm hf'+\frac{h^2}2f''\pm\frac{h^3}6f'''+\dots$, og leddene med odde potens kansellerer i de sentrale formlene.
\end{formel}

\begin{oppskrift}[differansemetoden generelt]
\begin{enumerate}
  \item \textbf{Gitter:} $x_i=x_0+ih$ ($i=0,\dots,M$, $h=L/M$) og $t_n=nk$. $U_i^n\approx u(x_i,t_n)$.
  \item Bytt hver deriverte med en differanse i de \emph{indre} punktene.
  \item Bruk \textbf{randbetingelsene} for $i=0$ og $i=M$. Dirichlet gir verdien direkte. Neumann og Robin: innfør et spøkelsespunkt $U_{-1}$ (eller $U_{M+1}$), uttrykk det med en sentral differanse for $u_x$, og sett inn.
  \item Bruk \textbf{initialbetingelsene}: $U_i^0=f(x_i)$, og eventuelt en ekstra betingelse for første tidssteg.
  \item Skriv skjemaet på formen $U_i^{n+1}=\dots$ (eksplisitt) eller $C\,\vb U^{n+1}=D\,\vb U^n+\dots$ (implisitt).
\end{enumerate}
\end{oppskrift}

\begin{formel}[Randverdiproblem i to punkter (fra oppsummeringen, F10a)]
Generelt: $u''+p(x)u'+q(x)u=r(x)$ på $[a,b]$. Innsatt i indre punkter:
\[
\frac{U_{i-1}-2U_i+U_{i+1}}{h^2}+p_i\frac{U_{i+1}-U_{i-1}}{2h}+q_iU_i=r_i,\qquad i=1,\dots,M-1,
\]
som gir et tridiagonalt lineært system. Dirichlet $u(a)=\alpha$: sett $U_0=\alpha$ og flytt det over til høyre side. Neumann $u'(a)=\sigma$: $\frac{U_1-U_{-1}}{2h}=\sigma$ gir $U_{-1}=U_1-2h\sigma$, og likningen i $i=0$ tas med. Robin $u'(a)+\beta u(a)=\gamma$ gir $U_{-1}=U_1-2h(\gamma-\beta U_0)$. Feilen er $O(h^2)$. Se Matte 4N vår 2024, oppg.~7, i oppgaveboksen under.
\end{formel}

\begin{formel}[5-punktsformelen for Laplace (prøveeksamen, oppg.~9)]
På et kvadratisk gitter med steglengde $h$:
\[
\Delta u(x_i,y_j)\approx\frac{U_{i-1,j}+U_{i+1,j}+U_{i,j-1}+U_{i,j+1}-4U_{i,j}}{h^2},\qquad i,j=1,\dots,N-1,
\]
med feil $O(h^2)$. I 3D får du 7-punktsformelen med $-6U_{i,j,k}$ og de seks naboene. Med Dirichlet-randverdier gir $\Delta u=f$ ett lineært system for de indre ukjente.
\end{formel}

\subsection{Numerikk for bølgelikningen}
\begin{formel}[Eksplisitt skjema for $u_{tt}=c^2u_{xx}$ med $u=0$ i endene]
Med $r=\frac{ck}h$ (Courant-tallet):
\begin{align*}
U_i^0&=f(x_i),\\
U_i^1&=U_i^0+k\,g(x_i)+\frac{r^2}2\bigl(U_{i+1}^0-2U_i^0+U_{i-1}^0\bigr),\\
U_i^{n+1}&=2U_i^n-U_i^{n-1}+r^2\bigl(U_{i+1}^n-2U_i^n+U_{i-1}^n\bigr),\qquad n\ge1,
\end{align*}
for $i=1,\dots,M-1$, og $U_0^n=U_M^n=0$. Feilen er $O(h^2+k^2)$.\\
\textbf{Første steg:} den sentrale differansen $\frac{U_i^1-U_i^{-1}}{2k}=g(x_i)$ gir $U_i^{-1}=U_i^1-2kg(x_i)$. Sett inn i det vanlige skjemaet med $n=0$ og løs for $U_i^1$.
\end{formel}
\begin{teorem}[CFL-betingelsen]
Det eksplisitte skjemaet er stabilt hvis og bare hvis $c\,\frac kh\le1$. En bølge får ikke bevege seg mer enn én gitteravstand $h$ per tidssteg $k$. Fast $h$ gir et største tillatte $k$, og omvendt.
\end{teorem}
\begin{eksempel}[$r=1$]
Med $c=2$, $h=1$ og $k=0{,}5$ er $r=1$. Er $g=0$, blir første steg $U_i^1=\frac12\bigl(U_{i-1}^0+U_{i+1}^0\bigr)$, en ren middelverdi av naboene. Det passer med øving 10, oppg.~2b.
\end{eksempel}

\subsection{Numerikk for varmelikningen}
La $\alpha=\dfrac{c^2k}{h^2}$ og $u_t=c^2u_{xx}$.
\begin{formel}[Eksplisitt Euler (forover Euler)]
\[
U_i^{n+1}=U_i^n+\alpha\bigl(U_{i-1}^n-2U_i^n+U_{i+1}^n\bigr),\qquad i=1,\dots,M-1,
\]
med $U_i^0=f(x_i)$. Feilen er $O(k+h^2)$.
\begin{itemize}
  \item \textbf{Dirichlet} $u(0,t)=g_0(t)$, $u(L,t)=g_L(t)$: $U_0^{n+1}=g_0(t_{n+1})$ og $U_M^{n+1}=g_L(t_{n+1})$.
  \item \textbf{Neumann} $u_x(0,t)=g_0$, $u_x(L,t)=g_L$, via spøkelsespunkt:
  \begin{gather*}
  U_0^{n+1}=U_0^n+2\alpha\bigl(U_1^n-U_0^n-hg_0(t_n)\bigr),\\ U_M^{n+1}=U_M^n+2\alpha\bigl(U_{M-1}^n-U_M^n+hg_L(t_n)\bigr).
  \end{gather*}
  \item \textbf{Robin} $a_0u_x(0,t)+b_0u(0,t)=g_0$ og $a_Lu_x(L,t)+b_Lu(L,t)=g_L$:
  \begin{gather*}
  U_0^{n+1}=U_0^n+2\alpha\Bigl(U_1^n-U_0^n+\tfrac{hb_0}{a_0}U_0^n-\tfrac h{a_0}g_0(t_n)\Bigr),\\
  U_M^{n+1}=U_M^n+2\alpha\Bigl(U_{M-1}^n-U_M^n-\tfrac{hb_L}{a_L}U_M^n+\tfrac h{a_L}g_L(t_n)\Bigr).
  \end{gather*}
  Fortegnet foran $\frac{hb_0}{a_0}$ er rettet fra lysarket og kontrollert mot øving 11 og numerisk.
  \item \textbf{Matriseform} (Neumann): $\vb U^{n+1}=(I+\alpha A)\vb U^n+2\alpha h\,\vb a^n$ med $\vb a^n=(-g_0(t_n),0,\dots,0,g_L(t_n))^\T$ og
  \[
  A=\begin{psmallmatrix}-2&2&&&\\1&-2&1&&\\&\ddots&\ddots&\ddots&\\&&1&-2&1\\&&&2&-2\end{psmallmatrix}\in\R^{(M+1)\times(M+1)}.
  \]
  For Dirichlet bruker du $B=\mathrm{tridiag}(1,-2,1)\in\R^{(M-1)\times(M-1)}$ for de indre punktene.
\end{itemize}
\end{formel}
\begin{teorem}[stabilitet for eksplisitt Euler]
Skjemaet er stabilt hvis og bare hvis $\alpha=\dfrac{c^2k}{h^2}\le\dfrac12$. (Iterasjonsmatrisen $Q=I+\alpha A$ må ikke forsterke noen egenvektor.) Halveres $h$, må $k$ deles på 4. Tidobbel romlig oppløsning krever et hundre ganger kortere tidssteg.
\end{teorem}

\begin{formel}[Implisitt Euler (bakover Euler)]
Bakoverdifferanse i tid:
\[
U_i^{n+1}-\alpha\bigl(U_{i-1}^{n+1}-2U_i^{n+1}+U_{i+1}^{n+1}\bigr)=U_i^n .
\]
\textbf{Dirichlet} (ukjente $U_1,\dots,U_{M-1}$):
\[
(I-\alpha B)\vb U^{n+1}=\vb U^n+\alpha\vb b^{n+1},\qquad \vb b^{n+1}=\bigl(g_0(t_{n+1}),0,\dots,0,g_L(t_{n+1})\bigr)^\T.
\]
\textbf{Neumann} (ukjente $U_0,\dots,U_M$):
\[
(I-\alpha A)\vb U^{n+1}=\vb U^n+2\alpha h\,\vb a^{n+1},\qquad \vb a^{n+1}=\bigl(-g_0(t_{n+1}),0,\dots,0,g_L(t_{n+1})\bigr)^\T.
\]
Fortegnene i $\vb a$ er rettet fra lysarket og testet numerisk.\\
Hvert tidssteg krever løsning av et tridiagonalt system, som koster $O(M)$. Skjemaet er \textbf{ubetinget stabilt}. Feilen er $O(k+h^2)$, så for nøyaktighet bør $k\sim h^2$.
\end{formel}

\begin{formel}[Crank--Nicolson]
Gjennomsnittet av eksplisitt og implisitt Euler:
\[
\Bigl(I-\tfrac\alpha2B\Bigr)\vb U^{n+1}=\Bigl(I+\tfrac\alpha2B\Bigr)\vb U^n+\tfrac\alpha2\bigl(\vb b^n+\vb b^{n+1}\bigr)\quad\text{(Dirichlet)},
\]
\[
\Bigl(I-\tfrac\alpha2A\Bigr)\vb U^{n+1}=\Bigl(I+\tfrac\alpha2A\Bigr)\vb U^n+\alpha h\bigl(\vb a^n+\vb a^{n+1}\bigr)\quad\text{(Neumann)}.
\]
Feilen er $O(h^2+k^2)$ og prisen omtrent som for implisitt Euler. Skjemaet er også ubetinget stabilt (standardresultat).
\end{formel}
\begin{center}
\small
\begin{tabular}{@{}llll@{}}
\toprule
Skjema & Feil & Stabilitet & Pris per tidssteg\\
\midrule
Eksplisitt Euler & $O(k+h^2)$ & $\alpha\le\frac12$ & billig (bare oppdatering)\\
Implisitt Euler & $O(k+h^2)$ & alltid & tridiagonalt system\\
Crank--Nicolson & $O(k^2+h^2)$ & alltid & tridiagonalt system\\
Bølge, eksplisitt & $O(k^2+h^2)$ & $ck/h\le1$ (CFL) & billig\\
\bottomrule
\end{tabular}
\end{center}

\begin{oving}
\begin{itemize}[leftmargin=5.6em,labelwidth=5em,labelsep=0.6em,align=left]
\item[\lett] \href{https://github.com/gustavbv/Matte4C---eksamen/blob/main/oppgavesett/F10-varmelikning-og-differanser/1-TMA4432-ovinger-og-proveeksamen/002-TMA4432-2026-proveeksamen-oppg9-oppgave.pdf}{Prøveeksamen, oppg.~9} {\color{graa}(\href{https://github.com/gustavbv/Matte4C---eksamen/blob/main/oppgavesett/F10-varmelikning-og-differanser/1-TMA4432-ovinger-og-proveeksamen/002-TMA4432-2026-proveeksamen-oppg9-svar.pdf}{LF})}: 5-punktsformelen for Laplace i 2D og 3D
\item[\lett] \href{https://github.com/gustavbv/Matte4C---eksamen/blob/main/oppgavesett/F10-varmelikning-og-differanser/2-Bergmann-eksamener/021-TMA4125-30-2023-kont-oppg9-oppgave.pdf}{Matte 4N kont.\ 2023, oppg.~9} {\color{graa}(\href{https://github.com/gustavbv/Matte4C---eksamen/blob/main/oppgavesett/F10-varmelikning-og-differanser/2-Bergmann-eksamener/021-TMA4125-30-2023-kont-oppg9-svar.pdf}{LF})}: eksplisitt skjema, to tidssteg for hånd
\item[\middels] \href{https://github.com/gustavbv/Matte4C---eksamen/blob/main/oppgavesett/F10-varmelikning-og-differanser/2-Bergmann-eksamener/015-TMA4125-30-2025-kont-oppg7-oppgave.pdf}{Matte 4N kont.\ 2025, oppg.~7} {\color{graa}(\href{https://github.com/gustavbv/Matte4C---eksamen/blob/main/oppgavesett/F10-varmelikning-og-differanser/2-Bergmann-eksamener/015-TMA4125-30-2025-kont-oppg7-svar.pdf}{LF})}: eksplisitt skjema og stabilitet
\item[\middels] \href{https://github.com/gustavbv/Matte4C---eksamen/blob/main/oppgavesett/F10-varmelikning-og-differanser/1-TMA4432-ovinger-og-proveeksamen/004-TMA4432-oving-10-oppg2-oppgave.pdf}{Øving 10, oppg.~2} {\color{graa}(\href{https://github.com/gustavbv/Matte4C---eksamen/blob/main/oppgavesett/F10-varmelikning-og-differanser/1-TMA4432-ovinger-og-proveeksamen/004-TMA4432-oving-10-oppg2-svar.pdf}{LF})}: differanseskjema for bølgelikningen
\item[\middels] \href{https://github.com/gustavbv/Matte4C---eksamen/blob/main/oppgavesett/F10-varmelikning-og-differanser/2-Bergmann-eksamener/018-TMA4125-2024-var-oppg7-oppgave.pdf}{Matte 4N vår 2024, oppg.~7} {\color{graa}(\href{https://github.com/gustavbv/Matte4C---eksamen/blob/main/oppgavesett/F10-varmelikning-og-differanser/2-Bergmann-eksamener/018-TMA4125-2024-var-oppg7-svar.pdf}{LF})}: randverdiproblem med Neumann og Dirichlet
\item[\vanskelig] \href{https://github.com/gustavbv/Matte4C---eksamen/blob/main/oppgavesett/F10-varmelikning-og-differanser/1-TMA4432-ovinger-og-proveeksamen/009-TMA4432-oving-11-oppg3-oppgave.pdf}{Øving 11, oppg.~3} {\color{graa}(\href{https://github.com/gustavbv/Matte4C---eksamen/blob/main/oppgavesett/F10-varmelikning-og-differanser/1-TMA4432-ovinger-og-proveeksamen/009-TMA4432-oving-11-oppg3-svar.pdf}{LF})}: Robin- og Neumann-rand, eksplisitt og implisitt
\item[\vanskelig] \href{https://github.com/gustavbv/Matte4C---eksamen/blob/main/oppgavesett/F10-varmelikning-og-differanser/3-hele-oppgaven-i-pensum/066-TMA4130-35-2018-host-oppg9-oppgave.pdf}{TMA4130 høst 2018, oppg.~9} {\color{graa}(\href{https://github.com/gustavbv/Matte4C---eksamen/blob/main/oppgavesett/F10-varmelikning-og-differanser/3-hele-oppgaven-i-pensum/066-TMA4130-35-2018-host-oppg9-svar.pdf}{LF})}: randverdiproblem med Robin-rand
\item[\vanskelig] \href{https://github.com/gustavbv/Matte4C---eksamen/blob/main/oppgavesett/F10-varmelikning-og-differanser/3-hele-oppgaven-i-pensum/079-TMA4130-2016-host-oppg10-oppgave.pdf}{TMA4130 høst 2016, oppg.~10} {\color{graa}(\href{https://github.com/gustavbv/Matte4C---eksamen/blob/main/oppgavesett/F10-varmelikning-og-differanser/3-hele-oppgaven-i-pensum/079-TMA4130-2016-host-oppg10-svar.pdf}{LF})}: Poisson-likningen på en trekant
\end{itemize}
\end{oving}


\subsection{PDE-er på hele \texorpdfstring{$\R$}{R} med Fouriertransform}
\begin{oppskrift}[varmelikningen på $\R$: $u_t=c^2u_{xx}$, $u(x,0)=f(x)$, $u\to0$ når $\abs x\to\infty$]
\begin{enumerate}
  \item Transformer i $x$: $\hat u_t(\omega,t)=-c^2\omega^2\hat u(\omega,t)$, en ODE i $t$ for hver $\omega$.
  \item Løs og bruk initialbetingelsen: $\hat u(\omega,t)=\hat f(\omega)\,\ee^{-c^2\omega^2t}$.
  \item Et produkt i frekvens er en konvolusjon i rom. Med $\ee^{-c^2\omega^2t}=\sqrt{2\pi}\,\hat g$ og tabellen for gaussfunksjoner får vi
  \[
  u(x,t)=\frac1{2c\sqrt{\pi t}}\int_{-\infty}^{\infty}f(y)\,\ee^{-\frac{(x-y)^2}{4c^2t}}\dd y\qquad\text{(varmekjernen, kontrollregnet)}.
  \]
\end{enumerate}
Er $f$ gaussisk (øving 11, oppg.~2), kan du regne $\Four^{-1}(\hat f\ee^{-c^2\omega^2t})$ direkte med tabellen.
\end{oppskrift}
\begin{oppskrift}[andre PDE-er på $\R$ (øving 11)]
\begin{itemize}
  \item \textbf{Bølgelikningen:} $\hat u_{tt}=-c^2\omega^2\hat u$ gir $\hat u=\hat f\cos(c\omega t)+\hat g\frac{\sin(c\omega t)}{c\omega}$. Transformert tilbake gir det d'Alemberts formel.
  \item \textbf{Første orden} $u_t=cu_x$: $\hat u_t=\ii c\omega\hat u$ gir $\hat u=\hat f\ee^{\ii c\omega t}$, og forskyvningsregelen gir $u(x,t)=f(x+ct)$.
  \item Generelt: transformer i $x$, løs ODE-en i $t$, og transformer tilbake med tabell, regneregler og konvolusjon.
\end{itemize}
\end{oppskrift}

\begin{oving}
\begin{itemize}[leftmargin=5.6em,labelwidth=5em,labelsep=0.6em,align=left]
\item[\lett] \href{https://github.com/gustavbv/Matte4C---eksamen/blob/main/oppgavesett/F10-varmelikning-og-differanser/1-TMA4432-ovinger-og-proveeksamen/008-TMA4432-oving-11-oppg2-oppgave.pdf}{Øving 11, oppg.~2} {\color{graa}(\href{https://github.com/gustavbv/Matte4C---eksamen/blob/main/oppgavesett/F10-varmelikning-og-differanser/1-TMA4432-ovinger-og-proveeksamen/008-TMA4432-oving-11-oppg2-svar.pdf}{LF})}: varmelikningen på $\R$ med gaussisk startverdi
\item[\middels] \href{https://github.com/gustavbv/Matte4C---eksamen/blob/main/oppgavesett/F10-varmelikning-og-differanser/1-TMA4432-ovinger-og-proveeksamen/007-TMA4432-oving-11-oppg1-oppgave.pdf}{Øving 11, oppg.~1} {\color{graa}(\href{https://github.com/gustavbv/Matte4C---eksamen/blob/main/oppgavesett/F10-varmelikning-og-differanser/1-TMA4432-ovinger-og-proveeksamen/007-TMA4432-oving-11-oppg1-svar.pdf}{LF})}: bølgelikningen på $\R$ med FT
\item[\vanskelig] \href{https://github.com/gustavbv/Matte4C---eksamen/blob/main/oppgavesett/F08-fouriertransform/1b-TMA4432-ekstraoppgaver/004-TMA4432-oving-11-oppg5-oppgave.pdf}{Øving 11, oppg.~5 (ekstra)} {\color{graa}(\href{https://github.com/gustavbv/Matte4C---eksamen/blob/main/oppgavesett/F08-fouriertransform/1b-TMA4432-ekstraoppgaver/004-TMA4432-oving-11-oppg5-svar.pdf}{LF})}: telegraflikningen med FT
\end{itemize}
\end{oving}


\section{F11 · Splines, B-splines og Bézierkurver}

\textbf{Kjernen:} Polynomer av høy grad oscillerer (Runge). Bruk i stedet stykkevise polynomer av lav grad som er så glatte som mulig i skjøtene. De utgjør et vektorrom. B-splines er en lokal basis for rommet. Bézierkurver bygges av Bernsteinpolynomer og styres med kontrollpunkter.

\subsection{Polynominterpolasjon (repetisjon fra Matte 3C)}
Gitt $n+1$ punkter $(x_i,y_i)$ med ulike $x_i$ finnes nøyaktig ett $p\in\mathcal P_n$ med $p(x_i)=y_i$. Lagrange-formen er
\[
p_n(x)=\sum_{i=0}^ny_i\,\ell_i(x),\qquad \ell_i(x)=\prod_{j\neq i}\frac{x-x_j}{x_i-x_j},\qquad \ell_i(x_j)=\delta_{ij}.
\]
\textbf{Runges eksempel:} $f(x)=\frac1{1+x^2}$ på $[-5,5]$ med ekvidistante noder gir store svingninger ved endene når $n$ vokser. Chebyshev-noder hjelper litt, men problemet med høy grad består. Det er motivasjonen for splines.

\subsection{Splinerommet}
\begin{defn}[spline av grad $k$]
La $\Delta=[x_0,\dots,x_n]$ med $a=x_0<x_1<\dots<x_n=b$. Da er $s\in\mathcal S_{k,\Delta}$ hvis
\begin{enumerate}
  \item $s|_{[x_j,x_{j+1}]}\in\mathcal P_k$ for $j=0,\dots,n-1$ (polynom av grad høyst $k$ på hvert intervall), og
  \item $s\in C^{k-1}([a,b])$ (funksjonen og de deriverte opp til orden $k-1$ er kontinuerlige, også i knutene).
\end{enumerate}
$\mathcal S_{k,\Delta}$ er et vektorrom. $k=1$ gir lineære, $k=2$ kvadratiske og $k=3$ kubiske splines.
\end{defn}
\begin{teorem}[dimensjon]
$\dim\mathcal S_{k,\Delta}=n+k$ ($n$ intervaller, altså $n+1$ knuter). Det er $(k+1)n$ koeffisienter og $k(n-1)$ glatthetskrav i de indre knutene. Kubisk: $4n-3(n-1)=n+3$.
\end{teorem}

\begin{oppskrift}[ligger $f$ i $\mathcal S_{k,\Delta}$? (prøveeksamen, oppg.~10)]
Sjekk begge kravene og begrunn hvert av dem:
\begin{enumerate}
  \item \textbf{Polynom på hvert intervall:} et polynom av grad $\le k$ ligger alltid i rommet. Grad $>k$ (for eksempel $x^4$ i $\mathcal S_3$) gjør det ikke. En knekk som \emph{ikke} er i en knute, for eksempel $\abs{x-\frac12}$ når $\frac12$ ikke er en knute, betyr at $f$ ikke er et polynom på det intervallet.
  \item \textbf{Glatthet i knutene:} sammenlign venstre- og høyregrenser for $f,f',\dots,f^{(k-1)}$ i hver indre knute.
  \item \textbf{Trunkerte potenser:} $(x-x_j)_+^m$ (lik $(x-x_j)^m$ for $x\ge x_j$, ellers $0$) er $C^{m-1}$ i $x_j$. Med $x_j$ som knute og $m=k$ ligger den i $\mathcal S_{k,\Delta}$. Er $m<k$, er den ikke glatt nok. For eksempel er $(x)_+^2$ bare $C^1$ i $0$, så den ligger ikke i $\mathcal S_3$.
  \item \textbf{Ukjente konstanter} (bestem $a,b$): still opp likningene for kontinuitet av $f,f',\dots,f^{(k-1)}$ i skjøten og løs.
\end{enumerate}
\end{oppskrift}

\subsection{Lineære og kubiske splines}
\begin{defn}[hattefunksjoner (lineære splines)]
For $i=1,\dots,n-1$:
\[
\Lambda_i(x)=\begin{cases}\frac{x-x_{i-1}}{x_i-x_{i-1}},&x_{i-1}\le x<x_i,\\ \frac{x_{i+1}-x}{x_{i+1}-x_i},&x_i\le x<x_{i+1},\\0&\text{ellers},\end{cases}
\]
og $\Lambda_0$, $\Lambda_n$ er «halve hatter» i endene. Siden $\Lambda_i(x_j)=\delta_{ij}$, interpolerer $s_1(x)=\sum_{i=0}^ny_i\Lambda_i(x)$ dataene.
\end{defn}

\begin{defn}[randbetingelser for kubisk interpolasjon]
Interpolasjon gir $n+1$ krav, og dimensjonen er $n+3$, så to krav må til:
\begin{itemize}
  \item \textbf{Naturlig:} $s''(x_0)=s''(x_n)=0$ (ingen krumning i endene).
  \item \textbf{Innspent (clamped):} $s'(x_0)$ og $s'(x_n)$ er gitt (på lysarket $=0$, altså vannrette tangenter).
  \item \textbf{Periodisk:} $s(x_0)=s(x_n)$, $s'(x_0)=s'(x_n)$ og $s''(x_0)=s''(x_n)$.
  \item \textbf{Not-a-knot:} $x_1$ og $x_{n-1}$ fjernes som knuter, men brukes fortsatt som interpolasjonspunkter.
\end{itemize}
Den naturlige kubiske splinen er også den $C^2$-interpolanten som minimerer $\int_a^b(s'')^2\dd x$, altså den som krummer minst (standardresultat).
\end{defn}

\begin{formel}[Kubisk interpolasjon med momentene $M_i=s''(x_i)$]
Med $h_i=x_i-x_{i-1}$ og $f_i=y_i$, for $i=1,\dots,n-1$:
\begin{gather*}
\mu_iM_{i-1}+2M_i+\lambda_iM_{i+1}=d_i,\qquad \lambda_i=\frac{h_{i+1}}{h_i+h_{i+1}},\qquad \mu_i=\frac{h_i}{h_i+h_{i+1}},\\
d_i=\frac{6}{h_i+h_{i+1}}\Bigl(\frac{f_{i+1}-f_i}{h_{i+1}}-\frac{f_i-f_{i-1}}{h_i}\Bigr).
\end{gather*}
Naturlig spline: $M_0=M_n=0$ gir et tridiagonalt $(n-1)\times(n-1)$-system. Er alle $h_i=h$, blir likningen $M_{i-1}+4M_i+M_{i+1}=\frac6{h^2}\bigl(f_{i+1}-2f_i+f_{i-1}\bigr)$.
Splinen på $[x_{i-1},x_i]$ er (siden $s''$ er lineær der):
\[
s(x)=M_{i-1}\frac{(x_i-x)^3}{6h_i}+M_i\frac{(x-x_{i-1})^3}{6h_i}+\Bigl(f_{i-1}-\frac{M_{i-1}h_i^2}6\Bigr)\frac{x_i-x}{h_i}+\Bigl(f_i-\frac{M_ih_i^2}6\Bigr)\frac{x-x_{i-1}}{h_i}.
\]
Begge formlene er kontrollert mot \texttt{scipy}. Med andre randbetingelser enn de naturlige endres første og siste rad i systemet.
\end{formel}

\subsection{B-splines}
\begin{defn}[B-spline $B_{i,k+1}$ av orden $m=k+1$ (grad $k$)]
Den entydige (opp til skalering) splinen i $\mathcal S_k$ med \textbf{støtte} $\supp B_{i,k+1}=[x_i,x_{i+k+1}]$, skalert slik at B-splinene summerer til 1.
\begin{gather*}
B_{i,1}(x)=\begin{cases}1,&x\in[x_i,x_{i+1}),\\0,&\text{ellers},\end{cases}\\
B_{i,k+1}(x)=\frac{x-x_i}{x_{i+k}-x_i}B_{i,k}(x)+\frac{x_{i+k+1}-x}{x_{i+k+1}-x_{i+1}}B_{i+1,k}(x)\qquad\text{(Cox--de Boor)}.
\end{gather*}
\end{defn}
\begin{eksempel}[lav grad (kontrollregnet)]
$B_{i,2}$ er hatten med topp $1$ i $x_{i+1}$. Den kvadratiske er
\[
B_{i,3}(x)=\begin{cases}
\dfrac{(x-x_i)^2}{(x_{i+1}-x_i)(x_{i+2}-x_i)},&[x_i,x_{i+1}),\\[8pt]
\dfrac{(x-x_i)(x_{i+2}-x)}{(x_{i+2}-x_i)(x_{i+2}-x_{i+1})}+\dfrac{(x_{i+3}-x)(x-x_{i+1})}{(x_{i+3}-x_{i+1})(x_{i+2}-x_{i+1})},&[x_{i+1},x_{i+2}),\\[8pt]
\dfrac{(x_{i+3}-x)^2}{(x_{i+3}-x_{i+1})(x_{i+3}-x_{i+2})},&[x_{i+2},x_{i+3}),\\[4pt]
0&\text{ellers}.
\end{cases}
\]
(Lysarket har bare det første leddet på det midterste intervallet.) For uniforme knuter har den kubiske $B_{i,4}$ verdiene $\frac16,\frac46,\frac16$ i de tre indre knutene i støtten.
\end{eksempel}
\begin{teorem}[egenskaper]
\begin{enumerate}
  \item $B_{i,m}$ er et polynom av grad $k$ på hvert intervall og ligger i $\mathcal S_k$.
  \item $\supp B_{i,m}=[t_i,t_{i+m}]$ og $B_{i,m}>0$ på $(t_i,t_{i+m})$. Med doble knuter kan støtten bli kortere enn du venter.
  \item \textbf{Partisjon av enheten:} $\sum_iB_{i,m}(x)=1$ på $[a,b]$, når de fiktive knutene er med.
  \item I et punkt er høyst $m-1=k$ B-splines ulik null (ved enkle knuter), så interpolasjonssystemet blir glissent (båndmatrise).
\end{enumerate}
\end{teorem}
\begin{defn}[fiktive knuter og sammenfallende knuter]
Uten ekstra knuter får vi bare $n-k$ B-splines, og summen er 1 bare på $[x_k,x_{n-k}]$. Legg derfor til $k$ fiktive knuter på hver side: $x_{-k}<\dots<x_{-1}<x_0$ og $x_n<x_{n+1}<\dots<x_{n+k}$. Da har vi $n+k$ lineært uavhengige B-splines $B_{-k,k+1},\dots,B_{n-1,k+1}$, en basis for $\mathcal S_{k,\Delta}$ (som stemmer med $\dim=n+k$):
\[
s(x)=\sum_{i=-k}^{n-1}c_iB_{i,k+1}(x).
\]
\textbf{Sammenfallende knuter} ($t_j\le t_{j+1}$, for eksempel $[0,0,0,1,1,2,3,4,5,5,5]$): bruk Cox--de Boor med konvensjonen $\frac{x-t_i}{t_{i+k}-t_i}:=0$ hvis $t_i=t_{i+k}$, og $B_{i,1}\equiv0$ hvis $t_i=t_{i+1}$. Lysarket skriver dette som $B_{i,k+1}=\lambda_{i,k+1}B_{i,k}+(1-\lambda_{i+1,k+1})B_{i+1,k}$. Gjentatte endeknuter gir «halve hatter» i endene. En knute med multiplisitet $\mu$ gir bare $C^{k-\mu}$ der (standardresultat).
\end{defn}
\begin{oppskrift}[interpolasjon med B-splines]
Ansatsen $s(x)=\sum_ic_iB_{i,k+1}(x)$ og kravene $s(x_j)=y_j$ gir et lineært system $\sum_ic_iB_{i,k+1}(x_j)=y_j$ med båndmatrise.
\begin{itemize}
  \item \textbf{Lineært:} $B_{i,2}(x_{i+1})=1$ og $B_{i,2}$ er null i de andre knutene, så $c_i=y_{i+1}$ (øving 12, oppg.~1).
  \item \textbf{Kubisk:} $n+1$ likninger og $n+3$ ukjente. Legg til to krav (naturlig, periodisk eller Hermite: $s'(a)=y_0'$ og $s'(b)=y_n'$).
\end{itemize}
\end{oppskrift}

\subsection{Bernsteinpolynomer og Bézierkurver}
\begin{defn}[Bernsteinpolynomer]
$b_{i,n}(t)=\binom ni t^i(1-t)^{n-i}$, $i=0,\dots,n$. De er $>0$ på $(0,1)$, summerer til 1 (binomialformelen $1=(t+(1-t))^n$), og oppfyller rekursjonen $b_{i,n}(t)=t\,b_{i-1,n-1}(t)+(1-t)\,b_{i,n-1}(t)$.
For $n=3$: $(1-t)^3$, $3t(1-t)^2$, $3t^2(1-t)$ og $t^3$.
\end{defn}
\begin{defn}[Bézierkurve]
Med kontrollpunkter $\vb p_0,\dots,\vb p_n\in\R^m$:
\[
\vb b(t)=\vb b(t;\vb p_0,\dots,\vb p_n)=\sum_{i=0}^n\vb p_i\,b_{i,n}(t)=\sum_{i=0}^n\binom ni(1-t)^{n-i}t^i\,\vb p_i,\qquad t\in[0,1].
\]
Hver komponent er et polynom av grad $n$, og kurven ligger i den konvekse innhyllingen av kontrollpunktene (vektene er positive og summerer til 1).
\end{defn}
\begin{teorem}[egenskaper (rettet, kontrollregnet)]
\[
\vb b(0)=\vb p_0,\quad \vb b(1)=\vb p_n,\quad \vb b'(0)=n(\vb p_1-\vb p_0),\quad \vb b'(1)=n(\vb p_n-\vb p_{n-1}),
\]
\[
\vb b''(0)=n(n-1)(\vb p_0-2\vb p_1+\vb p_2),\qquad \vb b''(1)=n(n-1)(\vb p_n-2\vb p_{n-1}+\vb p_{n-2}).
\]
Kurven går gjennom $\vb p_0$ og $\vb p_n$. Punktene $\vb p_1$ og $\vb p_{n-1}$ bestemmer start- og sluttretningen.
\end{teorem}
\begin{oppskrift}[de Casteljaus algoritme (stabil evaluering i $t_0$)]
\begin{enumerate}
  \item Start med $\vb b_{i,0}=\vb p_i$, $i=0,\dots,n$.
  \item For $k=1,\dots,n$ og $i=0,\dots,n-k$: $\vb b_{i,k}=(1-t_0)\vb b_{i,k-1}+t_0\vb b_{i+1,k-1}$.
  \item $\vb b(t_0)=\vb b_{0,n}$.
\end{enumerate}
Algoritmen bygger på $\vb b(t;\vb p_0,\dots,\vb p_n)=(1-t)\vb b(t;\vb p_0,\dots,\vb p_{n-1})+t\,\vb b(t;\vb p_1,\dots,\vb p_n)$.\\
\textbf{Eksempel} (kontrollregnet): $\vb p=(0,0),(1,2),(3,3),(4,0)$ og $t_0=\frac12$. Første nivå er $(0{,}5;1)$, $(2;2{,}5)$, $(3{,}5;1{,}5)$, andre nivå $(1{,}25;1{,}75)$, $(2{,}75;2)$, og resultatet er $\vb b(\frac12)=(2;1{,}875)$.
\end{oppskrift}
\begin{teorem}[sammensatte Bézierkurver]
To kubiske kurver $\vb b(t;\vb p_0,\dots,\vb p_3)$ for $t\in[0,1)$ og $\vb b(t-1;\vb q_0,\dots,\vb q_3)$ for $t\in[1,2]$:
\begin{itemize}
  \item kontinuerlig hvis $\vb p_3=\vb q_0$,
  \item deriverbar hvis i tillegg $\vb q_1-\vb q_0=\vb p_3-\vb p_2$, altså $\vb q_1=2\vb q_0-\vb p_2$ ($\vb p_2$, $\vb p_3=\vb q_0$ og $\vb q_1$ ligger med lik avstand på en linje).
\end{itemize}
\textbf{Hermite-data (øving 12, oppg.~3):} skal $\vb c$ på $[t_a,t_b]$ med $\Delta t=t_b-t_a$ ha gitte verdier og deriverte i endene, bruk $\vb c(t)=\vb b\bigl(\frac{t-t_a}{\Delta t}\bigr)$. Kjerneregelen gir
\[
\vb p_0=\vb y_a,\qquad \vb p_1=\vb y_a+\tfrac{\Delta t}3\vb y_a',\qquad \vb p_2=\vb y_b-\tfrac{\Delta t}3\vb y_b',\qquad \vb p_3=\vb y_b\qquad\text{(kontrollregnet)}.
\]
\end{teorem}
\textbf{Bruk:} TrueType-fonter, SVG og CAD. Videre emner er NURBS, Bézierflater og B-spline-flater.

\begin{oving}
\begin{itemize}[leftmargin=5.6em,labelwidth=5em,labelsep=0.6em,align=left]
\item[\lett] \href{https://github.com/gustavbv/Matte4C---eksamen/blob/main/oppgavesett/F11-splines-og-bezier/3-hele-oppgaven-i-pensum/023-TMA4215-2005-kont-oppg1-oppgave.pdf}{TMA4215 kont.\ 2005, oppg.~1} {\color{graa}(\href{https://github.com/gustavbv/Matte4C---eksamen/blob/main/oppgavesett/F11-splines-og-bezier/3-hele-oppgaven-i-pensum/023-TMA4215-2005-kont-oppg1-svar.pdf}{LF})}: bestem $a,b$ så $f$ blir en kvadratisk spline
\item[\lett] \href{https://github.com/gustavbv/Matte4C---eksamen/blob/main/oppgavesett/F11-splines-og-bezier/1-TMA4432-ovinger-og-proveeksamen/002-TMA4432-oving-12-oppg1-oppgave.pdf}{Øving 12, oppg.~1} {\color{graa}(\href{https://github.com/gustavbv/Matte4C---eksamen/blob/main/oppgavesett/F11-splines-og-bezier/1-TMA4432-ovinger-og-proveeksamen/002-TMA4432-oving-12-oppg1-svar.pdf}{LF})}: hattefunksjoner og lineær interpolasjon
\item[\middels] \href{https://github.com/gustavbv/Matte4C---eksamen/blob/main/oppgavesett/F11-splines-og-bezier/1-TMA4432-ovinger-og-proveeksamen/001-TMA4432-2026-proveeksamen-oppg10-oppgave.pdf}{Prøveeksamen, oppg.~10} {\color{graa}(\href{https://github.com/gustavbv/Matte4C---eksamen/blob/main/oppgavesett/F11-splines-og-bezier/1-TMA4432-ovinger-og-proveeksamen/001-TMA4432-2026-proveeksamen-oppg10-svar.pdf}{LF})}: hvilke funksjoner ligger i $\mathcal S_{3,\Delta}$?
\item[\middels] \href{https://github.com/gustavbv/Matte4C---eksamen/blob/main/oppgavesett/F11-splines-og-bezier/1-TMA4432-ovinger-og-proveeksamen/003-TMA4432-oving-12-oppg2-oppgave.pdf}{Øving 12, oppg.~2} {\color{graa}(\href{https://github.com/gustavbv/Matte4C---eksamen/blob/main/oppgavesett/F11-splines-og-bezier/1-TMA4432-ovinger-og-proveeksamen/003-TMA4432-oving-12-oppg2-svar.pdf}{LF})}: splinerom, dimensjon og ukjente koeffisienter
\item[\vanskelig] \href{https://github.com/gustavbv/Matte4C---eksamen/blob/main/oppgavesett/F11-splines-og-bezier/3-hele-oppgaven-i-pensum/022-TMA4215-2006-kont-oppg1-oppgave.pdf}{TMA4215 kont.\ 2006, oppg.~1} {\color{graa}(\href{https://github.com/gustavbv/Matte4C---eksamen/blob/main/oppgavesett/F11-splines-og-bezier/3-hele-oppgaven-i-pensum/022-TMA4215-2006-kont-oppg1-svar.pdf}{LF})}: kvadratiske splines med et minimeringskrav
\end{itemize}
\end{oving}

\begin{oving}
\begin{itemize}[leftmargin=5.6em,labelwidth=5em,labelsep=0.6em,align=left]
\item[\middels] \href{https://github.com/gustavbv/Matte4C---eksamen/blob/main/oppgavesett/F11-splines-og-bezier/1-TMA4432-ovinger-og-proveeksamen/004-TMA4432-oving-12-oppg3-oppgave.pdf}{Øving 12, oppg.~3} {\color{graa}(\href{https://github.com/gustavbv/Matte4C---eksamen/blob/main/oppgavesett/F11-splines-og-bezier/1-TMA4432-ovinger-og-proveeksamen/004-TMA4432-oving-12-oppg3-svar.pdf}{LF})}: sammensatt Bézierkurve og B-spline-system
\item[\vanskelig] \href{https://github.com/gustavbv/Matte4C---eksamen/blob/main/oppgavesett/F11-splines-og-bezier/2-Bergmann-eksamener/014-TMA4215-2021-host-oppg3-oppgave.pdf}{TMA4215 høst 2021, oppg.~3} {\color{graa}(\href{https://github.com/gustavbv/Matte4C---eksamen/blob/main/oppgavesett/F11-splines-og-bezier/2-Bergmann-eksamener/014-TMA4215-2021-host-oppg3-svar.pdf}{LF})}: kubisk Bézierkurve fra samplingsverdier
\end{itemize}
\end{oving}


\section{F12 · Minste kvadraters metode}

\textbf{Kjernen:} Flere målinger enn parametre gir et overbestemt system som ikke kan løses eksakt. Vi minimerer summen av kvadrerte avvik. Er modellen lineær i parametrene, gir det normallikningene (eller QR, eller $A^\dagger$). Er modellen ikke-lineær, bruker vi Gauss--Newton eller Levenberg--Marquardt.

\subsection{Lineære modeller}
En modell $y=m(t;\vb x)$ der parametrene $\vb x\in\R^n$ inngår \emph{lineært}, gir $m>n$ likninger $m(t_i;\vb x)=y_i$, altså $A\vb x\approx\vb b$.
\begin{itemize}
  \item Rett linje $a+bt$: rad $i$ er $(1,\ t_i)$. Polynom $a+bt+ct^2$: rad $(1,t_i,t_i^2)$ (Vandermonde-matrise).
  \item Periodisk: $a+b\cos(2\pi t)+c\sin(2\pi t)$, som er lineær i $(a,b,c)$ selv om den ikke er lineær i $t$.
  \item Eksponentiell $ae^{bt}$: ta logaritmen, $\ln y=\ln a+bt$, som er lineær i $(\tilde a,b)$ med $a=\ee^{\tilde a}$. Da minimerer du avvik i $\ln y_i$, ikke i $y_i$, og svaret blir et annet enn med den ikke-lineære modellen.
  \item B-spline- og Bézier-tilpasning (F11) er også lineære i koeffisientene.
\end{itemize}

\begin{defn}[lineært minste kvadraters problem]
For $A\in\R^{m\times n}$ med $m>n$ og full rang $n$:
\[
\vb x^*=\argmin_{\vb x\in\R^n}\norm{A\vb x-\vb b}_2^2=\argmin\sum_{i=1}^m\bigl(m(t_i;\vb x)-y_i\bigr)^2 .
\]
\textbf{Tolkning:} $A\vb x^*=\sum_jx_j^*\vb a_j$ er det punktet i $\ran A$ (kolonnerommet) som ligger nærmest $\vb b$, altså den ortogonale projeksjonen. Residualet $A\vb x^*-\vb b$ står vinkelrett på alle kolonnene.
\end{defn}
\begin{teorem}[tre måter å løse det på]
\begin{enumerate}
  \item \textbf{Normallikningene:} $A^\T A\,\vb x=A^\T\vb b$. Entydig løsning når $\rank A=n$. Dette er betingelsen $\nabla\frac12\norm{A\vb x-\vb b}^2=A^\T(A\vb x-\vb b)=\vb 0$. Numerisk ulempe: $\operatorname{cond}(A^\T A)=\operatorname{cond}(A)^2$, for eksempel med Vandermonde-matriser.
  \item \textbf{Moore--Penrose:} $\vb x^*=A^\dagger\vb b$ (F1). Fungerer også uten full rang, og gir da minimum-norm-løsningen.
  \item \textbf{QR:} $A=QR$ med ortogonal $Q\in\R^{m\times m}$ og øvre triangulær $R$. La $\tilde R\in\R^{n\times n}$ være den øverste delen av $R$, sett $\vb d=Q^\T\vb b$ og løs $\tilde R\vb x=(d_1,\dots,d_n)^\T$ ved tilbakesubstitusjon. Kvadratfeilen er $\norm{A\vb x^*-\vb b}_2^2=d_{n+1}^2+\dots+d_m^2$. Det følger av at $\norm{A\vb x-\vb b}=\norm{Q^\T(A\vb x-\vb b)}=\norm{R\vb x-\vb d}$.
\end{enumerate}
\end{teorem}
\begin{eksempel}[øving 1, oppg.~3: linje gjennom $(0;1{,}25)$, $(1;3)$, $(2;4{,}8)$, $(3;6{,}9)$ (kontrollregnet)]
$A^\T A=\begin{psmallmatrix}4&6\\6&14\end{psmallmatrix}$ og $A^\T\vb b=\begin{psmallmatrix}15{,}95\\33{,}3\end{psmallmatrix}$ gir
$\vb x=\frac1{20}\begin{psmallmatrix}14&-6\\-6&4\end{psmallmatrix}\begin{psmallmatrix}15{,}95\\33{,}3\end{psmallmatrix}=\begin{psmallmatrix}1{,}175\\1{,}875\end{psmallmatrix}$. Linjen er $p(t)=1{,}175+1{,}875t$, med feil $\norm{A\vb x-\vb b}_2\approx0{,}184$.
\end{eksempel}
\begin{fallgruve}[overtilpasning]
Med mange parametre (polynom av høy grad) kan modellen treffe støyen i stedet for trenden. Et polynom av grad $m-1$ går gjennom alle $m$ punktene, men svinger mellom dem. I øving 1 skjer dette med kubisk modell.
\end{fallgruve}

\begin{oving}
\begin{itemize}[leftmargin=5.6em,labelwidth=5em,labelsep=0.6em,align=left]
\item[\lett] \href{https://github.com/gustavbv/Matte4C---eksamen/blob/main/oppgavesett/F12-minste-kvadrater/3-hele-oppgaven-i-pensum/009-TMA4115-2025-var-oppg2-oppgave.pdf}{TMA4115 vår 2025, oppg.~2} {\color{graa}(\href{https://github.com/gustavbv/Matte4C---eksamen/blob/main/oppgavesett/F12-minste-kvadrater/3-hele-oppgaven-i-pensum/009-TMA4115-2025-var-oppg2-svar.pdf}{LF})}: rett linje gjennom tre punkter
\item[\middels] \href{https://github.com/gustavbv/Matte4C---eksamen/blob/main/oppgavesett/F12-minste-kvadrater/1-TMA4432-ovinger-og-proveeksamen/002-TMA4432-oving-1-oppg3-oppgave.pdf}{Øving 1, oppg.~3} {\color{graa}(\href{https://github.com/gustavbv/Matte4C---eksamen/blob/main/oppgavesett/F12-minste-kvadrater/1-TMA4432-ovinger-og-proveeksamen/002-TMA4432-oving-1-oppg3-svar.pdf}{LF})}: lineær regresjon med Moore--Penrose
\item[\middels] \href{https://github.com/gustavbv/Matte4C---eksamen/blob/main/oppgavesett/F12-minste-kvadrater/3-hele-oppgaven-i-pensum/031-TMA4115-2018-kont-oppg9-oppgave.pdf}{TMA4115 kont.\ 2018, oppg.~9} {\color{graa}(\href{https://github.com/gustavbv/Matte4C---eksamen/blob/main/oppgavesett/F12-minste-kvadrater/3-hele-oppgaven-i-pensum/031-TMA4115-2018-kont-oppg9-svar.pdf}{LF})}: andregradspolynom tilpasset fem punkter
\item[\vanskelig] \href{https://github.com/gustavbv/Matte4C---eksamen/blob/main/oppgavesett/F12-minste-kvadrater/2-Bergmann-eksamener/006-TMA4215-2021-host-oppg1-oppgave.pdf}{TMA4215 høst 2021, oppg.~1} {\color{graa}(\href{https://github.com/gustavbv/Matte4C---eksamen/blob/main/oppgavesett/F12-minste-kvadrater/2-Bergmann-eksamener/006-TMA4215-2021-host-oppg1-svar.pdf}{LF})}: QR med Givens-rotasjoner og minste kvadrater
\item[\vanskelig] \href{https://github.com/gustavbv/Matte4C---eksamen/blob/main/oppgavesett/F12-minste-kvadrater/2-Bergmann-eksamener/005-TMA4215-2022-host-oppg2-oppgave.pdf}{TMA4215 høst 2022, oppg.~2} {\color{graa}(\href{https://github.com/gustavbv/Matte4C---eksamen/blob/main/oppgavesett/F12-minste-kvadrater/2-Bergmann-eksamener/005-TMA4215-2022-host-oppg2-svar.pdf}{LF})}: gradient descent sammenlignet med normallikningen
\end{itemize}
\end{oving}


\subsection{Ikke-lineære minste kvadrater}
\begin{defn}
Med et residual $\vb R(\vb x)=(r_1(\vb x),\dots,r_m(\vb x))^\T$ (for eksempel $r_i=m(t_i;\vb x)-y_i$):
\[
\argmin_{\vb x\in\R^n}E(\vb x),\qquad E(\vb x)=\tfrac12\norm{\vb R(\vb x)}_2^2=\tfrac12\sum_{i=1}^mr_i(\vb x)^2,
\]
\[
\nabla E(\vb x)=\vb R'(\vb x)^\T\vb R(\vb x),\qquad \Hess E(\vb x)=\vb R'(\vb x)^\T\vb R'(\vb x)+\sum_{i=1}^mr_i(\vb x)\Hess r_i(\vb x).
\]
\end{defn}
\begin{center}
\small
\begin{tabular}{@{}lll@{}}
\toprule
& lineær & ikke-lineær\\
\midrule
Residual $\vb R(\vb x)$ & $A\vb x-\vb b$ & vilkårlig\\
Jacobimatrise $\vb R'(\vb x)$ & $A$, konstant & avhenger av $\vb x$\\
$\nabla E=\vb R'^\T\vb R$ & $A^\T(A\vb x-\vb b)$, lineær & ikke-lineær\\
$\nabla E=\vb 0$ & normallikningene & ikke-lineært system\\
Løses med & ett lineært system & iterasjon\\
Eksempel & $a+bx+cx^2$ & $a\ee^{bx}$\\
\bottomrule
\end{tabular}
\end{center}

\begin{oppskrift}[Gauss--Newton]
Newton på $\nabla E=\vb 0$, men uten andrederiverte: $\Hess E\approx J^\T J$ med $J=\vb R'$. Fra $\vb x^{(0)}$, for $k=0,1,2,\dots$:
\begin{enumerate}
  \item Regn ut $\vb R(\vb x^{(k)})$ og $J^{(k)}=\vb R'(\vb x^{(k)})$.
  \item Løs $(J^{(k)})^\T J^{(k)}\Delta^{(k)}=-(J^{(k)})^\T\vb R(\vb x^{(k)})$ (med QR eller $A^\dagger$). Det er det samme som det lineære minste kvadraters problemet $\Delta^{(k)}=\argmin_\Delta\norm{J^{(k)}\Delta+\vb R(\vb x^{(k)})}_2^2$.
  \item $\vb x^{(k+1)}=\vb x^{(k)}+\Delta^{(k)}$. Stopp når $\norm{\vb R}$ eller $\norm{\Delta}$ er liten.
\end{enumerate}
\textbf{Svakheter:} metoden konvergerer bare lokalt og kan divergere. Lineariseringen $\vb R(\vb x+\Delta)\approx\vb R(\vb x)+J\Delta$ er dårlig for store $\Delta$, så $E$ kan til og med øke. $J^\T J$ kan mangle full rang, og da er løsningen ikke entydig.
\end{oppskrift}
\begin{oppskrift}[Levenberg--Marquardt (Gauss--Newton med tillitsområde)]
Straff store steg: $E_\lambda(\Delta)=\frac12\norm{J\Delta+\vb R}_2^2+\frac\lambda2\norm\Delta_2^2=\frac12\Bigl\|\begin{psmallmatrix}J\\ \sqrt\lambda I\end{psmallmatrix}\Delta-\begin{psmallmatrix}-\vb R\\ \vb 0\end{psmallmatrix}\Bigr\|_2^2$. Normallikningen blir
\[
\bigl(J^\T J+\lambda I\bigr)\Delta=-J^\T\vb R,
\]
og matrisen er inverterbar for $\lambda>0$. Stor $\lambda$ gir korte steg, omtrent gradient descent med steglengde $1/\lambda$. Liten $\lambda$ gir omtrent Gauss--Newton. $1/\lambda$ fungerer som en tillitsradius.\\
\textbf{Algoritmen:} gitt $\vb x^{(0)}$ og $\lambda_0>0$, løs for $\Delta^{(k)}$. Hvis $E(\vb x^{(k)}+\Delta^{(k)})<E(\vb x^{(k)})$: ta steget og sett $\lambda_{k+1}=0{,}8\lambda_k$. Ellers: bli stående og sett $\lambda_{k+1}=2\lambda_k$.
\end{oppskrift}

\begin{oving}
Oppgavebanken har ingen gamle eksamensoppgaver om Gauss--Newton eller Levenberg--Marquardt. Bruk denne til å øve:\\
\textbf{Oppgave.} Modellen $m(t;a,b)=a\ee^{bt}$ skal tilpasses $(t_i,y_i)=(0;2)$, $(1;3)$, $(2;5)$. (a) Skriv opp $\vb R$ og $J$. (b) Ta ett Gauss--Newton-steg fra $(a,b)=(2;0{,}5)$. (c) Hva endres i (b) med Levenberg--Marquardt og $\lambda=1$?\\
\textbf{Fasit (kontrollregnet):} $r_i=a\ee^{bt_i}-y_i$, og rad $i$ i $J$ er $(\ee^{bt_i},\ at_i\ee^{bt_i})$. I startpunktet er $\vb R\approx(0;\,0{,}2974;\,0{,}4366)$ og $J\approx\begin{psmallmatrix}1&0\\1{,}6487&3{,}2974\\2{,}7183&10{,}873\end{psmallmatrix}$. Gauss--Newton gir $\Delta\approx(-0{,}0768;\,-0{,}0235)$, altså $(a,b)\approx(1{,}923;\,0{,}477)$, og $E$ synker fra $0{,}140$ til $0{,}0077$. Med $\lambda=1$ blir $\Delta\approx(-0{,}0507;\,-0{,}0304)$, et kortere steg i en litt annen retning ($E\approx0{,}0083$).
\end{oving}

\section{Oversikt: øvinger, prøveeksamen og sjekkliste}

\subsection{Øvingene}
\begin{center}
\small
\begin{tabularx}{\textwidth}{@{}l l X l@{}}
\toprule
Øving & Tema & Oppgaver (ekstraoppgaver i \emph{kursiv}) & Forelesning\\
\midrule
1 & SVD & SVD av $2\times2$ og $3\times2$, pseudoinvers med rangfeil, lineær regresjon & F1, F12\\
2 & Tensorer & Moore--Penrose (overbestemt), Kronecker/Khatri--Rao/Hadamard, ALS i Python & F1, F2\\
3 & Kurver og flater & ellipse, skisser, sykloide, partiellderiverte, kule, ellipsoide, graf & F3\\
4 & Jacobi og Newton & linearisering, singulær Jacobimatrise, Lotka--Volterra, stav og fjær & F4\\
5 & Optimering & kjerneregel, gradienter og Hessematriser, Tikhonov, gradient descent, Lagrange & F5\\
6 & Fourierrekker & periode $2L$, reell og kompleks form, stykkevise funksjoner, \emph{$\cos^2$, $\abs{\sin}$, egenskaper} & F6\\
7 & Fourierrekker II & halvperiodeutvidelse, konvolusjon, $c_k(f')$, \emph{komplekse rekker, $\ee^{-x^2}*\ee^{-2x^2}$, $f*f$} & F6, F8\\
8 & Fouriertransform & FT av $x^2\chi$, $f*f=g$, DFT i Python og filtrering, \emph{sinc$*$sinc, gaussintegral, DFT-bevis} & F7, F8\\
9 & PDE & klassifisering, superposisjon, streng med trapesprofil, \emph{4.\ orden, bølge på $\R$, demping} & F9\\
10 & PDE II & d'Alembert, differanseskjema for bølgen, varmelikning (Dirichlet og Neumann), \emph{d'Alembert, stiv stav} & F10\\
11 & PDE III & bølge og varme på $\R$ med FT, Robin og Neumann numerisk, \emph{1.\ orden, telegraflikningen} & F10\\
12 & Splines & hatter og B-splines, kubisk splinerom, sammensatt Bézier og B-spline-system & F11\\
\bottomrule
\end{tabularx}
\end{center}

\subsection{Prøveeksamenen (10 oppgaver, 100 poeng)}
\begin{center}
\small
\begin{tabular}{@{}rlll@{}}
\toprule
Nr. & Tema & Poeng & Forelesning\\
\midrule
1 & Moore--Penrose: SVD, $A^\dagger$, minste feil, $(A^\T)^\dagger$ & 10 & F1\\
2 & Tucker vs.\ CP vs.\ SVD (forståelse) & 8 & F2\\
3 & Tangentplan til en flate gitt implisitt & 8 & F3, F5\\
4 & Ett Newton-steg for et $2\times2$-system & 7 & F4\\
5 & Gradient descent, steglengde og startpunkt, Newton & 10 & F5\\
6 & Fourierkoeffisienter ved avlesning og konvolusjon & 12 & F6\\
7 & DFT: båndbegrensning og rekonstruksjon & 10 & F7\\
8 & Bølgelikning på $[-1,1]$ med separasjon & 14 & F9\\
9 & 5- og 7-punktsformelen for Laplace & 10 & F10\\
10 & Hvilke funksjoner ligger i $\mathcal S_{3,\Delta}$? Dimensjon, naturlig & 11 & F11\\
\bottomrule
\end{tabular}
\end{center}
Poengene summerer til 100.

\subsection{Sjekkliste: dette må du kunne uten formelarket}
\begin{multicols}{2}
\small
\begin{itemize}
  \item Regne ut SVD for hånd, $\vb u_k=A\vb v_k/\sigma_k$
  \item $A^\dagger=\sum\sigma_k^{-1}\vb v_k\vb u_k^\T$, minste kvadrater og minst norm
  \item Eckart--Young: feilen er $\sigma_{k+1}$ (eller roten av halesummen)
  \item rang, bilde og kjerne fra SVD
  \item Kronecker, Khatri--Rao og Hadamard: definisjon og størrelse
  \item matrisering og $\times_k$, $X_{(1)}=A(C\odot B)^\T$
  \item ALS og HOSVD i ord, CP som spesialtilfelle av Tucker
  \item glatthet for kurver og flater, tangentlinje og -plan, normal
  \item buelengde, fart og akselerasjon
  \item deriverbarhet og linearisering
  \item Newton-steg for hånd, når metoden feiler
  \item kjerneregelen med matrisestørrelser
  \item $D_{\vb r}f=\ip{\nabla f}{\vb r}$, $\nabla f\perp$ nivåmengder
  \item gradient descent og Newton med Hessematrisen
  \item Lagrange: $\nabla f=\lambda\nabla g$, $g=0$
  \item Fourierrekker med periode $2L$, like og odde utvidelser
  \item reell $\leftrightarrow$ kompleks: $c_n=\frac{a_n-\ii b_n}2$
  \item konvergens i hopp (midtpunktet), Gibbs
  \item Parseval og rekkesummer
  \item $c_k(f*g)=2\pi c_kc_k$, $c_k(f')=\ii kc_k$, forskyvning $\ee^{-\ii kx_0}$
  \item Dirichlet-kjernen
  \item DFT-egenskaper, aliasing, $\FN^{-1}=\frac1N\overline{\FN}$, $\mathcal F_4$
  \item FFT: idé og $O(N\log N)$
  \item $\Four(f')=\ii\omega\hat f$, $\Four(f*g)=\sqrt{2\pi}\hat f\hat g$
  \item klassifisering av PDE-er, superposisjon
  \item utledning med separasjon av variable (tre tilfeller av $k$)
  \item løsningen av varmelikningen (sinus og cosinus)
  \item d'Alemberts formel
  \item Laplace på et rektangel (sinh)
  \item varmekjernen på $\R$
  \item differanseformler med orden
  \item bølgeskjemaet med første steg, CFL
  \item eksplisitt Euler, implisitt Euler, Crank--Nicolson, $\alpha\le\frac12$
  \item spøkelsespunkt for Neumann og Robin
  \item 5-punktsformelen
  \item splinedefinisjonen, $\dim=n+k$, randbetingelser
  \item Cox--de Boor og egenskapene til B-splines
  \item Bernstein, Bézier og de Casteljau
  \item $C^1$-skjøting av Bézierkurver
  \item normallikningene, QR og $A^\dagger$
  \item Gauss--Newton og Levenberg--Marquardt
\end{itemize}
\end{multicols}


\end{document}
